\documentclass[11pt]{amsart}

\usepackage[T1]{fontenc}
\usepackage{lmodern}
\usepackage{amsmath,amssymb,amsthm}
\usepackage{booktabs}
\usepackage{needspace}
\usepackage{microtype}
\usepackage[margin=2.8cm]{geometry}
\usepackage{xcolor}
\usepackage[colorlinks=true,linkcolor=blue!50!black,citecolor=blue!50!black,urlcolor=blue!50!black]{hyperref}
\usepackage{siunitx}

\newtheorem{theorem}{Theorem}[section]
\newtheorem{proposition}[theorem]{Proposition}
\newtheorem{lemma}[theorem]{Lemma}
\newtheorem{corollary}[theorem]{Corollary}
\newtheorem{maintheorem}{Theorem}
\renewcommand{\themaintheorem}{\Alph{maintheorem}}
\theoremstyle{definition}
\newtheorem{definition}[theorem]{Definition}
\newtheorem{remark}[theorem]{Remark}

\newcommand{\Rr}{R}
\newcommand{\Kk}{K}
\newcommand{\eee}{\mathrm{e}}
\newcommand{\ii}{\mathrm{i}}
\newcommand{\chat}{\hat{c}_1}

\hypersetup{pdftitle={Kneser's tetration continued around the cusp e^(1/e): separation from regular iteration and the inverse horn map of e^u-1},pdfauthor={Guanghao Li}}
\begin{document}

\title[Kneser's tetration around the cusp and the horn map of $\eee^u-1$]{Kneser's tetration continued around the cusp $\eee^{1/\eee}$: separation from regular iteration and the inverse horn map of $\eee^u-1$}

\author{Guanghao Li}
\email{wsmslgh@gmail.com}

\date{\today}

\begin{abstract}
For $1<b<\eta=\eee^{1/\eee}$ the equation $F(z+1)=b^{F(z)}$, $F(0)=1$, has the
regular solution $\Rr_b$ at the attracting fixed point.  We prove that
Kneser's real solution for $b>\eta$, continued analytically in the base around
the cusp $\eta$, crosses the Shell--Thron boundary (jointly holomorphically in
base and height, whatever the arithmetic of the neutral multiplier), extends
single-valued to the whole upper half of the Shell--Thron region, and has
boundary values $\Kk_b$ on $(1,\eta)$ along every path there.  These are a solution $\Kk^W_b$ obtained by sewing the two regular
charts, characterised, for $b$ near $\eta$, by the position of its seam.  Moreover, with a
computer-assisted certificate, Kneser's family continues to a single-valued
holomorphic family on the whole upper half $b$-plane: Kneser's own straight
initial region survives on the entire exterior, and the family crosses every
point of the open upper Shell--Thron boundary arc.  So the continuation is path
independent, it coincides with Paulsen's base-continued family, and it is
holomorphic across every real base in $(0,\eta)\setminus\{1\}$ (in particular
at the neutral endpoint $\eee^{-\eee}$).  Writing
$\Kk_b=\Rr_b\circ(\mathrm{id}+\sum_{n\ge0}c_n\eee^{2\pi\ii nz})$ and
$\Lambda=\exp(4\pi^2/\log\lambda)$, we show $c_n/\Lambda^n\to B_n\eee^{2\pi\ii na}$ as
$b\to\eta^-$ for every $n$, where $B_n$ are the Fourier coefficients of the
inverse \'Ecalle--Voronin horn map of $\eee^u-1$ and $a$ is the Fatou coordinate
of the base point; $B_1\ne0$.  The proof combines a confluence theorem for the
transition map between the two regular charts, a contraction in a weighted
Wiener algebra, and a family of tilted initial regions whose quotients are
uniformly quasiconformal to a flat cylinder.  The corrections have an
asymptotic expansion to all orders in the unfolding parameter with
coefficients given by convergent parabolic series.  Numerically, the sewn
solution reproduces Paulsen's $120$-digit value at $b=\sqrt2$ in all quoted
digits.
\end{abstract}

\maketitle

\section{Introduction}\label{sec:intro}

The tetration equation
\begin{equation}\label{eq:abel}
  F(z+1)=b^{F(z)},\qquad F(0)=1,
\end{equation}
has, for real $b>\eta:=\eee^{1/\eee}$, a distinguished real-analytic solution
constructed by Kneser~\cite{Kneser1950}.  It tends to the two complex-conjugate
fixed points of $w\mapsto b^w$ as $\operatorname{Im}z\to\pm\infty$.
Its conformal uniformisation is uniquely characterised by the
initial-region criterion of~\cite{TrappmannKouznetsov2011}.  For
$1<b<\eta$ the fixed points are real: $L=\eee^\lambda$ is attracting, with
$0<\lambda<1$, and $L_2=\eee^{\lambda_2}$ is repelling, with $\lambda_2>1$.  The
Koenigs construction at $L$ yields the regular solution $\Rr_b$.  This is the
solution usually meant in that range~\cite{KouznetsovTrappmann2010}.
Paulsen~\cite{Paulsen2019} proposed continuing Kneser's two-fixed-point
construction through the upper half of the Shell--Thron region.  If this
base continuation $\kappa$ has a boundary value, denote it by
\begin{equation}\label{eq:K}
  \Kk_b=\lim_{y\downarrow0}\kappa_{b+\ii y}.
\end{equation}
High-precision algorithms for the relevant pointwise constructions are
developed in~\cite{Kouznetsov2009,KouznetsovTrappmann2010,KouznetsovTrappmann2012}.
A different line of work constructs non-regular solutions by infinite
compositions: Nixon~\cite{Nixon2022} builds, for $b$ in the interior of the
Shell--Thron region, a family of solutions of period $2\pi\ii/\lambda$ and
compares them with the regular solution through a $1$-periodic
``$\theta$-mapping'' $F=\Rr_b\circ(\mathrm{id}+\theta)$, the normal form we
also use below.  That work does not treat Kneser's solution for $b<\eta$, the
limit $b\to\eta$, or continuation in the base, and for $b>\eta$ the link
between its construction and Kneser's solution is left as a conjecture.
Here the $1$-periodic correction is determined by the transition map between
the two regular charts, and its Fourier coefficients are shown to be
asymptotic to $\Lambda^nB_n\eee^{2\pi\ii na}$ (notation below).
We construct a two-sided sewing solution $\Kk^W_b$ independently and
determine its difference from $\Rr_b$.  We then prove that the
continuation of Kneser's solution through the upper half-disc about
$\eta$, and then near the real axis, has boundary values $\Kk_b=\Kk^W_b$ (Theorem~\ref{mthm:F}), so every
result about $\Kk^W_b$ holds for these boundary values.

\subsection*{Notation}
Let $h=2\pi/|\log\lambda|$ be the imaginary period of $\Rr_b$, and put
\[
 \Lambda=\eee^{-2\pi h}=\exp\bigl(4\pi^2/\log\lambda\bigr).
\]
Both $\Rr$ and $\Kk^W$ are normalised by $F(0)=1$.  The map
$P=\Rr^{-1}\circ\Kk^W-\mathrm{id}$
is $1$-periodic with $P(0)=0$.  In the relevant region it has only
non-negative modes, $P=\sum_{n\ge0}c_n\eee^{2\pi\ii nz}$, and the separation is
\[
 \Kk^W_b(z)-\Rr_b(z)=c_1\Rr_b'(z)\bigl(\eee^{2\pi\ii z}-1\bigr)+O(c_2)+O(c_1^2).
\]
At $b=\eta$ the map becomes the parabolic germ $f(u)=\eee^u-1$ in
$w=\eee(1+u)$.  Let $\Phi_{\mathrm{att}},\Phi_{\mathrm{rep}}$ be its Fatou
coordinates, both normalised by the formal Abel function
$\alpha(u)=-2/u+\tfrac13\log(-u)+O(u)$.  Let
\[
 \Phi_{\mathrm{att}}\circ\Phi_{\mathrm{rep}}^{-1}(z)=z+\sum_{n\ge1}B_n\eee^{2\pi\ii nz}
\]
be the inverse horn map on the upper gate~\cite{Ecalle1981,Voronin1981}, and $a=\Phi_{\mathrm{att}}(1/\eee-1)$
the attracting Fatou coordinate of the base point $w=1$.  Numerically
\[
 \begin{aligned}
 |B_1|&=0.0890584364122133\ldots, &
 \frac{B_2}{B_1^2}&=-0.0253204931+1.5742190652\,\ii,\\
 a&=3.0292972144180\ldots
 \end{aligned}
\]
(Section~\ref{sec:setup}).  The transition map between the two regular
solutions is $T_b=\Rr_b^{-1}\circ S_b$, where $S_b$ is the regular solution at
the repelling point.  Write $T_b(z)-z=\sum t_n\eee^{2\pi\ii nz}$ and
$\tau_n=t_n\eee^{-2\pi\ii nt_0}$ (Section~\ref{sec:transition}).

\emph{Symbols with more than one use.}  A few letters are unavoidably
overloaded; their meaning is fixed within each section.
\begin{itemize}
\item $\varepsilon$: $1-\lambda$ in Proposition~\ref{prop:saddle} and, as a complex
  parameter, in Section~\ref{sec:cusp}; $|\log\lambda|$ in
  Sections~\ref{sec:W2}--\ref{sec:char}.  The two agree to first order.
\item $\tau$: the repelling Koenigs map (Section~\ref{sec:transition}); the
  invariants $\tau_n$; and $T-\mathrm{id}-t_0$ in Theorem~\ref{thm:W2}.
\item $a$: the Fatou coordinate $\Phi_{\mathrm{att}}(1/\eee-1)$ above; $a=\log b$
  only inside Remark~\ref{rem:paulsen} and
  Propositions~\ref{prop:exterior-real-parameter}
  and~\ref{prop:backward-log-monodromy}.
\item $h$: the period $2\pi/|\log\lambda|$; the horn map and its inverse are
  written $h^{\pm1}$ only in Theorem~\ref{thm:reduction}, Section~\ref{sec:proofC}
  and Section~\ref{sec:first}.
\item $s$: $1-\eee\log b$ in Section~\ref{sec:first}; a shift $s(b)$ in
  Theorems~\ref{thm:reduction} and~\ref{thm:C}; an interpolation parameter in Section~\ref{sec:cusp}.
\item $\kappa$: the base-continued Kneser family $\kappa_b$; $\kappa_n=c_n/c_1^n$ and the
  first-order coefficients $\kappa^{(n)}$ carry indices.
\end{itemize}

\subsection*{Main results}

\begin{maintheorem}[Confluence]\label{mthm:A}
For every $n\ge1$, $\tau_n(b)\to B_n\eee^{2\pi\ii na}$ as $b\to\eta^-$.
Moreover, $B_1\ne0$ and $t_n/t_1^n\to B_n/B_1^n$ for every $n\ge1$.
\end{maintheorem}

\begin{maintheorem}[The sewing solution]\label{mthm:B}
There is $b_1<\eta$ such that for $b\in(b_1,\eta)$ the following holds.  There
is exactly one solution $\Kk^W_b$ of~\eqref{eq:abel} with the following three
properties:
\begin{itemize}
\item it has a representation $\Rr_b\circ(\mathrm{id}+P)$ on an upper half-plane;
\item it has a representation $S_b\circ(\mathrm{id}+Q)$ on a lower half-plane;
\item $P$ and $Q$ are bounded and $1$-periodic, and the seam lies half a period
  below the real axis.
\end{itemize}
The precise statement is Definition~\ref{def:class}.  Its coefficients satisfy
\[
 \Bigl|\frac{c_n}{\Lambda^n}-\tau_n(b)\Bigr|\le C_n\Lambda
\]
with explicit constants $C_n$ independent of $b$.
\end{maintheorem}

\begin{corollary}\label{cor:limits}
For $\Kk^W_b$, $\hat c_n:=c_n/\Lambda^n\to B_n\eee^{2\pi\ii na}$ for every $n\ge1$.
In particular
\[
 \lim_{b\to\eta^-}\frac{|c_1|}{\Lambda}=|B_1|.
\]
The limiting phase is $\arg B_1+2\pi a$, and
$c_n/c_1^n\to B_n/B_1^n$ for every $n\ge1$.
\end{corollary}

\begin{maintheorem}[First-order term]\label{mthm:D}
Let $s=1-\eee\log b$ and $p=|\log\lambda|\log\lambda_2=2s+O(s^2)$.  For every
$n\ge1$ with $B_n\ne0$ (in particular $n=1$ by Proposition~\ref{prop:horn-nonzero}),
\[
 \log\frac{\tau_n(b)}{B_n\eee^{2\pi\ii na}}=\kappa^{(n)}p+O(p^{2})\qquad(b\to\eta^-),
\]
where $\kappa^{(n)}$ is given by an absolutely convergent series over orbits of
the parabolic germ $\eee^u-1$ (Theorem~\ref{thm:D}; the remainder $O(p^2)$ is
Theorem~\ref{thm:E} with $m=1$).  More generally,
$\log(\tau_n/B_n\eee^{2\pi\ii na})$ has an asymptotic expansion
$\sum_{j\ge1}\kappa^{(n)}_jp^j$ to all orders, and every coefficient is built from
convergent parabolic series of the same kind (Theorem~\ref{thm:E}).
\end{maintheorem}

The sewn solution is defined for $b<\eta$ without reference to Kneser's
construction.  The next theorem shows that it is exactly what Kneser's
solution becomes when continued in the base around the cusp.

\begin{maintheorem}[Continuation around the cusp]\label{mthm:F}
Kneser's real-normalised solution for real $b>\eta$ continues analytically in the
base through the upper half-disc about $\eta$, across the Shell--Thron boundary arc
there (whatever the arithmetic of the neutral multiplier), and single-valued
along every path in the upper half of the Shell--Thron region.  Its boundary
values on $(1,\eta)$ are $\Kk_b=\Kk^W_b$, so Corollary~\ref{cor:limits} holds for $\Kk_b$:
$c_n/\Lambda^n\to B_n\eee^{2\pi\ii na}$ for every $n\ge1$, and $|c_1|/\Lambda\to|B_1|>0$.  Continuation
through the lower half-disc gives the complex-conjugate branch.
\end{maintheorem}

\begin{maintheorem}[The whole upper half-plane; computer-assisted]\label{mthm:G}
The continuation of Theorem~\ref{mthm:F} is a single-valued holomorphic family on
the whole upper half $b$-plane.  It crosses every point of the open upper
Shell--Thron boundary arc holomorphically, it is holomorphic at the neutral
endpoint $b=\eee^{-\eee}$, and it extends holomorphically across every point of
$(0,\eta)\setminus\{1\}$; $b=0$ and $b=1$ are genuine singularities.  In particular the continuation
is path independent, and it is the base-continued family of~\cite{Paulsen2019}.
\end{maintheorem}

The precise statements are Theorems~\ref{thm:cusp} and~\ref{thm:global} for
Theorem~\ref{mthm:F}, and Theorems~\ref{thm:far}, \ref{thm:last-arc}, \ref{thm:upper},
\ref{thm:endpoint}, \ref{thm:positive-real-axis} and
Proposition~\ref{prop:zero-one-obstruction} for Theorem~\ref{mthm:G}.
\begin{itemize}
\item \emph{Around the cusp} (Section~\ref{sec:cusp}), Kneser's initial region (the chord
  between the fixed points and its image) is tilted continuously along the path.
  In the coordinate $\log\frac{u-u_1}{u-u_2}$ it is a thin strip whose quotient is uniformly
  quasiconformal to a flat cylinder, so the neutral boundary is no obstacle.  At
  real $b<\eta$ it is a fundamental domain of the sewn surface, and the seam
  condition of Definition~\ref{def:class} identifies the limit with $\Kk^W_b$.  A
  quasiconformal transport of the dynamics extends the family over the whole
  upper Shell--Thron region (Section~\ref{sec:global}).
\item \emph{Away from the cusp} (Sections~\ref{sec:far}--\ref{sec:two-cycle}), Kneser's own
  straight chord remains a crescent with a canonical marked point on the whole
  exterior part of the upper half-plane and across the boundary arc.  This is
  verified by interval arithmetic, in a parametrisation by the multiplier of
  the continued fixed point.
\item The family of~\cite{Paulsen2019} is defined by analytic continuation in the
  base, but there the crossing of the Shell--Thron boundary is argued only
  numerically, and the uniqueness statement needs an additional
  root-selection condition (Remark~\ref{rem:paulsen}).  Theorem~\ref{mthm:G} settles
  the first point.  The nine-digit agreement at $b=\sqrt2$ is a numerical check.
\end{itemize}
By Lemma~\ref{lem:param} and Proposition~\ref{prop:saddle}, $\Lambda$ has the
saddle-node form
\[
 \Lambda=\exp\bigl(-C_\eta(\eta-b)^{-1/2}(1+o(1))\bigr),\qquad
 C_\eta=\frac{4\pi^2}{\sqrt{2\eee^{1-1/\eee}}}=20.3508\ldots .
\]
Thus the coefficients of the separation have the stated exponentially
small upper bounds.  Proposition~\ref{prop:horn-nonzero} makes the
first-mode asymptotic sharp with a strictly positive constant, and
Proposition~\ref{prop:horn-intervals} gives a certified enclosure of it.

\subsection*{Numerical results}
Theorems~\ref{mthm:A} and~\ref{mthm:B} are confirmed and extended numerically
in Section~\ref{sec:numerics}.
\begin{enumerate}
\item Ladder computations of $\Kk_b$ reproduce $\tau_1$ and $t_2/t_1^2$ to their
  full precision, arguments included; for example, eight digits at $b=1.1$.
\item At $b=\sqrt2$, far from $\eta$, the sewing equation gives
  $\operatorname{Im}\Kk^W_{\sqrt2}(\tfrac12)=-1.18899697185\cdot10^{-48}$, which agrees with
  Paulsen's $120$-digit value $-1.18899697\cdot10^{-48}$ for his base-continued
  family~\cite[\S6]{Paulsen2019} in all nine quoted digits, sign included, as
  Theorem~\ref{mthm:G} predicts.
\item $\tau_1$, $t_2/t_1^2$ and $t_3/t_1^3$ have been followed up to
  $\eta-b=3\cdot10^{-5}$.
\item The parabolic series of Theorem~\ref{mthm:D} gives
  \[
   \log\frac{\tau_1}{B_1\eee^{2\pi\ii a}}=(-0.010199006345897402441-0.013250956182506712250\,\ii)\,p+O(p^{2}),
  \]
  in agreement with fits from the side $b<\eta$ to their eleven digits.  The
  second coefficient,
  $\kappa^{(1)}_2=-7.88992074991511334\cdot10^{-5}+3.6554727325799459\cdot10^{-4}\,\ii$,
  agrees with the fits to within one unit in their sixth digit.
\end{enumerate}

\subsection*{Method}
Four observations drive the proofs.
\begin{enumerate}
\item The separation is a property of the transition map $T_b$ between two
  classical Koenigs coordinates.  A welding identity expresses $c_n$ through
  the Fourier modes of $T_b$ (Theorem~\ref{thm:welding}).  The scale $\Lambda$
  appears because the seam sits half a period below the real axis, and the
  gate carrying the positive modes a further half period below it.
\item The limit of $T_b$ is a confluence problem for the saddle-node unfolding
  in the direction where both fixed points are real (the ``Glutsyuk''
  direction~\cite{Glutsyuk2001,MRR2004,ChristopherRousseau}).  In the
complementary Lavaurs direction, the corresponding statements are the classical
results of parabolic implosion~\cite{Lavaurs1989,Shishikura2000,Oudkerk1999}.
For vector fields, see~\cite{RousseauTeyssier2008}.  We prove the
  needed statement directly (Section~\ref{sec:proofC}), using the model time
  $A_1\log(u-u_1)+A_2\log(u-u_2)$, whose one-step error vanishes at both fixed
  points.
\item The Kneser-type solution is the conformal sewing of two half-cylinders
  along $T_b$.  A contraction in a weighted Wiener algebra gives it with
  explicit bounds (Theorem~\ref{thm:W2}).  A proper-map argument proves
  uniqueness in the class of Definition~\ref{def:class}
  (Theorem~\ref{thm:char}).
\item Kneser's initial region can be tilted continuously around the cusp.
  In the coordinate $\log\frac{u-u_1}{u-u_2}$ it is a thin strip on which the map is a
  translation up to relative error $O(|\varepsilon|)$, so its quotient is uniformly
  quasiconformal to a flat cylinder, whatever the multiplier.  The
  Ahlfors--Bers theorem gives holomorphic dependence on the base, and at
  $b<\eta$ the tilted region is a fundamental domain of the sewn surface
  (Section~\ref{sec:cusp}).
\end{enumerate}

\subsection*{Organisation}
Section~\ref{sec:setup} collects the parametrisation and the parabolic data.
Section~\ref{sec:transition} introduces $T_b$, the welding identity and the
sewing construction.  Section~\ref{sec:proofC} proves Theorem~\ref{mthm:A}, and
Section~\ref{sec:first} proves Theorem~\ref{mthm:D}.
Sections~\ref{sec:W2}--\ref{sec:char} prove Theorem~\ref{mthm:B} (existence
with Wiener bounds, uniqueness) and its base dependence, and Section~\ref{sec:cusp}
proves Theorem~\ref{mthm:F}, including the crossing away from the cusp
(Section~\ref{sec:far}).  Section~\ref{sec:general}, placed before it, states
Theorems~\ref{mthm:A}, \ref{mthm:B} and~\ref{mthm:D} for general real saddle-node
unfoldings and proves that the moduli $|\tau_n|$ do not depend on the base point.
Section~\ref{sec:numerics} reports the numerics, and
Section~\ref{sec:open} lists open problems.

\begin{remark}[On the definition and uniqueness in~\cite{Paulsen2019}]\label{rem:paulsen}
In~\cite{Paulsen2019}, $\kappa_b$ is defined by analytic continuation in $b$ of
Kneser's solution $\kappa_b(z_0)$, $\operatorname{Re}z_0>-2$, from $(\eta,\infty)$
(Prop.~3 there).  That proposition asserts real-analytic dependence of
Kneser's conformal uniformisation on the real base without a separate
parameter-dependence proof; Proposition~\ref{prop:exterior-real-parameter}
supplies such a proof locally along the real exterior ray.  Near the cusp, Theorem~\ref{mthm:F} carries the family through
the Shell--Thron boundary by a different route.  Away from the cusp,
Theorems~\ref{thm:far} and~\ref{thm:last-arc} cross the boundary arcs with
Kneser's untilted chord.  Off that boundary, Lemma~1 of
that paper asserts two
regular representations $\psi_1^{-1}(\rho_1)$ and
$\psi_2^{-1}(\rho_2)$ from the limits at $\pm\ii\infty$.
We impose their degree-one form explicitly in
Definition~\ref{def:class}.  A fixed-point limit alone leaves an integer
end label open: near an end one only gets
$\sigma(F(z))=\eee^{\ell z}q^mh(q)$ with $q=\eee^{\pm2\pi\ii z}$, $h(0)\ne0$ and some
integer $m\ge0$, and only $m=0$ is a time change $\mathrm{id}+(\text{bounded
periodic})$.  The continuation across the Shell--Thron boundary is
argued in~\cite[\S5]{Paulsen2019} by the continuity of a numerical scheme; it is
not proved there.

The uniqueness statement~\cite[Prop.~2]{Paulsen2019} requires only
analyticity on $\operatorname{Re}z>-2$, $F(0)=1$ and the two limits.  As stated it is false for
every real $b>\eta$.  Let $K=K_b$ be Kneser's solution, $a=\log b>1/\eee$, and
$u_n=K(n)\to\infty$.
\begin{itemize}
\item \emph{Large image discs.}  Choose $N$ with $au_{N+1}>1$, and $\delta>0$ with
  $B(u_N,\delta)\subset K(B(N,\tfrac12))$ and $u_{N+1}(1-\eee^{-a\delta})>\delta$.  The inclusion
  $E_b(B(t,\delta))\supset B(\eee^{at},\eee^{at}(1-\eee^{-a\delta}))$ (apply the principal logarithm and
  $|\log(1+v)|\le-\log(1-|v|)$) propagates these discs.  So $K(B(n+1,\tfrac12))$ contains a
  disc about $u_{n+1}$ of radius $u_{n+1}(1-\eee^{-a\delta})\to\infty$.
\item \emph{A shifted solution.}  For large $n$ and $k=\lceil a\exp(au_{n+1})/(2\pi)\rceil$, the point
  $w_k=a^{-1}(\log(2\pi k/a)+\ii\pi/2)$ lies in that disc, and $E_b(w_k)=2\pi\ii k/a$,
  $E_b^2(w_k)=1$.  Choose $z_k\in B(n+1,\tfrac12)$ with $K(z_k)=w_k$ and put
  $F_k(z)=K(z+z_k+2)$.  Then $F_k$ is analytic on $\operatorname{Re}z>-2$, has $F_k(0)=1$ and the
  same vertical limits as $K$, but $F_k(-1)=2\pi\ii k/a\ne0=K(-1)$.
\end{itemize}
Since $k$ is unbounded, this gives infinitely many solutions.
The gap in~\cite[Prop.~2]{Paulsen2019} is the last step of its proof,
which infers $C=0$ from $F_1(C)=1$.  For real $b>\eta$ the uniqueness
theorem of~\cite{PaulsenCowgill2017} is not affected: it also assumes
analyticity on $\mathbb C\setminus(-\infty,-2]$ and $F(\bar z)=\overline{F(z)}$,
both of which $F_k$ violates.  Indeed $z_k$ is nonreal, because $K$ is real on
the real axis and $w_k$ is not.  And $K$ cannot extend holomorphically through
$-2$, since $K(-1)=0$ whereas an exponential never vanishes.  Hence $F_k$ is
singular at $-4-z_k$, off the standard cut.
For complex bases, some root-selection
condition must be added to~\cite[Prop.~2]{Paulsen2019}.  The seam-height
condition of Definition~\ref{def:class} is such a condition, and
Theorem~\ref{thm:char} shows that it suffices near $(1,\eta)$.
\end{remark}

\section{Setup and parabolic data}\label{sec:setup}

Throughout, $1<b<\eta$ and $E(w)=b^w$.

\begin{lemma}[Parametrisation]\label{lem:param}
Let $L$ be a fixed point of $E$ with multiplier $\lambda=L\log b$.  Then
$L=\eee^\lambda$ and $\log b=\lambda\eee^{-\lambda}$.  The map
$\lambda\mapsto\lambda\eee^{-\lambda}$ increases on $(0,1)$ and decreases on
$(1,\infty)$, with maximum $\eee^{-1}=\log\eta$.  Hence $\lambda\in(0,1)$ and
$\lambda_2\in(1,\infty)$ parametrise $(1,\eta)$ through the attracting and the
repelling fixed point, and the two branches meet at $b=\eta$.
\end{lemma}

\begin{proof}
$L=b^L$ gives $\log L=L\log b=\lambda$.
\end{proof}

\begin{proposition}\label{prop:saddle}
With $\lambda=1-\varepsilon$ one has
$\varepsilon=\sqrt{2\eee(\log\eta-\log b)}\,(1+O(\varepsilon))$, and as $b\to\eta^-$
\[
 \Lambda=\exp\Bigl(-\frac{C_\eta}{\sqrt{\eta-b}}(1+o(1))\Bigr),\qquad
 C_\eta=\frac{4\pi^2}{\sqrt{2\eee^{1-1/\eee}}}.
\]
\end{proposition}

\begin{proof}
Expand $\lambda\eee^{-\lambda}=\eee^{-1}(1-\varepsilon^2/2-\varepsilon^3/3+O(\varepsilon^4))$ and use
$|\log\lambda|=\varepsilon+O(\varepsilon^2)$ and $\log\eta-\log b\sim(\eta-b)/\eta$.
\end{proof}

\begin{remark}[Translation invariance]\label{rem:invariance}
A simultaneous shift $\Kk(\cdot+\delta)$, $\Rr(\cdot+\delta)$ sends
$c_n\mapsto c_n\eee^{2\pi\ii n\delta}$.  Thus $\kappa_n=c_n/c_1^n$, when
defined, is invariant under complex translations, while $|c_1|$ is invariant
only under real ones.
The normalisation $F(0)=1$ fixes the imaginary origin.  This is why
normalised first-mode limits require a specified base point, and why the
constant $a$ enters Corollary~\ref{cor:limits}.
\end{remark}

\subsection{The parabolic germ}\label{sec:horn}
At $b=\eta$, put $w=\eee(1+u)$.  Then $E$ becomes $f(u)=\eee^u-1$, with attracting
petal on $\operatorname{Re}u<0$ and formal Abel function
\begin{equation}\label{eq:alpha}
 \alpha(u)=-\frac2u+\frac{\log(-u)}3+\sum_{k\ge1}c^{\alpha}_ku^k,\qquad
 \alpha(f(u))=\alpha(u)+1 .
\end{equation}
The Fatou coordinates are
$\Phi_{\mathrm{att}}(u)=\lim[\alpha(f^nu)-n]$ and
$\Phi_{\mathrm{rep}}(u)=\lim[\alpha(f^{-n}u)+n]$, with $\log(-u)$ continued along
orbits.  On the upper gate, the horn map
$\Phi_{\mathrm{rep}}\circ\Phi_{\mathrm{att}}^{-1}(z)=z+\sum_{n\ge1}A_n\eee^{2\pi\ii nz}$
and its inverse
$\Phi_{\mathrm{att}}\circ\Phi_{\mathrm{rep}}^{-1}(z)=z+\sum_{n\ge1}B_n\eee^{2\pi\ii nz}$
have no constant mode (Section~\ref{sec:proofC}, item (d)).  Formal inversion
gives
\begin{equation}\label{eq:Bn}
 B_1=-A_1,\qquad B_2=-A_2+2\pi\ii A_1^2 .
\end{equation}
\begin{proposition}\label{prop:horn-nonzero}
The first coefficient of the inverse upper horn map satisfies $B_1\ne0$.
\end{proposition}
\begin{proof}
Let $A$ be the immediate parabolic basin of $f(u)=\eee^u-1$.  The negative
real ray is contained in $A$, since $u<f(u)<0$ for $u<0$.  The entire map
$f$ has no critical points and exactly one finite singular value, $-1$,
its asymptotic value along that ray.  By~\cite[Lemma~87]{Cheritat2022},
the singular values of $f|_A:A\to A$ are contained in $A\cap\{-1\}$;
the negative ray shows that $-1$ does occur.  Hence $f$ belongs to the
one-petal unisingular class $S_\infty$ of~\cite[Definition~8]{Cheritat2022}.
The parabolic-renormalisation theorem~\cite[Theorem~9 and
Corollary~11]{Cheritat2022} gives a renormalisation $R[f]$ with exactly
one attracting petal at its parabolic fixed point.  The upper
renormalisation is the exponential projection of
$\Phi_{\mathrm{att}}\circ\Phi_{\mathrm{rep}}^{-1}$, up to nonzero linear
scalings of source and target~\cite[Appendix~A]{Cheritat2022}.  In our
normalisation, writing $\zeta=\eee^{2\pi\ii z}$, this projection is
$H(\zeta)=\zeta+2\pi\ii B_1\zeta^2+O(\zeta^3)$.  Since $H'(0)=1$,
nonzero linear scalings that again give multiplier one preserve whether
the quadratic coefficient vanishes.  One attracting petal means that
this coefficient is nonzero.  Therefore $B_1\ne0$.
\end{proof}

\begin{remark}[The mechanism]\label{rem:horn-mechanism}
The underlying reason is the classical one.  The map $f$ is of finite type
in the sense of Epstein: it has no critical points and a single finite
asymptotic value.  Its horn maps are again of finite type, with singular
values in $\mathbb C^*$ given by the projections of the single singular
orbit~\cite{Epstein1993}.  So the upper horn map has exactly one singular
value in $\mathbb C^*$.  By the Fatou-type theorem for finite-type
maps~\cite{Epstein1993}, the immediate basin of each attracting petal of a
parabolic fixed point contains a singular value other than the fixed point.
If $B_1=0$, then $H(\zeta)-\zeta=O(\zeta^3)$, and $H$ would have at least two
attracting petals at $0$, each needing its own singular value.
\end{remark}
Computing both maps directly on horizontal lines gives
\begin{equation}\label{eq:B12}
\begin{gathered}
 |B_1|=0.089058436412213331563,\qquad |B_2|=0.012487384111564700811,\\
 \kappa_2^\infty:=\frac{B_2}{B_1^2}=-0.02532049309975755123+1.5742190652319516267\,\ii,\\
 \frac{B_3}{B_1^3}=-2.254799124523803111-0.024408714296883171\,\ii .
\end{gathered}
\end{equation}
Also $a=\Phi_{\mathrm{att}}(1/\eee-1)=3.0292972144180360989\ldots$, so the limiting
phase in Corollary~\ref{cor:limits} is $\arg B_1+2\pi a=1.5865330757\ldots$
(mod $2\pi$).  These values are stable to all digits shown under changes of working
precision ($40$--$60$ digits), transform size ($N=16,24,32$) and sampling
height ($\operatorname{Im}z=1,1.5,2$).  Digits beyond the intervals below
remain numerical.  The first three horn coefficients have independent certified
enclosures of Proposition~\ref{prop:horn-intervals};
nonvanishing of $B_1$ also follows from Proposition~\ref{prop:horn-nonzero}.
The base-point coordinate $a$ has the certified enclosure
\[
 3.0292972144180360989249938<a<3.0292972144180360989249940.
\]
(Appendix~\ref{app:horn}).

\begin{proposition}[Certified first horn coefficients]\label{prop:horn-intervals}
The Fourier coefficients of the inverse upper horn map satisfy
\[
 0.0890584364122133315631<|B_1|<0.0890584364122133315633,
 \qquad
 0.01248738411156469<|B_2|<0.01248738411156471,
 \qquad
 0.00159278992<|B_3|<0.00159278994.
\]
The ratios used in the numerical comparisons also satisfy
\[
\begin{aligned}
 -0.025320493100&<\operatorname{Re}(B_2/B_1^2)<-0.025320493099,&
 1.574219065231&<\operatorname{Im}(B_2/B_1^2)<1.574219065233,\\
 -2.25479914&<\operatorname{Re}(B_3/B_1^3)<-2.25479910,&
 -0.02440873&<\operatorname{Im}(B_3/B_1^3)<-0.02440869.
\end{aligned}
\]
\end{proposition}

\begin{proof}
Appendix~\ref{app:horn}.
\end{proof}

\begin{remark}[Formal invariants]\label{rem:formal}
Put $\zeta=\eee^{2\pi\ii z}$.  The inverse horn map becomes the cylinder germ
$F(\zeta)=\zeta\exp(2\pi\ii\sum B_n\zeta^n)=\zeta+a_2\zeta^2+a_3\zeta^3+\cdots$.
Since $B_1\ne0$, its formal invariants under tangent-to-identity conjugacy
include $a_2=2\pi\ii B_1$ and the iterative residue
$\rho=1-a_3/a_2^2$~\cite{BuffEpstein2002}.  One has
\begin{equation}\label{eq:kapparho}
 \kappa_2^\infty=2\pi\ii\bigl(\tfrac12-\rho\bigr),\qquad
 \rho=0.249455254258832-0.004029881638351\,\ii .
\end{equation}
When defined, the higher ratios $B_n/B_1^n$ are not formal invariants of $F$.  They are,
however, invariants under the translations that preserve the normalisation of
Fatou coordinates, and Theorem~\ref{mthm:A} compares them for every $n$.
\end{remark}

\section{The separation is a transition-map invariant}\label{sec:transition}

The separation problem compares a Kneser-type solution $\Kk_b$ with $\Rr_b$.
This section removes $\Kk_b$ from the problem; Section~\ref{sec:cusp}
identifies it with the sewn solution constructed below.  The separation coefficients turn out to be determined, up to a
relative error of order $\Lambda$, by the transition map between the two
\emph{regular} solutions at the two real fixed points --- an object built from
two convergent Schr\"oder series in a fraction of a second, at any $b<\eta$.

\subsection{The transition map}
Let $\tau$ be the Koenigs coordinate at the repelling point $L_2$,
$\tau(E(w))=\lambda_2\tau(w)$, $\tau'(L_2)=1$, and put
\[
 S(z)=\tau^{-1}\bigl(-\lambda_2^{\,z}\bigr),
\]
continued to an entire function by $S(z+1)=E(S(z))$.  It maps $\mathbb R$
decreasingly onto the segment $(L,L_2)$ and has period $2\pi\ii/\log\lambda_2$.
Let $\Rr^{-1}(w)=\log(\sigma(w)/\sigma(1))/\log\lambda$ be the attracting Abel
coordinate, $\Rr^{-1}(1)=0$, which is multivalued with period
$\ii h=2\pi\ii/|\log\lambda|$.  On the horizontal strip where $S$ takes values
in the immediate basin of $L$ and a continuous branch is chosen, put
\[
 T(z):=\Rr^{-1}(S(z)),\qquad T(z+1)=T(z)+1,\qquad
 T(z)-z=\sum_{n\in\mathbb Z}t_n\eee^{2\pi\ii nz}.
\]

\begin{lemma}[Real structure]\label{lem:realT}
For real $b\in(1,\eta)$, $\operatorname{Im}T(x)\equiv h/2\pmod h$ for every
real $x$.  Consequently, on the branch with $\operatorname{Im}t_0=h/2$,
$T(x)-x-\ii h/2$ is real, $t_{-n}=\overline{t_n}$, and the line
$\operatorname{Im}w=-h/2$ of attracting time is the image of $\mathbb R$ under
the branch $T-\ii h$.
\end{lemma}

\begin{proof}
$S(x)\in(L,L_2)$, where $\sigma>0$, while $1<L$ gives $\sigma(1)<0$.  Hence
$\sigma(S(x))/\sigma(1)$ is negative and
$\operatorname{Im}\log(\sigma(S(x))/\sigma(1))\in\pi+2\pi\mathbb Z$; divide
by $\log\lambda=-2\pi/h$.
\end{proof}

Translations of $S$ ($S\mapsto S(\cdot+\delta)$, $\delta\in\mathbb C$) send
$t_0\mapsto t_0+\delta$, $t_n\mapsto t_n\eee^{2\pi\ii n\delta}$, so
\begin{equation}\label{eq:taun}
 \tau_n(b):=t_n\,\eee^{-2\pi\ii n t_0}\qquad(\operatorname{Im}t_0=h/2)
\end{equation}
does not depend on the normalisation of $S$; it depends on that of $\Rr$
only through the real translation fixed by $\Rr(0)=1$.  The condition
$\operatorname{Im}t_0=h/2$ refers to the real normalisation of $S$ and fixes the
branch of $\Rr^{-1}$; a branch shift by $\ii kh$ would multiply $\tau_n$ by
$\Lambda^{-nk}$.  The branch is fixed once in the real normalisation and then
carried along.

\subsection{The welding identity}

\begin{theorem}\label{thm:welding}
Let $\Kk$ be a solution of~\eqref{eq:abel}, holomorphic on the sets below
where it is represented, and suppose there are
$1$-periodic holomorphic maps $\varphi_\pm=\mathrm{id}+(\text{periodic})$ with
\begin{enumerate}
\item $\Kk=\Rr\circ\varphi_+$ on an upper half-plane $U_+$, where
  $P=\varphi_+-\mathrm{id}=\sum_{n\ge0}c_n\eee^{2\pi\ii nz}$ converges on $U_+$;
\item $\Kk=S\circ\varphi_-$ on a lower half-plane $U_-$, where $\varphi_-$ is
  univalent, $G:=\varphi_-^{-1}-\mathrm{id}$ is bounded on
  $V_-=\varphi_-(U_-)$, which contains a lower half-plane, and
  $G(w)\to g_0$ as $\operatorname{Im}w\to-\infty$;
\item $U_+\cap U_-$ contains a horizontal line $\ell$ with $\varphi_-(\ell)$
  inside the strip of $T$.
\end{enumerate}
Then, with the branch of $T$ fixed by continuity from $\varphi_+\circ\varphi_-^{-1}$,
for every $n\ge1$
\begin{equation}\label{eq:welding}
 t_n=\sum_{m\ge n}c_m\bigl[\eee^{2\pi\ii mG}\bigr]_{n-m},\qquad
 t_0=g_0+\sum_{m\ge0}c_m\bigl[\eee^{2\pi\ii mG}\bigr]_{-m},
\end{equation}
where $[F]_k$ is the $k$-th Fourier coefficient; in particular
$[\eee^{2\pi\ii mG}]_0=\eee^{2\pi\ii mg_0}$ and $[\eee^{2\pi\ii mG}]_k=0$ for $k>0$.
\end{theorem}

\begin{proof}
On $\ell$, $\Rr\circ\varphi_+=S\circ\varphi_-=\Rr\circ T\circ\varphi_-$, so
$\varphi_+-T\circ\varphi_-$ takes values in the discrete period lattice of
$\Rr$ and is constant; absorb it into the branch of $T$.  Hence
$T=\varphi_+\circ\varphi_-^{-1}$ on the curve $\varphi_-(\ell)$, i.e.\
$T(w)-w=G(w)+P(w+G(w))$ there.  The Fourier coefficients are computed by
Cauchy's theorem on this curve.  The region below it lies in $V_-$, and
$\sum c_m\eee^{2\pi\ii m\varphi_-^{-1}(w)}$ converges uniformly on the curve
$\varphi_-(\ell)$ itself, because $\ell\subset U_+$.  Since $G$ is bounded and holomorphic on $V_-$, with
limit $g_0$ at $-\ii\infty$, so is $\eee^{2\pi\ii mG}$, whose Fourier series
therefore contains only modes $k\le0$ and has constant term
$\eee^{2\pi\ii mg_0}$.  Expanding
$P(w+G)=\sum_mc_m\eee^{2\pi\ii mw}\eee^{2\pi\ii mG(w)}$ and collecting modes
gives~\eqref{eq:welding}; the Fourier coefficients of $T-\mathrm{id}$ do not depend
on the line inside its strip of holomorphy.
\end{proof}

\begin{corollary}\label{cor:cn}
Under the hypotheses of Theorem~\ref{thm:welding}, for every $n\ge1$,
\[
 t_n\,\eee^{-2\pi\ii ng_0}
 =c_n+\eee^{-2\pi\ii ng_0}
   \sum_{m>n}c_m[\eee^{2\pi\ii mG}]_{n-m}.
\]
If $c_n\ne0$, this may be written as
\[
 t_n\,\eee^{-2\pi\ii ng_0}=c_n\bigl(1+\rho'_n\bigr),\qquad
 |\rho'_n|\le\sum_{m>n}\frac{|c_m|}{|c_n|}\,
 \bigl|[\eee^{2\pi\ii mG}]_{n-m}\bigr|\,\bigl|\eee^{-2\pi\ii ng_0}\bigr|,
\]
and $g_0=t_0-c_0-\sum_{m\ge1}c_m[\eee^{2\pi\ii mG}]_{-m}$.  Heuristically, if the
$c_m$ decay like $\Lambda^m$ and $G$ is bounded on a line at height comparable to
$h$ above the gluing line, both corrections are $O(\Lambda)$.  The precise
statement for the sewing solution is Theorem~\ref{thm:W2}(4).  On the
branch of Lemma~\ref{lem:realT}, $g_0=t_0-\ii h+O(\Lambda)$, so that
for the sewing solution, by Theorem~\ref{thm:W2}(4),
\begin{equation}\label{eq:chat-tau}
 \hat c_n:=\frac{c_n}{\Lambda^n}=\tau_n(b)+O(\Lambda).
\end{equation}
Since $B_1\ne0$, also $c_1\ne0$ for $b$ sufficiently close to $\eta$ and
$\kappa_n:=c_n/c_1^{\,n}=t_n/t_1^{\,n}+O(\Lambda)$ for each fixed $n$.
\end{corollary}

The identity $\eee^{-2\pi\ii n(t_0-\ii h)}=\Lambda^n\eee^{-2\pi\ii nt_0}$
is where the scale $\Lambda$ enters: \emph{the $\Lambda$ law is the statement
that the gluing line sits at attracting time $-h/2$, half a period below the
real axis, while the gate carrying the positive modes of $T$ sits a further
half period below it.}  Nothing in the argument uses Kneser's construction.

\subsection{Welding construction}

\begin{proposition}\label{prop:weld}
For $b\in(1,\eta)$ glue the half-cylinders
$C_-=\{\operatorname{Im}w'\le0\}/\mathbb Z$ (repelling time) and
$C_+=\{\operatorname{Im}w\ge-h/2\}/\mathbb Z$ (attracting time) along their
boundaries by $w=T(w')-\ii h$.  The result is conformally $\mathbb C^*$; its
universal-cover coordinate $z$ defines a solution
\[
 \Kk^{W}(z)=\Rr(w(z))\ \text{on }C_+,\qquad \Kk^{W}(z)=S(w'(z))\ \text{on }C_-,
\]
unique up to a translation in $z$, which satisfies the hypotheses of
Theorem~\ref{thm:welding}.  It is holomorphic except where $w(z)$ meets the
singular set $\bigcup_{k}((-\infty,-2]+\ii kh)$ of $\Rr$.
\end{proposition}

\noindent
Unlike $S$, which is entire (the Poincar\'e function of a repelling point),
$\Rr$ is not: its inverse Koenigs map extends only through inverse branches of
$E$.  So $\Kk^W$ carries, besides the cut near $(-\infty,-2]$, the copies
near $(-\infty,-2]+\ii kh$, $k\ge1$, lying in its $\Rr$-represented region.
The same holds for the Kneser-type solutions at complex bases inside the
Shell--Thron region, whose upper representation is through the attracting
regular superfunction.

\begin{proof}
By Lemma~\ref{lem:realT} the gluing map is a real-analytic diffeomorphism
between the boundary circles that extends holomorphically to a neighbourhood,
so the glued surface $X$ is a Riemann surface; each end is a punctured disc,
hence $X\cup\{\text{two points}\}$ is a simply connected compact Riemann
surface and $X\cong\mathbb C^*$.  The definitions agree on the seam because
$\Rr\circ(T-\ii h)=S$.  The generator of $\pi_1(X)$ lifts to $z\mapsto z+1$
and acts as $w\mapsto w+1$, $w'\mapsto w'+1$, whence
$\Kk^W(z+1)=E(\Kk^W(z))$.  Near each puncture the change of coordinate is a
holomorphic map of punctured discs fixing the puncture, which is the
statement that $w-z$ (resp.\ $w'-z$) has only non-negative (resp.\
non-positive) Fourier modes.  Uniqueness is the uniqueness of the uniformisation.
\end{proof}

\begin{proposition}[Local parameter dependence of the sewn family]
\label{prop:weld-parameter}
For every $b_0\in(1,\eta)$, the normalised geometric sewing of
Proposition~\ref{prop:weld} extends to a holomorphic family over a complex
disc $\Delta_{b_0}$ about $b_0$.  On every common regular domain of the
universal-cover coordinate, $\Kk_b^W(z)$ is jointly holomorphic in $(b,z)$.
These local families agree on overlaps and therefore define the inner sewn
family on an open complex neighbourhood $N_W$ of the entire real interval
$(1,\eta)$, the union of the discs $\Delta_{b_0}$.  No uniform width of that
neighbourhood is asserted.
\end{proposition}

\begin{proof}
Fix $b_0$.  Both fixed points and their Koenigs coordinates vary
holomorphically near $b_0$.  The compact set
$S_{b_0}([0,1])\subset(L,L_2)$ lies in the attracting basin and avoids its
fixed point.  Thus the chosen branch of
$T_b=\Rr_b^{-1}\circ S_b$ is jointly holomorphic on a parameter disc times
a narrow strip about $[0,1]$; the relation $T_b(w'+1)=T_b(w')+1$
extends it to the whole strip.  The real circle map is orientation
preserving: $\tau_b'(S_b(x))>0$ by inverse iteration to $L_2$,
$\sigma_b'(S_b(x))>0$ by forward iteration to $L$, and hence
$S_b'(x)<0$ and
$(\Rr_b^{-1})'(S_b(x))=\sigma_b'(S_b(x))/
 (\sigma_b(S_b(x))\log\lambda_b)<0$.
Consequently $T_{b_0}'(x)>0$.  Compactness of the circle and local
injectivity give a smaller strip on which $T_{b_0}$ is injective modulo
$1$; shrinking the parameter disc preserves this property and makes the
image of the real circle an analytic Jordan curve in the cylinder.

Put $U_b=T_b+\varpi_b$, where
$\varpi_b=2\pi\ii/\log\lambda_b$; at $b_0$ this is $T_{b_0}-\ii h$.
Use $U_b$ to identify a thin open annular collar in the repelling
cylinder with its image in the attracting cylinder.  More explicitly,
for a sufficiently small $\delta>0$ take the repelling chart
$\{\operatorname{Im}w'<\delta/2\}/\mathbb Z$ and the attracting chart
above the curve $U_b(\{\operatorname{Im}w'=-\delta/2\})$; identify them
on $|\operatorname{Im}w'|<\delta/2$.  After shrinking the parameter disc,
these charts and the transition $(b,w')\mapsto(b,U_b(w'))$ form a
Hausdorff holomorphic family of cylinders.  At $b=b_0$ it is the
analytic-collar version of the boundary welding in
Proposition~\ref{prop:weld}.  Adjoin the two ends using the parameter
independent local coordinates $\eee^{2\pi\ii w}$ and
$\eee^{-2\pi\ii w'}$.  This yields a proper holomorphic submersion
whose fibres are compact Riemann surfaces of genus zero.  The two added
points and the upper-chart point $w=0$ are three disjoint holomorphic
sections (the last remains above the seam after shrinking the disc).

Local holomorphic triviality of a proper family with mutually
biholomorphic compact fibres~\cite{FischerGrauert1965}, followed by the
unique fibrewise M\"obius map sending these three sections to $0,\infty,1$,
gives a jointly holomorphic coordinate $q_b$ on the compactified fibres.
Choose the local lift $z=(2\pi\ii)^{-1}\log q_b$ for which the marked
upper point $w=0$ has $z=0$.  On the lift of the overlap,
$\Rr_b(U_b(w'))=S_b(w')$ because $\varpi_b$ is a period of $\Rr_b$.
The two regular superfunctions therefore define $\Kk_b^W$, jointly
holomorphic wherever their represented branches are regular, and
$\Kk_b^W(0)=\Rr_b(0)=1$.  On the real parameter slice this is exactly
the geometric welding.  Its marked coordinate makes the local families
unique, so the identity theorem patches them on overlapping discs.
\end{proof}


For non-real $b$ sufficiently close to any fixed real $b_0\in(1,\eta)$,
the sewn solution thus depends holomorphically on $b$.
Theorem~\ref{thm:W1} gives the stronger quantitative Wiener bounds near
$\eta$.  For $\operatorname{Im}b>0$ near $(b_1,\eta)$ it tends (Theorem~\ref{thm:W1}(2)) to $L$ and $L_2$ as
$\operatorname{Im}z\to\pm\infty$, respectively.  Those limits alone do not characterise
it: Proposition~2 of~\cite{Paulsen2019} has the counterexample in
Remark~\ref{rem:paulsen}.  The base-continued Kneser family is identified with $\Kk^W$ in
Theorem~\ref{mthm:F}; the seam condition of Definition~\ref{def:class} is what
selects the right root there.  Table~\ref{tab:tau} gives numerical evidence:
$\tau_1$ and $t_2/t_1^2$ reproduce the ladder values of $\hat c_1$ and
$\kappa_2$ --- moduli \emph{and} arguments --- to the precision of the
ladder extrapolation, e.g.\ eight digits at $b=1.10$.

\begin{table}[htbp]
\centering
\caption{Kneser ladders (Sections~\ref{sec:invariant},~\ref{sec:mode2}, extrapolated
to $\operatorname{Im}b=0$) against the Kneser-free invariants
$\tau_1=t_1\eee^{-2\pi\ii t_0}$ and $t_2/t_1^2$ (script \texttt{transition\_map.py}).}
\label{tab:tau}
\small\setlength{\tabcolsep}{3.5pt}
\begin{tabular}{lllll}
\toprule
$b$ & $\hat c_1$ (ladder) & $\tau_1$ & $\kappa_2$ (ladder) & $t_2/t_1^2$\\
\midrule
$1.05$ & $0.0850235\,\eee^{1.534136\ii}$ & $0.0850226\,\eee^{1.534165\ii}$ & $1.374049\,\eee^{1.413420\ii}$ & $1.374336\,\eee^{1.413645\ii}$\\
$1.10$ & $0.08642720\,\eee^{1.55109214\ii}$ & $0.08642720\,\eee^{1.55109214\ii}$ & $1.440851\,\eee^{1.477690\ii}$ & $1.440851\,\eee^{1.477690\ii}$\\
$1.15$ & $0.0872053\,\eee^{1.561012\ii}$ & $0.0872052\,\eee^{1.561013\ii}$ & $1.479111\,\eee^{1.511255\ii}$ & $1.479107\,\eee^{1.511258\ii}$\\
$1.30$ & $0.088434$ & $0.0884321\,\eee^{1.577580\ii}$ & $1.541516\,\eee^{1.561918\ii}$ & $1.541515\,\eee^{1.561919\ii}$\\
\bottomrule
\end{tabular}
\end{table}

\subsection{Reduction to confluence}

Write $a:=\Phi_{\mathrm{att}}(1/\eee-1)=3.0292972144180360989\ldots$ for the
$\alpha$-normalised attracting Fatou coordinate of the base point $w=1$ at
$b=\eta$ (real, since the orbit of $u=1/\eee-1$ stays on the negative axis).

\begin{theorem}\label{thm:reduction}
Assume the following confluence statement \textup{(C)}: there are constants
$s(b)\in\mathbb C$ such that, as $b\to\eta^-$,
\[
 \Rr_b^{-1}\bigl(\eee(1+u)\bigr)\to\Phi_{\mathrm{att}}(u)-a,\qquad
 S_b^{-1}\bigl(\eee(1+u)\bigr)+s(b)\to\Phi_{\mathrm{rep}}(u),
\]
uniformly on compact subsets of the upper gate, with the branches continued as
in Section~\ref{sec:horn}.  Then for every $n\ge1$
\[
 \tau_n(b)\longrightarrow B_n\,\eee^{2\pi\ii na}.
\]
Consequently, also
$t_n/t_1^{\,n}\to B_n/B_1^{\,n}$ for every $n\ge1$.
Together with Corollary~\ref{cor:cn} and the identification of
$\Kk_b$ with $\Kk^W_b$ in Theorem~\ref{thm:cusp}, this gives
$\hat c_n\to B_n\eee^{2\pi\ii na}$ for all $n$.  The limiting phase
$\arg\hat c_1\to\arg B_1+2\pi a=1.5865330757\ldots\pmod{2\pi}$.
\end{theorem}

\begin{proof}
Replace $S_b$ by $S_b(\cdot-s(b))$, which leaves $\tau_n$ unchanged.  By (C),
$T_b\to\Phi_{\mathrm{att}}\circ\Phi_{\mathrm{rep}}^{-1}-a=h^{-1}-a$ uniformly
on $[0,1]+\ii Y_g$ for fixed $Y_g$ in the gate, so the Fourier coefficients on
that line converge: $t_n\to B_n$ ($n\ge1$) and $t_0\to-a$, the latter because
$B_0=0$ in the $\alpha$ normalisation.  Hence
$\tau_n=t_n\eee^{-2\pi\ii nt_0}\to B_n\eee^{2\pi\ii na}$.
\end{proof}

Statement (C) is a confluence theorem for the saddle-node unfolding
$u\mapsto \eee^{\lambda\eee^{-\lambda}\eee(1+u)}/\eee-1$ in the direction where the
two fixed points are real (the ``Glutsyuk'' direction): the Koenigs
coordinates at the two hyperbolic points converge, on the gate, to the
sectorial Fatou coordinates of the parabolic germ.  Results of this type are
due to Glutsyuk~\cite{Glutsyuk2001} and, for the full unfolded modulus,
Marde\v{s}i\'c--Roussarie--Rousseau~\cite{MRR2004} and
Christopher--Rousseau~\cite{ChristopherRousseau}.  Rather than matching their
normalisations, we prove (C) directly for this family in
Section~\ref{sec:proofC}, so Theorem~\ref{thm:reduction} holds unconditionally.

\section{Proof of Theorem~\ref{mthm:A}: confluence}\label{sec:proofC}

We prove (C) directly for the family at hand; no normalisation from the
literature is needed.  In $u$-coordinates ($w=\eee(1+u)$, $\mu=\eee\log b$)
\[
 f_b(u)=\eee^{\mu-1+\mu u}-1,\qquad f_\eta(u)=\eee^u-1 .
\]
For $b<\eta$ close to $\eta$ the map has, in a fixed disc $D_r$, exactly the two
fixed points $u_1=\eee^{\lambda-1}-1<0<u_2=\eee^{\lambda_2-1}-1$ (Rouch\'e), and
\begin{equation}\label{eq:prepared}
 f_b(u)-u=(u-u_1)(u-u_2)\,k_b(u),\qquad k_b\to k_\eta,\quad
 k_\eta(u)=\tfrac12+\tfrac u6+\cdots,
\end{equation}
uniformly on $D_r$, with $k_b$ zero-free there.  Put $A_j=1/\log\lambda_j$,
so $A_1<0<A_2$.

\begin{lemma}[Residues]\label{lem:res}
As $b\to\eta^-$: $A_1(u_1-u_2)\to2$, $A_1+A_2\to\nu:=1-4k_\eta'(0)=\tfrac13$, hence
$A_1u_1+A_2u_2\to2$ and $A_ju_j^2\to0$.
\end{lemma}

\begin{proof}
$\lambda_1=1+(u_1-u_2)k_b(u_1)$ and $\lambda_2=1+(u_2-u_1)k_b(u_2)$.  With
$d=u_1-u_2\to0$ and $1/\log(1+x)=1/x+\tfrac12+O(x)$,
$A_1+A_2=\frac{k_b(u_2)-k_b(u_1)}{d\,k_b(u_1)k_b(u_2)}+1+O(d)\to1-4k_\eta'(0)$,
and $A_1d\to1/k_\eta(0)=2$.
\end{proof}

\noindent
Numerically $A_1+A_2=0.3333331469$ and $A_1(u_1-u_2)=1.99676$ at $\lambda=0.99$.
Define the \emph{model time}
\[
 \Psi_b(u)=A_1\log(u-u_1)+A_2\log(u-u_2),
\]
which is exactly the time of the vector field $(u-u_1)(u-u_2)/(\ldots)$ whose
time-one map has multipliers $\lambda_1,\lambda_2$, and let
$\alpha_2(u)=-2/u+\tfrac13\log(-u)$.  By Lemma~\ref{lem:res}, for $u,p$ in a
simply connected region at positive distance from $0$, with logarithms
continued along the same path,
\begin{equation}\label{eq:psilimit}
 \Psi_b(u)-\Psi_b(p)\longrightarrow\alpha_2(u)-\alpha_2(p).
\end{equation}

\begin{lemma}[Uniform one-step error]\label{lem:F}
$F_b:=\Psi_b\circ f_b-\Psi_b-1$ is single-valued and holomorphic on a fixed disc
$D_{r'}$, $F_b\to F_\eta:=\alpha_2\circ f_\eta-\alpha_2-1$ uniformly there, and
\[
 |F_b(u)|\le M\,|u-u_1|\,|u-u_2|\qquad(u\in D_{r'/2})
\]
with $M$ independent of $b$.
\end{lemma}

\begin{proof}
By~\eqref{eq:prepared}, $(f(u)-u_1)/(u-u_1)=1+(u-u_2)k(u)$ and
$(f(u)-u_2)/(u-u_2)=1+(u-u_1)k(u)$, so
\[
 F=A_1\log\frac{1+(u-u_2)k}{1+(u-u_1)k}+(A_1+A_2)\log\bigl(1+(u-u_1)k\bigr)-1,
\]
holomorphic for $|u|<r'$ small.  The first logarithm is $O(|u_1-u_2|)$
uniformly, and $A_1(u_1-u_2)$, $A_1+A_2$ are bounded (Lemma~\ref{lem:res}); so
$|F|\le M_0$ on $|u|=r'$ uniformly, and the displayed formula converges.  At
$u=u_1$ the first term is $A_1\log\lambda_1=1$ and the second vanishes, so
$F(u_1)=0$; likewise $F(u_2)=0$.  The maximum principle applied to
$F/((u-u_1)(u-u_2))$ on $D_{r'}$ gives the bound, and letting $b\to\eta$ gives
$F_\eta=O(u^2)$.
\end{proof}

\begin{lemma}[Uniform petals and summability]\label{lem:petal}
Put $Z_b(u)=A_1\log\frac{u-u_1}{u-u_2}$ and
$\Omega_{R,b}=\{u:\operatorname{Re}Z_b(u)>R\}$ (a disc of Apollonius about
$u_1$).  There is $R_0$ such that for $R\ge R_0$ and all $b$ near $\eta$:
$\Omega_{R,b}\subset D_{r'/2}$, $f_b(\Omega_{R,b})\subset\Omega_{R,b}$,
$\operatorname{Re}Z_b(f_b^ku)\ge\operatorname{Re}Z_b(u)+\tfrac34k$, and
\[
 \sum_{k\ge N}|f_b^ku-u_1|\,|f_b^ku-u_2|\le\frac{C}{R+\tfrac34N}
 \qquad(u\in\Omega_{R,b}),
\]
with $C$ independent of $b$ and $u$.  As $b\to\eta$, $\Omega_{R,b}$ converges on
compacta to the petal $\{\operatorname{Re}(-2/u)>R\}$.
\end{lemma}

\begin{proof}
$Z(f(u))-Z(u)=A_1\log\frac{1+(u-u_2)k}{1+(u-u_1)k}=1+O(|u|+|u_1|+|u_2|)$
uniformly, by the expansion in the proof of Lemma~\ref{lem:F}; on
$\Omega_{R,b}\subset D_{C/R+|u_1|}$ this is within $\tfrac14$ of $1$ once $R$ is large,
which gives invariance and the growth of $\operatorname{Re}Z$.  With
$q=(u-u_1)/(u-u_2)$, $|q|=\eee^{-\sigma}$, $\sigma=|\log\lambda_1|\operatorname{Re}Z$,
one has $u-u_2=(u_2-u_1)/(q-1)$, whence
\[
 |u-u_1||u-u_2|=\frac{|q|\,|u_2-u_1|^2}{|1-q|^2}
 \le\Bigl(\frac{|u_2-u_1|}{|\log\lambda_1|}\Bigr)^2
 \frac{(1+\sigma)^2\eee^{-\sigma}}{(\operatorname{Re}Z)^2},
\]
using $1-\eee^{-\sigma}\ge\sigma/(1+\sigma)$.  The first factor is bounded
(Lemma~\ref{lem:res}) and $(1+\sigma)^2\eee^{-\sigma}\le4$; summing
$(R+\tfrac34k)^{-2}$ over $k\ge N$ gives the bound.  The last assertion follows
from $\operatorname{Re}Z_b(u)\to\operatorname{Re}(-2/u)$, i.e.\
$A_1(u_2-u_1)\to-2$.
\end{proof}

\begin{theorem}[Confluence]\label{thm:C}
Statement \textup{(C)} holds.  More precisely, with branches continued along
paths in the respective parabolic basins:
\begin{enumerate}
\item $\Rr_b^{-1}(\eee(1+u))\to\Phi_{\mathrm{att}}(u)-a$ uniformly on compact
  subsets of the parabolic basin of $f_\eta$;
\item for a fixed $p'$ in the gate (below), putting
  $s(b)=\Phi_{\mathrm{rep}}(p')-\mathcal S_b(p')$,
  $\mathcal S_b(u)+s(b)\to\Phi_{\mathrm{rep}}(u)$ uniformly on compact
  subsets of the repelling basin (backward orbits tending to $0$).  Here
  $\mathcal S_b$ is the real-normalised repelling time, i.e.\ the branch of
  $S_b^{-1}$ that is real on the segment $(u_1,u_2)$ and is continued through
  the band of Lemma~\ref{lem:band}, then along backward-invariant paths.
\end{enumerate}
The \emph{gate} of height $Y_*$ is $\mathcal G(Y_*)=\Phi_{\mathrm{rep}}^{-1}
(\{Y_*\le\operatorname{Im}z\le Y_*+1,\ -1\le\operatorname{Re}z\le2\})$.  For $Y_*$ large it is a compact
subset of both parabolic basins of $f_\eta$.  For $b$ near $\eta$ it lies in
the band $\mathcal B_b(Y_*-2)$, because $Z_b\to-2/u$ uniformly on it
(Lemma~\ref{lem:res}) and $\operatorname{Im}(-2/u)=\operatorname{Im}\alpha(u)+O(1)
=\operatorname{Im}\Phi_{\mathrm{rep}}(u)+O(1)$ there.
\end{theorem}

\begin{proof}
(1) On $\Omega_{R,b}$ put
$\mathcal A_b(u)=\lim_n[\Psi_b(f_b^nu)-n]=\Psi_b(u)+\sum_{k\ge0}F_b(f_b^ku)$;
the series converges absolutely by Lemmas~\ref{lem:F} and~\ref{lem:petal}.
Since $f^nu\to u_1$ and the Koenigs map satisfies
$\sigma(f^nu)=\lambda_1^n\sigma(u)$,
$\Psi_b(f^nu)-n=A_1\log\sigma(u)+A_2\log(u_1-u_2)+o(1)$, because
$A_1\log\lambda_1=1$.  Thus $\mathcal A_b$ is the Koenigs Abel coordinate up to
an additive constant, and
$\Rr_b^{-1}(\eee(1+u))=\mathcal A_b(u)-\mathcal A_b(u^*)$, $u^*=1/\eee-1$.
Let $K$ be a compact subset of the petal
$P_R=\{\operatorname{Re}(-2/u)>R\}$ and $U\supset K\cup\{f_\eta^{N_0}u^*\}$ a
simply connected region in $P_R$ at positive distance from $0$.  For $b$ near
$\eta$, $U\subset\Omega_{R,b}$ and $u_1,u_2\notin U$.  For $u,p\in U$,
\[
 \mathcal A_b(u)-\mathcal A_b(p)=\Psi_b(u)-\Psi_b(p)
 +\sum_{k\ge0}\bigl[F_b(f_b^ku)-F_b(f_b^kp)\bigr].
\]
The first term converges by~\eqref{eq:psilimit}.  In the sum, each term
converges because $f_b^k\to f_\eta^k$ and $F_b\to F_\eta$ uniformly.  The tails
beyond $N$ are $\le 2MC/(R+\tfrac34N)$ uniformly in $b$ (Lemma~\ref{lem:petal}).
Hence the difference converges uniformly to
$\alpha_2(u)-\alpha_2(p)+\sum_k[F_\eta(f_\eta^ku)-F_\eta(f_\eta^kp)]
 =\Phi_{\mathrm{att}}(u)-\Phi_{\mathrm{att}}(p)$.  The last identity holds
because $\Phi_{\mathrm{att}}=\lim[\alpha(f^nu)-n]$ and
$\alpha-\alpha_2=O(u)\to0$ along orbits.  A general compact $K$ of the basin
satisfies $f_\eta^{N}(K)\Subset P_R$ for some $N$.  Then
$f_b^N(K)\Subset P_R$ for $b$ near $\eta$, and
$\mathcal A_b=\mathcal A_b\circ f_b^N-N$.  Taking $p=f_b^{N_0}u^*$, which
tends to $f_\eta^{N_0}u^*\in U$, and using uniform convergence on a
neighbourhood gives (1).  (On $\Omega_{R,b}$ the function $\mathcal A_b$ is
multivalued, since this punctured Apollonius disc surrounds $u_1$.
It is single-valued on $U$, which
avoids $u_1$.  On $K$, $\sigma_b\ne0$ because $f_b^N(K)$ avoids $u_1$.  On
non-immediate components of the basin the branch is defined by
$\mathcal A_b\circ f_b^N-N$.  The Koenigs map is normalised by
$\sigma_b'(u_1)=1$.)

(2) Apply (1)'s argument to $G_b(v)=-f_b^{-1}(-v)$, the principal inverse
branch reflected.  Then $G_\eta(v)=-\log(1-v)=v+\tfrac{v^2}2+\tfrac{v^3}3+\cdots$,
which is of the same prepared form with $k_\eta'(0)=\tfrac13$ and residue
$-\tfrac13$.  The attracting fixed point of $G_b$ is $-u_2$, with multiplier
$1/\lambda_2$.  Its Koenigs Abel coordinate is $-S_b^{-1}(\eee(1-v))$ up to a
constant, and its parabolic Fatou coordinate is $-\Phi_{\mathrm{rep}}(-v)$ up
to a constant (since $\alpha_{f}(-v)=-\alpha_{G}(v)+\text{const}$).  The free
constant is fixed by $p'$.
\end{proof}

With Theorem~\ref{thm:C}, Theorem~\ref{thm:reduction} becomes unconditional.
Four points need checking; each is settled here.  The proof of
Lemma~\ref{lem:band} uses only Lemmas~\ref{lem:res}--\ref{lem:petal}, so there is
no circularity.

\emph{(a) Inversion on the gate line.}  Fix $Y_*\ge Y_B+6$ and the line
$\ell=[0,1]+\ii Y_*$.  By Rouch\'e, the solution $v_b(z)$ of
$\mathcal S_b(v)+s(b)=z$ near $\Phi_{\mathrm{rep}}^{-1}(z)$ exists and converges
uniformly on $\ell$.  Moreover $S_b(z-s(b))=\eee(1+v_b(z))$, because $\mathcal S_b$
is a branch of $S_b^{-1}$.

\emph{(b) The gate line lies in the strip of holomorphy of $T$.}  The points
$v_b(z)$, $z\in\ell$, lie within $o(1)$ of the gate, hence in the band $\mathcal B_b(Y_B+3)$.
By Lemma~\ref{lem:band}(ii) their $S$-times $z-s(b)$ satisfy
$|\operatorname{Im}(z-s(b))|\le h/2-Y_B-3+\pi/3+C/Y_B+C\varepsilon<h/2-Y_B-2$.  So the
shifted line $\ell-s(b)$ lies in the strip $\Sigma$ of Lemma~\ref{lem:band}(iii),
where $T$ is holomorphic.  Therefore the Fourier coefficients of
$\tilde T=T(\cdot-s)$ on $\ell$ are $t_n\eee^{-2\pi\ii ns}$ and
$\tilde t_0=t_0-s$, with the $t_n$ those of $T$ on $\mathbb R$.  No prior
knowledge of $\operatorname{Im}s$ is used.

\emph{(c) The branch of $\Rr_b^{-1}$.}  Continue $\Rr_b^{-1}$ from $u^*$ along
the negative axis to $f_\eta^{N_0}u^*$, where it is real.  Then go inside
$P_R\setminus D_\delta(0)$ through the upper half-plane to the gate, and then
inside the band down to the segment $(u_1,u_2)$.  This path passes above $u_1$
and does not wind around $u_2$.  Along it $\Rr_b^{-1}=\Psi_b+\sum_kF_b\circ f_b^k+c_R$, so the change of
$\operatorname{Im}\Rr_b^{-1}$ is
$A_1\Delta\arg(u-u_1)+A_2\Delta\arg(u-u_2)+O(C/Y+\varepsilon)=h/2+0+\text{small}$.
Lemma~\ref{lem:realT} forces the exact value, so on the segment
$\operatorname{Im}\Rr_b^{-1}=+h/2$.  The branch in
Theorem~\ref{thm:C}(1) is therefore the one with $\operatorname{Im}t_0=h/2$ used
in~\eqref{eq:taun}.

\emph{(d) $B_0=0$.}  At $b=\eta$, on the gate,
$\Phi_{\mathrm{att}}-\alpha_2=\sum_{k\ge0}F_\eta(f^ku)$ and
$\Phi_{\mathrm{rep}}-\alpha_2=-\sum_{k\ge1}F_\eta(f^{-k}u)$, with the same
principal $\log(-u)$ in $\alpha_2$.  That logarithm stays continuous along both
orbits, since $-f^{\pm k}u$ remains in the lower half-plane.  By
Lemma~\ref{lem:F} and $f^{\pm k}u\approx-2/(x\pm k+\ii Y)$, both sums are
$O(1/Y)$.  Hence $\Phi_{\mathrm{att}}\circ\Phi_{\mathrm{rep}}^{-1}(z)-z\to0$ as
$\operatorname{Im}z\to+\infty$, so its constant Fourier mode vanishes:
$B_0=0$.  Numerically $|B_0|<10^{-59}$.

By (a)--(d) and Theorem~\ref{thm:C}, $\tilde T\to h^{-1}-a$ uniformly on
$\ell$, where $h^{-1}$ denotes the inverse horn map (not the period $h$).
So $\tilde t_n\to B_n$ and $\tilde t_0\to-a$, and
$\tau_n=\tilde t_n\eee^{-2\pi\ii n\tilde t_0}$ gives:

\begin{corollary}\label{cor:tauconv}
For every $n\ge1$, $\tau_n(b)\to B_n\eee^{2\pi\ii na}$ as $b\to\eta^-$.
Since $B_1\ne0$, also $t_n/t_1^n\to B_n/B_1^n$ for every $n\ge1$.
\end{corollary}

\noindent
A direct pointwise check of Theorem~\ref{thm:C} on two gate points
(\texttt{confluence\_check.py}) gives errors
$0.086,\,0.021,\,0.0050,\,0.00079$ (attracting side) at
$\lambda=0.8,0.9,0.95,0.98$, decreasing like $\varepsilon^2$.

\section{The first-order term: a parabolic series}\label{sec:first}

Theorem~\ref{mthm:A} gives the limit of $\tau_n$.  Here we prove that $\tau_n$ is
differentiable at $b=\eta$ from the side $b<\eta$.  The derivative is given by an
absolutely convergent series computed entirely at the parabolic germ.  Put
\[
 s=1-\eee\log b=1-\mu,\qquad f_s(u)=\eee^{-s+(1-s)u}-1 ,
\]
so $s\in(0,s_0]$ corresponds to $b<\eta$ near $\eta$ and $f_0(u)=\eee^u-1$.  The
map $(s,u)\mapsto f_s(u)$ is entire.

\subsection{Why termwise differentiation fails, and the remedy}
With the model time $\Psi_b$ of Section~\ref{sec:proofC}, the one-step error is
only $O(|u-u_1||u-u_2|)$ (Lemma~\ref{lem:F}).  Along a parabolic orbit
$u_k\approx-2/k$ one has $\partial_su_k\asymp k$.  Hence both $\partial_sF(u_k)$ and
$F'(u_k)\partial_su_k$ are of order one, and the differentiated series diverges.
The remedy is the idea behind the prepared form of
Marde\v{s}i\'c--Roussarie--Rousseau~\cite{MRR2004}.  We add to the model time a
polynomial correction, holomorphic in $s$, that makes the one-step error vanish
to order $N$ at both fixed points simultaneously.

By the Weierstrass preparation theorem,
\begin{equation}\label{eq:weier}
 f_s(u)-u=q_s(u)\,k_s(u),\qquad q_s(u)=u^2-\sigma_1(s)u+\sigma_2(s),
\end{equation}
on a polydisc $|s|<s_0$, $|u|<r$.  Here $\sigma_1,\sigma_2$ are holomorphic,
$k_s$ is holomorphic and zero-free, and $k_0(0)=\tfrac12$.  From
$f_s(0)=-s+O(s^2)$ and $f_s'(0)-1=-2s+O(s^2)$ we get $\sigma_2=-2s+O(s^2)$,
$\sigma_1=O(s)$, and the discriminant
\[
 D(s)=\sigma_1^2-4\sigma_2=8s+O(s^2)
\]
is biholomorphic near $0$.  For real $s>0$ the roots $u_1<0<u_2$ are the fixed
points of Section~\ref{sec:proofC}.  Put $t=\sqrt{D(s)}$.  Then
$u_{1,2}=\tfrac12(\sigma_1\mp t)$ are holomorphic functions of $t$, interchanged by
$t\mapsto-t$.  Since $f_s'(u_1)=1+(u_1-u_2)k_s(u_1)$, the residues are
$A_1=a_s(u_1,u_2)$ and $A_2=a_s(u_2,u_1)$, where
\[
 a_s(x,y)=\frac1{\log\bigl(1+(x-y)k_s(x)\bigr)}
 =\frac1{(x-y)k_s(x)}+\frac12+(x-y)\,h_s(x,y)
\]
with $h_s$ holomorphic.

\begin{lemma}[Uniform model]\label{lem:model}
Fix $N\ge1$.  After shrinking $s_0$ and $r$, there are holomorphic functions
$e_1,\dots,e_{2N-2}$ on $\{|s|<s_0\}$ with the following property.  Put
$P_s(u)=\sum_ie_i(s)u^i$ and, for $s\ne0$,
\[
 \Psi_s(u)=A_1\log(u-u_1)+A_2\log(u-u_2)+P_s(u).
\]
Then:
\begin{enumerate}
\item $\nu(s)=A_1+A_2$ extends holomorphically to $|s|<s_0$.  For every
  $\delta\in(0,r)$ there is $s_\delta>0$ such that
  $\Psi_s(u)-\nu(s)\log u$ extends to a holomorphic function of
  $(s,u)\in\{|s|<s_\delta\}\times\{\delta<|u|<r\}$.  Here, for $|u|>\delta$,
  the branches are fixed by $\log(u-u_j):=\log u+\operatorname{Log}(1-u_j/u)$
  with one common branch of $\log u$.  Individually $A_j\sim\pm1/\sqrt{2s}$, so
  any other convention would introduce $2\pi\ii A_j$-jumps that are not
  holomorphic in $s$.
\item The one-step error
  \[
   F_s(u)=A_1\log\frac{f_s(u)-u_1}{u-u_1}+A_2\log\frac{f_s(u)-u_2}{u-u_2}
   +P_s(f_s(u))-P_s(u)-1
  \]
  (principal logarithms) extends to a holomorphic function on
  $\{|s|<s_0\}\times\{|u|<r\}$, and $F_s=q_s^N\,\Gamma_s$ with $\Gamma$
  holomorphic and $|\Gamma|\le M$ there.
\end{enumerate}
\end{lemma}

\begin{proof}
\emph{Removability.}  Let $\Theta(s,x,y,u)=a_s(x,y)\ell(y)+a_s(y,x)\ell(x)$,
where $\ell$ is holomorphic in $(s,x,y,u)$.  By the expansion of $a_s$, its
singular part is
\[
 \frac{\ell(y)/k_s(x)-\ell(x)/k_s(y)}{x-y}.
\]
The numerator vanishes on $x=y$, so $\Theta$ extends holomorphically across the
diagonal.  Evaluate $\Theta$ at $(x,y)=(u_1,u_2)$.  The result is a
holomorphic function of $t$, even in $t$, hence a holomorphic function of
$t^2=D(s)$, hence of $s$.  The same holds for the other pairing
$\Theta'(s,x,y,u)=a_s(x,y)\ell(x)+a_s(y,x)\ell(y)$, whose singular part
$[\ell(x)/k_s(x)-\ell(y)/k_s(y)]/(x-y)$ is again removable and symmetric.
These arguments apply to the following choices of $\ell$:
\begin{itemize}
\item $\ell\equiv1$ in $\Theta'$, which gives $\nu$;
\item $\ell(x)=\log(1-x/u)$ in $\Theta'$, which gives $\Psi_s-\nu\log u-P_s$ for
  $|u|>\delta$ once $|u_j|<\delta/2$;
\item $\ell(x)=\log(1+(u-x)k_s(u))$ in $\Theta$, which gives
  $F^{\log}_s:=A_1\log\frac{f-u_1}{u-u_1}+A_2\log\frac{f-u_2}{u-u_2}-1$,
  because $(f_s(u)-u_1)/(u-u_1)=1+(u-u_2)k_s(u)$.
\end{itemize}
This proves~(1) and the holomorphy of $F^{\log}$.

\emph{Divisibility.}  For $s\ne0$, $F^{\log}_s(u_j)=A_j\log f_s'(u_j)-1=0$
(as in Lemma~\ref{lem:F}).  Weierstrass division by $q_s$ has a remainder
$r_0(s)+r_1(s)u$ that is holomorphic in $s$ and vanishes at two distinct points
for $s\ne0$.  So the remainder is $0$, and $F^{\log}=q_s\varphi$ with $\varphi$
holomorphic.  Likewise $D_i:=(f_s(u))^i-u^i=q_s d_i$ (powers, not iterates).  At $s=0$,
$D_i=u^i\bigl((1+\tfrac u2+\cdots)^i-1\bigr)$, so
$d_i=\tfrac i2u^{i-1}+O(u^i)$.

\emph{The linear system.}  Divide $\varphi+\sum e_id_i$ by the Weierstrass
polynomial $q_s^{N-1}$ of degree $2N-2$.  The $2N-2$ coefficients of the
remainder depend holomorphically on $s$ and affinely on $e$, with matrix
$M(s)$.  At $s=0$ the remainder is the Taylor polynomial modulo $u^{2N-2}$, so
$M(0)$ is lower triangular with diagonal entries $i/2\ne0$.  Hence
$e(s)=-M(s)^{-1}\,\mathrm{rem}(\varphi)$ is holomorphic near $0$.  Then
$F_s=F^{\log}_s+\sum e_iD_i=q_s(\varphi+\sum e_id_i)$ is divisible by
$q_s\cdot q_s^{N-1}$.  The bound on $\Gamma$ follows after shrinking the
polydisc.
\end{proof}

Only the one-step error enters the estimates below, and it is single-valued.
The logarithms in $\Psi_s$ are continued along orbits exactly as in the proof of
Theorem~\ref{thm:C}.  Since $P_s$ is single-valued, replacing $\Psi_b$ by
$\Psi_s$ changes no branch.  The path of item~(c) in Section~\ref{sec:proofC}
runs outside $D_\delta(0)$ and does not pass between $u_1$ and $u_2$.  So
along it the convention of Lemma~\ref{lem:model}(1) agrees with the continuation
used there.

\subsection{Orbits for complex parameters}
Let $P_R=\{\operatorname{Re}(-2/u)>R\}$ and put $\zeta=-2/u$.

\begin{lemma}[Finite orbits]\label{lem:finite}
Let $K\subset P_R$ be compact, with $R$ large, and let
$Z=\max_K|2/u|$.  There is $c>0$ with the following property.  For every
integer $J\ge1$, every $\sigma\in\mathbb C$ with $|\sigma|\le c/J^2$, every
$u\in K$ and every $0\le k\le J$, the iterate $u_k=f_\sigma^k(u)$ is defined
and $\zeta_k=-2/u_k$ satisfies
\[
 \operatorname{Re}\zeta_k\ge\operatorname{Re}\zeta_0+k/2,\qquad
 |\zeta_k|\le|\zeta_0|+2k .
\]
In particular $|u_k|\le2/(R+k/2)$.  For $\sigma=0$ this holds for all $k$.
\end{lemma}

\begin{proof}
For $|u|\le r$ we have $-2/f_0(u)=-2/u+1+\epsilon(u)$ with $|\epsilon(u)|\le C_0|u|$,
and $|f_0(u)|\ge|u|/2$.  Write $f_\sigma=f_0+\sigma g$ with $|g|\le G$.
Suppose the claim holds up to $k<J$.  Then $|u_k|\le2/R\le r$ and
$w:=\sigma g(u_k)/f_0(u_k)$ satisfies
\[
 |w|\le G|\sigma||\zeta_k|\le Gc(Z+2)\le\tfrac12 .
\]
Hence
\[
 \zeta_{k+1}=\frac{-2}{f_0(u_k)}(1+w)^{-1}=\zeta_k+1+e_k,\qquad
 |e_k|\le\frac{2C_0}R+4G|\sigma||\zeta_k|^2\le\frac{2C_0}R+4Gc(Z+2)^2 .
\]
This is at most $\tfrac12$ for $R\ge8C_0$ and $c$ small.  So
$\operatorname{Re}\zeta_{k+1}\ge\operatorname{Re}\zeta_k+\tfrac12$ and
$|\zeta_{k+1}|\le|\zeta_k|+\tfrac32$.
\end{proof}

\begin{lemma}\label{lem:derbounds}
Let $u\in K$, $u_k=f_0^k(u)$ and $\dot u_k=\partial_\sigma f_\sigma^k(u)|_{\sigma=0}$.
Then
\[
 |\dot u_k|\le C(k+1),\qquad |\partial_\sigma F_\sigma(v)|_{\sigma=0}|\le M'|v|^{2N-2},\qquad
 |F_0'(v)|\le M'|v|^{2N-1}\quad(|v|\le r/2).
\]
\end{lemma}

\begin{proof}
For the first bound, apply Cauchy's estimate on $|\sigma|=c/(k+1)^2$ together
with Lemma~\ref{lem:finite} (with $J=k+1$).  This gives
$|\dot u_k|\le\frac{2}{R+k/2}\cdot\frac{(k+1)^2}{c}$.  For the other two, write
$F_\sigma=q_\sigma^N\Gamma_\sigma$ with $q_0(v)=v^2$ and $\partial_\sigma q_\sigma$
bounded, and bound $\partial_\sigma\Gamma$ and $\Gamma_0'$ by Cauchy's estimates.
\end{proof}

\subsection{One-sided differentiability}
For $u$ in a compact $K\subset P_R$ and real $s\in[0,s_0]$ put
\[
 \mathcal A_s(u)=\Psi_s(u)+\sum_{k\ge0}F_s(f_s^ku),
\]
with a fixed branch of $\log u$ on a simply connected neighbourhood $U\subset P_R$
of $K$.  For $s>0$ this is the attracting Koenigs Abel coordinate up to an
additive constant, by the argument of Theorem~\ref{thm:C}(1).  For $s=0$,
$\mathcal A_0=\Phi_{\mathrm{att}}+\mathrm{const}$, because
$\Psi_0-\alpha=\mathrm{const}+O(u)$.

\begin{proposition}\label{prop:diff}
Let $N\ge3$.  Uniformly for $u\in K$, as $s\to0^+$,
\[
 \mathcal A_s(u)=\mathcal A_0(u)+s\,\dot{\mathcal A}(u)+O(s^{10/9}),
\]
where
\begin{equation}\label{eq:Adot}
 \dot{\mathcal A}(u)=\partial_s\Psi_s(u)\big|_{s=0}
 +\sum_{k\ge0}\Bigl[\partial_\sigma F_\sigma(u_k)\big|_{\sigma=0}
 +F_0'(u_k)\,\dot u_k\Bigr],\qquad u_k=f_0^k(u).
\end{equation}
The series converges absolutely and uniformly on $K$, with terms
$O(k^{2-2N})$.  Hence $\dot{\mathcal A}$ is holomorphic on the interior of
$K$.
\end{proposition}

\begin{proof}
The $\Psi$-part is holomorphic in $s$ by Lemma~\ref{lem:model}(1), so it
contributes $O(s^2)$.  Fix $J\ge1$ and split the sum at $J$.

\emph{Head.}  $H_J(\sigma)=\sum_{k<J}F_\sigma(f_\sigma^ku)$ is holomorphic on
$|\sigma|<c/J^2$.  There $|u_j(\sigma)|\le C\sqrt{|\sigma|}\le C\sqrt c/J$.
Since $k<J$ gives $1/J\le(R+1)/(R+k/2)$, Lemma~\ref{lem:finite} yields
$|q_\sigma(u_k)|\le C_3(R+k/2)^{-2}$.  Therefore
\[
 |H_J(\sigma)|\le MC_3^N\sum_k(R+k/2)^{-2N}=:M_1,
\]
uniformly in $J$.  Cauchy's estimate gives $|H_J''|\le8M_1J^4/c^2$ on
$[0,c/(2J^2)]$, so
\[
 |H_J(s)-H_J(0)-sH_J'(0)|\le4M_1c^{-2}s^2J^4 .
\]

\emph{Tails.}  We may assume $K\subset P_{R+1}$.  Since $Z_s\to-2/u$ uniformly
on $K$ (Lemma~\ref{lem:res}), $K\subset\Omega_{R,s}$ for all small $s>0$.  If $r$
was shrunk in Lemma~\ref{lem:model}, increase $R$ so that
$\Omega_{R,s}\subset D_{r/2}$, and reach a general $K$ by $f^m$ as below.  For
real $s\in(0,s_0]$ the proof of Lemma~\ref{lem:petal} gives
$|q_s(f_s^ku)|\le C_1(R+\tfrac34k)^{-2}$.  For $s=0$, Lemma~\ref{lem:finite}
gives the analogous bound $|u_k|^2\le4(R+k/2)^{-2}$.  Hence
$\bigl|\sum_{k\ge J}F_s(f_s^ku)\bigr|\le C_4J^{1-2N}$ for all $s\in[0,s_0]$.
By Lemma~\ref{lem:derbounds}, the terms of~\eqref{eq:Adot} are
$O((R+k/2)^{2-2N})$.  So $H_J'(0)$ differs from the series by at most
$C_5J^{3-2N}$.

\emph{Balance.}  For $s\le c/(2J^2)$,
\[
 \bigl|\mathcal A_s-\mathcal A_0-s\dot{\mathcal A}\bigr|
 \le Cs^2+4M_1c^{-2}s^2J^4+2C_4J^{1-2N}+C_5sJ^{3-2N}.
\]
Take $J=\lceil s^{-2/9}\rceil$, which is admissible for small $s$.  With
$N\ge3$ every term is $O(s^{10/9})$.
\end{proof}

\begin{remark}
Taking $J=s^{-2/(2N+3)}$ instead gives the error $O(s^{2-8/(2N+3)})$, so
$O(s^{2-\epsilon})$ for any $\epsilon>0$ with $N$ large.  The proof needs $N\ge3$
only because of the crude bound on the head ($N=2$ would give $s^{6/7}$).  The
series converges for $N\ge2$.

The derivative does not depend on $N$.  The condition $F_0=O(u^{2N})$ forces
$\Psi_0$ to agree with the formal Abel function $\alpha$ up to a constant and
$O(u^{2N-1})$, so the $e_i(0)$ are the coefficients of $\alpha$ and do not depend
on $N$.  Hence $\mathcal A^{(N)}_s-\mathcal A^{(N')}_s$, which equals the
difference of the two polynomials evaluated at $u_1(s)=O(\sqrt s)$, is a
constant plus $O(s^{3/2})$.  In any case additive constants cancel in
$\dot{\mathcal A}(v)-\dot{\mathcal A}(u^*)$.  Numerically, $N=3,4,6,8$ give the
same $\kappa^{(1)}$ to $16$--$20$ digits.
\end{remark}

\subsection{The derivative of $\tau_n$}
Compact subsets of the basin are handled by
$\mathcal A_s=\mathcal A_s\circ f_s^m-m$, where $m$ is chosen so that
$f_0^m(K)\Subset P_R$.  The iterate $f_s^m$ is jointly holomorphic, so the
expansion of Proposition~\ref{prop:diff} persists, with
$\dot{\mathcal A}=\dot{\mathcal A}\circ f_0^m+\mathcal A_0'\circ f_0^m\cdot\partial_sf_s^m$.

For the repelling side, apply everything to the reflected family
$G_s(v)=-f_s^{-1}(-v)=-(\log(1-v)+s)/(1-s)$, exactly as in
Theorem~\ref{thm:C}(2).  This family is again of the form~\eqref{eq:weier},
with fixed points $-u_2,-u_1$.  Let $\mathcal S_s(u)=-\mathcal A^G_s(-u)$.  By
Remark~\ref{rem:invariance} and~\eqref{eq:taun}, $\tau_n$ does not depend on the
additive normalisation of the repelling coordinate.  We use
$\mathcal S^\sharp_s=\mathcal S_s-\mathcal S_s(p')+\Phi_{\mathrm{rep}}(p')$.
This differs from the real normalisation by the constant called
$s(b)$ in Theorem~\ref{thm:C}(2), which we denote $s_*(b)$ here to avoid a clash
with the parameter $s$, and it is right-differentiable with derivative
$\dot{\mathcal S}^\sharp$.  On the attracting side,
$\Rr_b^{-1}(\eee(1+u))=\mathcal A_s(u)-\mathcal A_s(u^*)$, with the branch fixed
in item~(c) of Section~\ref{sec:proofC}.

\begin{theorem}\label{thm:D}
Let $\ell=[0,1]+\ii Y_*$ be the gate line of Section~\ref{sec:proofC}, and put
$v(z)=\Phi_{\mathrm{rep}}^{-1}(z)$ and $\tilde T_0=h^{-1}-a$.  Define
\[
 \delta T(z)=\dot{\mathcal A}(v(z))-\dot{\mathcal A}(u^*)
 -\tilde T_0'(z)\,\dot{\mathcal S}^\sharp(v(z)),\qquad
 \delta t_m=\int_\ell\delta T(z)\,\eee^{-2\pi\ii mz}\,dz .
\]
Then for every $n\ge1$, as $s\to0^+$,
\[
 \tau_n=B_n\eee^{2\pi\ii na}+s\,\dot\tau_n+O(s^{10/9}),\qquad
 \dot\tau_n=\eee^{2\pi\ii na}\bigl(\delta t_n-2\pi\ii n\,B_n\,\delta t_0\bigr).
\]
Moreover $p=-\log\lambda_1\log\lambda_2=2s+\tfrac{10}9s^2+O(s^3)$ is
holomorphic in $s$.  Consequently, if $B_n\ne0$,
\[
 \log\frac{\tau_n}{B_n\eee^{2\pi\ii na}}=\kappa^{(n)}p+O(p^{10/9}),\qquad
 \kappa^{(n)}=\frac{\delta t_n}{2B_n}-\pi\ii n\,\delta t_0 .
\]
\end{theorem}

\begin{proof}
By Proposition~\ref{prop:diff} and Cauchy's estimates, both $\mathcal A_s$ and
$\mathcal S^\sharp_s$ have expansions $X_0+s\dot X+O(s^{10/9})$ in the
$C^1$ topology on a compact neighbourhood of $v(\ell)$.  By Rouch\'e, the inverse
$v_s=(\mathcal S^\sharp_s)^{-1}$ exists near $v(\ell)$ and satisfies
\[
 v_s=v-s\,\dot{\mathcal S}^\sharp(v)/\Phi_{\mathrm{rep}}'(v)+O(s^{10/9}).
\]
Hence $\tilde T_s:=\Rr_b^{-1}\circ\eee(1+v_s)$ satisfies
$\tilde T_s=\tilde T_0+s\,\delta T+O(s^{10/9})$ uniformly on $\ell$, where we used
$\mathcal A_0'(v)/\Phi_{\mathrm{rep}}'(v)=\tilde T_0'$.  By items (a)--(c) of
Section~\ref{sec:proofC}, the Fourier coefficients $\tilde t_m(s)$ of
$\tilde T_s-\mathrm{id}$ on $\ell$ satisfy
$\tau_n=\tilde t_n\eee^{-2\pi\ii n\tilde t_0}$.  Differentiate, using
$\tilde t_n(0)=B_n$ and $\tilde t_0(0)=-a$.  The expansion of $p$ follows from
the removability argument of Lemma~\ref{lem:model} applied to
$\log\lambda_1\log\lambda_2$, together with a direct Taylor expansion.  Finally,
$\log(\tau_n/\tau_n(0))=s\,\dot\tau_n/\tau_n(0)+O(s^{10/9})$ and $s=p/2+O(p^2)$.
\end{proof}

\subsection{Higher orders}\label{sec:higher}
The argument of Proposition~\ref{prop:diff} extends to every order, at the
price of a larger $N$.

\begin{proposition}\label{prop:higher}
Let $m\ge1$ and $N\ge m^2+m+1$.  For $u\in K$ and $j\ge0$ put
\begin{equation}\label{eq:Aj}
 \mathcal A^{(j)}(u)=\partial_s^j\Psi_s(u)\big|_{s=0}
 +\sum_{k\ge0}\partial_\sigma^j\bigl[F_\sigma(f_\sigma^ku)\bigr]_{\sigma=0}.
\end{equation}
For $j\le m$ the series converges absolutely and uniformly on $K$, and its
terms are $O(k^{2j-2N})$.  Uniformly for $u\in K$, as $s\to0^+$,
\[
 \mathcal A_s(u)=\sum_{j=0}^m\frac{s^j}{j!}\mathcal A^{(j)}(u)+O(s^{m+\gamma}),\qquad
 \gamma=1-\frac{2(m+1)^2}{2N+2m+1}>0 .
\]
\end{proposition}

\begin{proof}
\emph{Term bounds.}  Fix $k\ge0$ and apply Lemma~\ref{lem:finite} with $J=k+1$.
For $|\sigma|\le c/(k+1)^2$ the head estimate in the proof of
Proposition~\ref{prop:diff} gives $|q_\sigma(f_\sigma^ku)|\le C_3(R+k/2)^{-2}$,
hence $|F_\sigma(f_\sigma^ku)|\le MC_3^N(R+k/2)^{-2N}$.  By Cauchy's estimate on
this disc,
\[
 \bigl|\partial_\sigma^j[F_\sigma(f_\sigma^ku)]_{\sigma=0}\bigr|
 \le j!\,MC_3^N\,c^{-j}(k+1)^{2j}(R+k/2)^{-2N}=O(k^{2j-2N}).
\]
This is summable for $j\le m<N-\tfrac12$.  For $m=1$ this recovers the bounds of
Lemma~\ref{lem:derbounds}, without computing $\dot u_k$.

\emph{Head.}  $H_J(\sigma)=\sum_{k<J}F_\sigma(f_\sigma^ku)$ is holomorphic and
bounded by $M_1$ on $|\sigma|<\rho:=c/J^2$.  For real $0\le s\le\rho/2$, the Taylor
remainder of order $m$ satisfies
\[
 \Bigl|H_J(s)-\sum_{j\le m}\frac{s^j}{j!}H_J^{(j)}(0)\Bigr|
 \le M_1\sum_{j>m}(s/\rho)^j\le2M_1c^{-m-1}\,s^{m+1}J^{2m+2}.
\]

\emph{Tails.}  As before,
$\bigl|\sum_{k\ge J}F_s(f_s^ku)\bigr|\le C_4J^{1-2N}$ for $s\in[0,s_0]$.  By
the term bounds, the tails of the $j$-th derivative series from $J$ on are
$O(J^{2j+1-2N})$.  Since $sJ^2\le c/2$, the tails contribute
$\sum_{j\le m}s^jJ^{2j+1-2N}=O(J^{1-2N})$.

\emph{Balance.}  Write $\mathcal T_J^{(j)}$ for the tail from $k=J$ on of the
$j$-th derivative series.  The identity
\[
 \mathcal A_s-\sum_{j\le m}\frac{s^j}{j!}\mathcal A^{(j)}
 =\Bigl[\Psi_s-\sum_{j\le m}\frac{s^j}{j!}\partial_s^j\Psi\Bigr]
 +\Bigl[H_J(s)-\sum_{j\le m}\frac{s^j}{j!}H_J^{(j)}(0)\Bigr]
 +\sum_{k\ge J}F_s(f_s^ku)-\sum_{j\le m}\frac{s^j}{j!}\mathcal T_J^{(j)}
\]
holds exactly.  The $\Psi$-part is holomorphic in $s$.  Altogether the error
is $O(s^{m+1}J^{2m+2}+J^{1-2N})$.  Take $J=\lceil s^{-\beta}\rceil$ with
$\beta=(m+1)/(2N+2m+1)<\tfrac12$, which is admissible.  Both terms are then
$O(s^{m+1-2(m+1)\beta})=O(s^{m+\gamma})$, and $\gamma>0$ exactly when
$N\ge m^2+m+1$.
\end{proof}

For $m=1$ this is Proposition~\ref{prop:diff}; the error there is
$O(s^{10/9})$ for $N=3$.

\begin{theorem}\label{thm:E}
For every $n\ge1$ with $B_n\ne0$, the function $\log(\tau_n/\tau_n(0))$ has an
asymptotic expansion to all orders in $p$ as $b\to\eta^-$:
\[
 \log\frac{\tau_n}{B_n\eee^{2\pi\ii na}}\sim\sum_{j\ge1}\kappa^{(n)}_jp^j,\qquad
 \kappa^{(n)}_1=\kappa^{(n)} .
\]
For each order $m$, the remainder after the $p^m$ term is $O(p^{m+1})$.
Each coefficient $\kappa^{(n)}_j$ is a polynomial in the following data:
\begin{itemize}
\item $B_n^{-1}$ and $a$;
\item the Taylor coefficients of $s(p)$;
\item finitely many Fourier coefficients on $\ell$ of explicit polynomial
  expressions in $\mathcal A^{(i)}\circ v$, $\mathcal S^{\sharp(i)}\circ v$
  ($i\le j$), their $u$-derivatives, and $1/\Phi_{\mathrm{rep}}'\circ v$.
\end{itemize}
Here $\mathcal A^{(i)}$ and $\mathcal S^{\sharp(i)}$ are given by the convergent
series~\eqref{eq:Aj} over parabolic orbits of $f_0$ and $G_0$.  The coefficients
do not depend on $N$.
\end{theorem}

\begin{proof}
Fix $m$ and $N\ge m^2+m+1$.  By Proposition~\ref{prop:higher} and the reduction
$\mathcal A_s=\mathcal A_s\circ f_s^{M}-M$ (with $f_s^{M}$ jointly holomorphic),
$\mathcal A_s$ and $\mathcal S^\sharp_s$ have expansions
$\sum_{j\le m}s^jX_j/j!+O(s^{m+\gamma})$.  These hold uniformly on a compact
neighbourhood of $v(\ell)\cup\{u^*\}$ and are holomorphic in $u$, so by Cauchy's
estimates they also hold in the $C^r$ topology on a smaller neighbourhood, for
every $r$.

\emph{Inverse.}  Formal inversion of
$\sum_js^j\mathcal S^\sharp_j/j!$ gives a polynomial approximant
$\tilde v_s=v+\sum_{1\le j\le m}s^jV_j$ with
$\mathcal S^\sharp_s(\tilde v_s)=z+O(s^{m+\gamma})$.  Since
$\mathcal S^\sharp_s\to\Phi_{\mathrm{rep}}$ in $C^1$ and $\Phi_{\mathrm{rep}}'\ne0$
near $v(\ell)$, the maps $\mathcal S^\sharp_s$ are uniformly bi-Lipschitz there,
and the true inverse satisfies $v_s=\tilde v_s+O(s^{m+\gamma})$.

\emph{Composition.}  Taylor-expand $\mathcal A_s$ about $v$, using the bounded
$u$-derivatives.  This gives $\tilde T_s=\sum_{j\le m}s^jT_j/j!+O(s^{m+\gamma})$
uniformly on $\ell$.  Hence the Fourier coefficients $\tilde t_i(s)$ and
$\tau_n=\tilde t_n\eee^{-2\pi\ii n\tilde t_0}$ have expansions to order $m$.
Since $\tau_n(0)\ne0$, so does $\log(\tau_n/\tau_n(0))$.

\emph{Change of parameter.}  Finally, $p$ is holomorphic in $s$ with
$p'(0)=2$, so $s$ is a holomorphic function of $p$.

\emph{Independence of $N$.}  An asymptotic expansion of a fixed function is
unique, so the coefficients do not depend on the choice of $N$.

\emph{Remainder.}  Applying the argument at order $m+1$ and absorbing the
$p^{m+1}$ term gives the remainder $O(p^{m+1})$.
\end{proof}

\begin{remark}
In particular, Theorem~\ref{thm:E} with $m=2$ improves the remainder in
Theorem~\ref{thm:D} from $O(p^{10/9})$ to $O(p^2)$.
Theorem~\ref{thm:E} gives explicit asymptotic coefficients in the normalisation
of $\tau_n$, but says nothing about summability.  The
$\tfrac12$-summability of~\cite{ChristopherRousseau}, in their normalisation,
suggests Gevrey-$2$ growth $|\kappa^{(n)}_j|\lesssim C^j(j!)^2$ in $p\propto\epsilon$
(Gevrey-$1$ in $\sqrt p$).  This matches the scale $\Lambda=\exp(-C/\sqrt p)$ of the
beyond-all-orders terms.
The constants in our proof grow faster, because $N$ must grow like $m^2$.
\end{remark}

\begin{remark}[Relation to the literature]
Christopher and Rousseau show that the unfolded modulus, in a particular
normalisation, is $\tfrac12$-summable in the canonical parameter
$\epsilon=(A_2-A_1)^{-2}$, with the Glutsyuk direction as the singular
direction~\cite[Cor.~4.9]{ChristopherRousseau}.  One checks the exact identity
$p=4\epsilon/(1-\nu^2\epsilon)$.  Their result gives the existence of an
asymptotic expansion.  Marde\v{s}i\'c--Roussarie--Rousseau proved continuity
at $\epsilon=0$ and conjectured $1$-summability in
$\sqrt\epsilon$~\cite[\S11]{MRR2004}; this is the result
of~\cite{ChristopherRousseau}, and Rib\'on proved multisummability of Fatou
coordinates and \'Ecalle--Voronin invariants for unfoldings of any finite
codimension~\cite{Ribon2010}.  These results give existence of the expansion
but not its coefficients.  Theorems~\ref{thm:D} and~\ref{thm:E} give the
coefficients, in the normalisation of $\tau_n$, as explicit convergent series
over parabolic orbits, with a direct proof that does not use summability
theory.  This parallels the expression of the \'Ecalle--Voronin invariants
themselves by convergent numerical series~\cite{DudkoSauzin2014}.
On the Lavaurs side, normalised perturbed Fatou coordinates and horn maps
depend holomorphically on the perturbation~\cite{Shishikura2000}
(see also~\cite[Thms.~1.3, 2.1]{InouShishikura}), so there too an expansion
exists, but we have not found its first-order coefficients in the literature.
The closest analogue is the splitting of separatrices of area-preserving maps
near a saddle-centre bifurcation: Gelfreich and Br\"annstr\"om prove that the
homoclinic invariant is $\eee^{-2\pi^2/\log\lambda}$ times an asymptotic series
$\sum_k a_k\delta^{2k}$, with $a_0$ the Stokes constant of the parabolic
limit~\cite{GelfreichBrannstrom,GelfreichSauzin2001}, and Gelfreich and Sim\'o
extract the higher $a_k$ numerically to high precision~\cite{GelfreichSimo2008}.
That setting concerns one Fourier mode of a two-dimensional map, and its
coefficients for $k\ge1$ are defined through auxiliary problems rather than
convergent orbit sums.  Within a parabolic family (no splitting of the fixed
point), orbit-sum expansions of the horn map in a parameter appear
in~\cite[Prop.~9.1]{BuffEcalleEpstein2013}.  We know of no
closed form.  An integer-relation search on $\kappa^{(1)}$ against
$\pi,\pi^2,\log2,\gamma$ returns only spurious relations.
\end{remark}

\subsection{Numerics}
The series~\eqref{eq:Adot} is evaluated with $N=3,\dots,8$ on $24$ points of
the line $\operatorname{Im}z=1$ or $1.5$, with orbits of length $300$--$3000$.
The $s$-derivative of each (finite, holomorphic) truncation is taken by a
central difference at $s=\pm10^{-20}$ to $\pm10^{-30}$ in $60$--$90$ digit
arithmetic (\texttt{mrr\_derivative.py}).  The results are stable to all
printed digits:
\[
\begin{aligned}
 \kappa^{(1)}&=-0.010199006345897402441-0.013250956182506712250\,\ii,\\
 \kappa^{(2)}&=-0.05105655620132933-0.06231578053849410\,\ii,\\
 \kappa^{(3)}&=-0.140044348129-0.169224463926\,\ii .
\end{aligned}
\]
These agree with the fits~\eqref{eq:kappa1} from the side $b<\eta$ to all of their
eleven digits.  At $s=0$ the same computation reproduces
$\tau_n(0)=B_n\eee^{2\pi\ii na}$ to the precision of the stored $B_n$.  The second
coefficients $\kappa^{(n)}_2$ of Theorem~\ref{thm:E} are computed by a
five-point stencil ($s=\pm10^{-12},\pm2\cdot10^{-12}$, $90$ digits) applied to
the same finite truncations (\texttt{mrr\_second.py}), with $N=7$ ($\operatorname{Im}z=1$,
orbits of length $2000$) and $N=10$ ($\operatorname{Im}z=1.5$, length $3000$).  The
series converge for $N\ge3$; the condition $N\ge m^2+m+1=7$ is needed only for
the proof.  The two computations agree to $18$, $16$ and $12$ digits
respectively:
\[
\begin{aligned}
 \kappa^{(1)}_2&=-7.88992074991511334\cdot10^{-5}+3.6554727325799459\cdot10^{-4}\,\ii,\\
 \kappa^{(2)}_2&=-1.895481905818775\cdot10^{-4}+1.106655243691997\cdot10^{-4}\,\ii,\\
 \kappa^{(3)}_2&=4.07137984831\cdot10^{-4}-8.49506527544\cdot10^{-3}\,\ii .
\end{aligned}
\]
The fits from the side $b<\eta$ (\texttt{deficit\_modes.py}) give
$-7.88991\cdot10^{-5}+3.65547\cdot10^{-4}\ii$,
$-1.89548\cdot10^{-4}+1.10666\cdot10^{-4}\ii$ and
$4.07137\cdot10^{-4}-8.49507\cdot10^{-3}\ii$, in agreement to within one unit in
the sixth digit.

\section{The sewing solution: decay estimate}\label{sec:W2}

Throughout, $b<\eta$ is real and close to $\eta$, $\varepsilon=|\log\lambda|$,
and $S$ has the real normalisation of Section~\ref{sec:transition}, so
$T(z)-z=\sum t_n\eee^{2\pi\ii nz}$ with $\operatorname{Im}t_0=h/2$ and
$t_{-n}=\overline{t_n}$ (Lemma~\ref{lem:realT}).

\begin{lemma}[The band]\label{lem:band}
Let $Z_b(u)=A_1\log\frac{u-u_1}{u-u_2}$ with the branch
$\operatorname{Im}Z_b\in(0,h)$ on $\mathbb C\setminus((-\infty,u_1]\cup[u_2,\infty))$.
This branch has $\operatorname{Im}Z_b=h/2$ on $(u_1,u_2)$ and
$\operatorname{Im}Z_b<h/2$ on $\operatorname{Im}u>0$.  For $Y>0$ let
\[
 \mathcal B_b(Y)=\{u:\ Y<\operatorname{Im}Z_b(u)<h-Y\},
\]
a lens bounded by two circular arcs through $u_1,u_2$.  There is an absolute
$Y_B$ such that for $Y\ge Y_B$ and all $b$ near $\eta$:
\begin{enumerate}
\item[(i)] every $u\in\mathcal B_b(Y)$ has its forward orbit in $D_{r'}$ tending
  to $u_1$ and its principal backward orbit in $D_{r'}$ tending to $u_2$, with
  $\sum_{k\in\mathbb Z}|F_b(f_b^ku)|\le C/Y+C\varepsilon$;
\item[(ii)] the real-normalised times $\mathcal R_b=\Psi_b+\sum_{k\ge0}F_b\circ f_b^k+c_R$
  and $\mathcal S_b=\Psi_b-\sum_{k\ge1}F_b\circ f_b^{-k}+c_S$ are holomorphic and
  injective on $\mathcal B_b(Y)$.  They are branches of $\Rr_b^{-1}$ and
  $S_b^{-1}$, with $\mathcal S_b$ real on $(u_1,u_2)$ and $\mathcal R_b$ as in
  item (c) above, and
  \[
   \bigl|\operatorname{Im}\mathcal S_b(u)-(\operatorname{Im}Z_b(u)-h/2)\bigr|
   \le\tfrac\pi3+C/Y+C\varepsilon ;
  \]
\item[(iii)] $S_b$ maps the strip $\Sigma=\{|\operatorname{Im}w'|<h/2-Y-2\}$
  into $\mathcal B_b(Y)$ with $\mathcal S_b\circ S_b=\mathrm{id}$.  Consequently
  $T=\mathcal R_b\circ S_b$ is holomorphic on $\Sigma$, and
  \[
   T(w')-w'=(c_R-c_S)+\sum_{k\in\mathbb Z}F_b\bigl(f_b^k(S_b(w'))\bigr).
  \]
\end{enumerate}
\end{lemma}

\begin{proof}
(i) Put $q=\eee^{Z/A_1}=(u-u_1)/(u-u_2)$ and $\zeta=Z/A_1=-\varepsilon Z$, so
$\operatorname{Im}\zeta\in(-2\pi+\varepsilon Y,-\varepsilon Y)$ on the band.  On
$-2\pi\le\operatorname{Im}\zeta\le0$ one has
$|\eee^\zeta-1|\ge c\min(1,\delta(\zeta))\max(1,|\eee^\zeta|)$ with
$\delta(\zeta)=\min(|\zeta|,|\zeta+2\pi\ii|)$ and an absolute $c>0$; the only
zeros are $0$ and $-2\pi\ii$.  On the band
$\delta(\zeta)=\varepsilon\min(|Z|,|Z-\ii h|)>\varepsilon Y$.  Since
$u-u_2=(u_2-u_1)/(q-1)$ and $|u_2-u_1|/\varepsilon\to2$
(Lemma~\ref{lem:res}), this gives
\begin{equation}\label{eq:bandsize}
 |u-u_2|\le C\bigl(\varepsilon+Y^{-1}\bigr),\qquad
 |u-u_1||u-u_2|=\frac{|q|\,|u_2-u_1|^2}{|q-1|^2}\le
 \begin{cases}C\,\min(|Z|,|Z-\ii h|)^{-2},&\delta(\zeta)\le1,\\
 C\varepsilon^2\eee^{-\varepsilon|\operatorname{Re}Z|},&\delta(\zeta)>1,\end{cases}
\end{equation}
using $|q|/\max(1,|q|)^2=\eee^{-|\operatorname{Re}\zeta|}$.  So
$\mathcal B_b(Y)\subset D_{C/Y+|u_1|+|u_2|}\subset D_{r'/2}$ for $Y$ large.  By the
proof of Lemma~\ref{lem:petal}, $Z\circ f-Z=1+O(|u|+\varepsilon)$, so
$\operatorname{Re}Z$ increases by at least $\tfrac34$ per step while the orbit is
in $D_{r'/2}$.

The imaginary part is controlled through
$\Psi_b=Z_b+(A_1+A_2)\log(u-u_2)$, with step error
$|F_b|\le M|u-u_1||u-u_2|$ (Lemma~\ref{lem:F}).  Summing~\eqref{eq:bandsize}
along $\operatorname{Re}Z_k\ge\operatorname{Re}Z_0+\tfrac34k$ bounds
$\sum_k|F_b(u_k)|$ by $C'/Y+C'\varepsilon$.  Points with
$\delta(\zeta_k)\le1$ contribute $\sum_k(x_k^2+Y^2)^{-1}\le C/Y$, and the others
$\varepsilon^2\sum\eee^{-\varepsilon|x_k|}\le C\varepsilon$.  Now
$\operatorname{Im}Z-\operatorname{Im}\Psi=-(A_1+A_2)\arg(u-u_2)$ with
$\arg(u-u_2)\in(0,2\pi)$ on the lens.  So $\operatorname{Im}Z$ drifts by at most
$C'/Y+C'\varepsilon+\tfrac{2\pi}3+o(1)$ along the orbit, which is $<Y/2$ for
$Y\ge Y_B\ge8$.  By induction the orbit stays in
$\mathcal B_b(Y/2)\subset D_{r'/2}$, and $\operatorname{Re}Z_k\to+\infty$,
i.e.\ $u_k\to u_1$.  The continued values of $Z$ along the orbit have
$\operatorname{Im}\in(Y/2,h-Y/2)\subset(0,h)$, so they agree with the fixed
slit-plane branch.  The backward orbit is the forward orbit of the reflected
inverse family $G_b(v)=-f_b^{-1}(-v)$, to which the same argument applies
(its prepared form is checked in the proof of Theorem~\ref{thm:C}(2)).  Its
period is $h_2=h+2\pi/3+o(1)$, so $\mathcal B_b(Y)$ is contained in the
corresponding band of $G_b$ at height $Y(1+O(\varepsilon))-1$.  This is harmless
for $Y\ge Y_B$.

(ii) As in the proof of Theorem~\ref{thm:C}(1), $\Psi_b(f^nu)-n\to A_1\log\sigma_b(u)+\mathrm{const}$
along forward orbits.  So $\mathcal R_b$ is a branch of the attracting Abel
coordinate, and $\mathcal S_b$ likewise of the repelling one.  Choose the
constants $c_R$ (by item (c) above) and $c_S$ so that $\mathcal S_b$ is real on
the segment.  Both $F_b$ and $f_b$ are real, so the sums are real there.  On
the segment $\operatorname{Im}Z=h/2$ and $\arg(u-u_2)=\pi$, which fixes
$\operatorname{Im}c_S$.  The displayed estimate then follows from
$\operatorname{Im}\Psi=\operatorname{Im}Z+(A_1+A_2)\arg(u-u_2)$,
$A_1+A_2\to\tfrac13$ and (i).  Injectivity: if $\mathcal S_b(u)=\mathcal S_b(u')$
then the principal backward orbits satisfy
$\tau_b(u_{-n})=\tau_b(u'_{-n})$ near $u_2$, where $\tau_b$ is injective, so
$u_{-n}=u'_{-n}$.  Moreover $f_b$ is injective on $D_{r'}$ (its non-injectivity
has period $2\pi\ii/\log b$ in $w$, far larger than the disc), so $u=u'$.  The
same argument with $\sigma_b$ and forward orbits proves injectivity of
$\mathcal R_b$.

(iii) Fix $x\in\mathbb R$.  $S_b(x)\in(u_1,u_2)\subset\mathcal B_b(Y)$ and
$\mathcal S_b(S_b(x))=x$, since both sides are the real-normalised repelling
time.  Move $y'$ from $0$ towards $\pm(h/2-Y-2)$.  As long as $S_b(x+\ii y')$
stays in the lens, continuation gives $\mathcal S_b(S_b(x+\ii y'))=x+\ii y'$.
It cannot leave through the corners $u_1,u_2$, where
$|\operatorname{Re}\mathcal S_b|\to\infty$.  It cannot leave through the arcs
either: there $\operatorname{Im}Z\in\{Y,h-Y\}$, so by (ii)
$|\operatorname{Im}\mathcal S_b|\ge h/2-Y-\pi/3-o(1)>|y'|$.  Hence
$S_b(\Sigma)\subset\mathcal B_b(Y)$, $T=\mathcal R_b\circ S_b$ on $\Sigma$, and
$T(w')-w'=\mathcal R_b(u)-\mathcal S_b(u)$ with $u=S_b(w')$ is the displayed
bi-infinite sum.
\end{proof}

\begin{lemma}[Uniform mode bound]\label{lem:modes}
Put $Y_0:=Y_B+3$.  There is $M$, which can be taken $\le1$, and $b_1<\eta$ such
that for $b\in(b_1,\eta)$ and all $n\ne0$,
\[
 |t_n|\le M\eee^{-2\pi|n|d},\qquad d:=\tfrac h2-Y_0-1 .
\]
\end{lemma}

\begin{proof}
Take $Y=Y_B+1$ in Lemma~\ref{lem:band}(iii), so that the estimate (ii) is
available on the closure of the lens.  $\Sigma$ has half-width
$h/2-Y_B-3=d+1$, so the lines $\operatorname{Im}w'=\pm d$ lie in $\Sigma$.  On them, by (iii) and (i),
$T-\mathrm{id}-t_0=\sum_kF_b(f_b^kS_b)-\text{mean}$ has modulus at most $2(C/Y_B+C\varepsilon)$.  Let $M$ be the supremum of this over
$b\in(b_1,\eta)$.  Cauchy's estimate on these lines gives the bound
for $n>0$ and $n<0$.  This does not use Theorem~\ref{thm:C}.
\end{proof}

For $H>0$ let $\mathfrak A_H$ be the space of $1$-periodic
$g=\sum g_n\eee^{2\pi\ii nz}$ with
$\|g\|_H=\sum|g_n|\eee^{2\pi|n|H}<\infty$.  It is a Banach algebra
(the weight is submultiplicative), its elements are holomorphic and bounded by
$\|g\|_H$ on $|\operatorname{Im}z|\le H$, and the projections
$\Pi_{<0},\Pi_{\ge0}$ onto negative and non-negative modes have norm $1$.  Also
$\|\eee^{g}\|_H\le\eee^{\|g\|_H}$ and
$\|\eee^{g_1}-\eee^{g_2}\|_H\le\|g_1-g_2\|_H\,\eee^{\max\|g_j\|_H}$.

\begin{theorem}[W2]\label{thm:W2}
Let $M,d$ be as in Lemma~\ref{lem:modes}, $\rho=\tfrac14$, and choose $D\ge1$
with $4\pi M\eee^{-2\pi D}\le\rho(1-\eee^{-2\pi D})^2$.  Put $H=d-D-\rho$ and
$\tau=T-\mathrm{id}-t_0$.  Assume $H>1$, which holds for $b$ near $\eta$ since
$d\sim\pi/\varepsilon$, and that $b$ is so close to $\eta$ that
$C_1\Lambda\eee^{4\pi(Y_0+D+1)}<\tfrac12$, where $C_1=4M+2\rho$.  Then:
\begin{enumerate}
\item The map $\mathcal T(Q)=-\Pi_{<0}[\tau\circ(\mathrm{id}+Q)]$ is a
  contraction, with constant $\le\tfrac12$, of the ball
  $\{Q\in\Pi_{<0}\mathfrak A_H:\|Q\|_H\le\rho\}$ into itself.  Its fixed point
  $Q^*$ and $P':=\Pi_{\ge0}[\tau\circ(\mathrm{id}+Q^*)]$ satisfy
  \[
   T\circ(\mathrm{id}+Q^*)=\mathrm{id}+t_0+P'\qquad\text{on }|\operatorname{Im}z|<H.
  \]
\item $\varphi_-=\mathrm{id}+Q^*$ is univalent on $\{\operatorname{Im}z<H-1\}$
  and $\varphi_+=\mathrm{id}+t_0+P'$ on $\{\operatorname{Im}z>-(H-1)\}$.  The
  function equal to $\Rr\circ(\varphi_+-\ii h)$ above and $S\circ\varphi_-$ below
  is a solution of~\eqref{eq:abel}; it is the sewing solution $\Kk^W_b$ of
  Proposition~\ref{prop:weld}.
\item The equation $\varphi_+(z_0)=\ii h$ has a unique solution with
  $|z_0-(\ii h-t_0)|<1$, and $|z_0-(\ii h-t_0)|\le C_1\Lambda\eee^{4\pi(Y_0+D+1)}$.
  After the translation $z\mapsto z+z_0$ one has $\Kk^W_b(0)=1$ and
  $P=\Rr^{-1}\circ\Kk^W_b-\mathrm{id}$ with $P(0)=0$.
\item For every $n\ge1$ the Fourier coefficients of this $P$ satisfy
  \[
   \Bigl|\frac{c_n}{\Lambda^n}-\tau_n(b)\Bigr|\le C_n\Lambda,
   \qquad C_n=C\,M(M+1)\,\eee^{2\pi(n+2)(Y_0+D+2)},
  \]
  with an absolute constant $C$.
\end{enumerate}
\end{theorem}

\begin{proof}
(1) Write $\tau\circ(\mathrm{id}+Q)=\sum_{m\ne0}t_m\eee^{2\pi\ii mz}\eee^{2\pi\ii mQ}$.
Then
\[
 \|\tau\circ(\mathrm{id}+Q)\|_H\le\sum_{m\ne0}|t_m|\eee^{2\pi|m|(H+\rho)}
 \le2M\sum_{m\ge1}\eee^{-2\pi mD}\le\rho
\]
by Lemma~\ref{lem:modes} and the choice of $D$.  Similarly
\[
 \|\tau\circ(\mathrm{id}+Q_1)-\tau\circ(\mathrm{id}+Q_2)\|_H\le
 \sum_{m\ne0}|t_m|\eee^{2\pi|m|(H+\rho)}2\pi|m|\,\|Q_1-Q_2\|_H
 \le\frac{4\pi M\eee^{-2\pi D}}{(1-\eee^{-2\pi D})^2}\|Q_1-Q_2\|_H ,
\]
which is at most $\tfrac12\|Q_1-Q_2\|_H$.  Since $|\operatorname{Im}(z+Q)|\le H+\rho<d$
on the strip, the Fourier series of $T$ converges at $z+Q(z)$ and equals $T$
there (Lemma~\ref{lem:band}).  The identity follows by splitting
$\tau\circ(\mathrm{id}+Q^*)=P'+\Pi_{<0}[\ldots]=P'-Q^*$.

(2) The series of $Q^*$ has only negative modes and converges on
$\operatorname{Im}z\le H$, so it converges on the whole lower half-plane
$\operatorname{Im}z<H$.  There $|{Q^*}'|\le2\pi\eee^{-2\pi}\rho\,/(1-\eee^{-2\pi})^2<1$ on
$\operatorname{Im}z\le H-1$, so $\operatorname{Re}\varphi_-'>0$ and $\varphi_-$
is univalent (Noshiro--Warschawski).  The same holds for $\varphi_+$ above.  On
the strip, $\Rr\circ(\varphi_+-\ii h)=\Rr\circ(T\circ\varphi_--\ii h)
=\Rr\circ T\circ\varphi_-=S\circ\varphi_-$, because $\ii h$ is a period of
$\Rr$ and $\Rr\circ T=S$.  The functional equation holds because $\varphi_\pm$
commute with $z\mapsto z+1$.  The resulting function realises the sewing of
Proposition~\ref{prop:weld}: the images of $\varphi_\pm$ are the two
half-cylinders and the seam in the $z$-plane is $\Gamma=\varphi_-^{-1}(\mathbb R)$, a curve in
$|\operatorname{Im}z|\le\rho$.  Since $\varphi_\pm-\mathrm{id}$ is bounded, the
argument principle shows that the restrictions of $\varphi_\pm$ to the two sides of
$\Gamma$ map onto the two half-cylinders.  By uniqueness of the uniformisation
it is $\Kk^W_b$ up to translation.  (The identity $\Rr\circ T=S$ is used on the
band, where orbits stay in $D_{r'}$ and $\mathcal R_b$ inverts $\Rr$;
Lemma~\ref{lem:band}.)

(3) For $\operatorname{Im}z\ge h/2-1$ we have
$|P'(z)-p'_0|\le\sum_{n\ge1}|p'_n|\eee^{-2\pi n(h/2-1)}$ and
$|p'_n|\le\rho\,\eee^{-2\pi nH}$.  Moreover
$p'_0=\sum_{m\ge1}t_m[\eee^{2\pi\ii mQ^*}]_{-m}$, and
$|[\eee^{2\pi\ii mQ^*}]_{-m}|\le\eee^{2\pi m\rho}\eee^{-2\pi mH}$, so
$|p'_0|\le2M\eee^{-2\pi(d+H-\rho)}\le2M\Lambda\eee^{2\pi(2Y_0+D+2+2\rho)}$, and
$\sum_{n\ge1}|p'_n|\eee^{-2\pi n(h/2-1)}\le2\rho\Lambda\eee^{2\pi(Y_0+D+2+\rho)}$.
Both are bounded by
\[
 \tfrac12 C_1\Lambda\eee^{4\pi(Y_0+D+1)}.
\]
Since
$\operatorname{Im}(\ii h-t_0)=h/2$ and $|{P'}'|\le\rho$ there, Rouch\'e's theorem
on the disc of radius $1$ gives the unique root and the bound, using the
smallness hypothesis.  Then
$\Kk^W_b(z_0)=\Rr(\varphi_+(z_0)-\ii h)=\Rr(0)=1$, and after translation
$P(z)=\varphi_+(z+z_0)-\ii h-z$, so $P(0)=0$.

(4) For $n\ge1$ the $n$-th mode of $P$ is $c_n=p'_n\eee^{2\pi\ii nz_0}$.
Because $Q^*$ has no constant or positive modes,
$[\eee^{2\pi\ii mQ^*}]_0=1$ and $[\eee^{2\pi\ii mQ^*}]_k=0$ for $k>0$.  Hence
$p'_n=t_n+\sum_{m>n}t_m[\eee^{2\pi\ii mQ^*}]_{n-m}$, and
\[
 |p'_n-t_n|\le\sum_{m>n}M\eee^{-2\pi md}\eee^{2\pi m\rho}\eee^{-2\pi(m-n)H}
 \le2M\eee^{-2\pi n(d-\rho)}\eee^{-2\pi(d+H-\rho)} .
\]
With $z_0=\ii h-t_0+\zeta$, $|\zeta|\le C_1\Lambda\eee^{4\pi(Y_0+D+1)}$ by (3), we have
$\eee^{2\pi\ii nz_0}=\Lambda^n\eee^{-2\pi\ii nt_0}\eee^{2\pi\ii n\zeta}$ and
$|t_n\eee^{-2\pi\ii nt_0}|=|\tau_n|$, so
$|t_n|=|\tau_n|\eee^{-\pi nh}$.  Dividing,
\[
 \frac{c_n}{\Lambda^n}-\tau_n=\tau_n\bigl(\eee^{2\pi\ii n\zeta}-1\bigr)
 +(p'_n-t_n)\eee^{-2\pi\ii nt_0}\eee^{2\pi\ii n\zeta}.
\]
In the first term $|\tau_n|=|t_n|\eee^{\pi nh}\le M\eee^{2\pi n(Y_0+1)}$ and
$|\eee^{2\pi\ii n\zeta}-1|\le4\pi n|\zeta|$ when $2\pi n|\zeta|\le1$; in
the second, $|\eee^{2\pi\ii n\zeta}|\le\eee$ in that case.  If instead
$2\pi n|\zeta|>1$, then
$\Lambda\ge(2\pi nC_1)^{-1}\eee^{-4\pi(Y_0+D+1)}$, and the trivial bound
$|c_n|\Lambda^{-n}+|\tau_n|\le C'M\eee^{2\pi n(Y_0+D+2)}$ is already
$\le C_n\Lambda$ (this uses $D\ge1$).  In the second term,
$|(p'_n-t_n)\eee^{-2\pi\ii nt_0}|\le2M\eee^{2\pi n(Y_0+1+\rho)}\Lambda
\eee^{2\pi(2Y_0+2+D+2\rho)}$, by $d=h/2-Y_0-1$ and $H=d-D-\rho$.  Adding the two
gives the stated bound, since
$n(Y_0+1+\rho)+2Y_0+2D+2.5\le(n+2)(Y_0+D+2)$.
\end{proof}

\begin{corollary}\label{cor:Kw}
For the sewing solution, $\hat c_n(\Kk^W_b)\to B_n\eee^{2\pi\ii na}$ for every
$n\ge1$ as $b\to\eta^-$.  Moreover,
$\kappa_n(\Kk^W_b)\to B_n/B_1^n$ and
$\lim\arg\hat c_1=\arg B_1+2\pi a$.
\end{corollary}

\begin{proof}
Theorem~\ref{thm:W2}(4) and Corollary~\ref{cor:tauconv}: $C_n$ does not depend
on $b$ and $\Lambda\to0$.  Proposition~\ref{prop:horn-nonzero} makes $c_1$
nonzero for $b$ near $\eta$, justifying division and the limiting phase.
\end{proof}

\begin{corollary}[Local parabolic limit of the sewn solution]
\label{cor:sewn-cusp-local}
Put $u^*=1/\eee-1$ and $a=\Phi_{\mathrm{att}}(u^*)$.  On some fixed disk
about $z=0$, the function
\[
 F_\eta^{\mathrm{att}}(z)
 =\eee\bigl(1+\Phi_{\mathrm{att}}^{-1}(z+a)\bigr)
\]
is holomorphic, satisfies $F_\eta^{\mathrm{att}}(0)=1$, and has
$(F_\eta^{\mathrm{att}})'(0)\ne0$.  As $b\to\eta^-$ through real bases,
$\Kk_b^W\to F_\eta^{\mathrm{att}}$ locally uniformly on that disk,
with all $z$-derivatives.  Consequently there are $r>0$ and $b_2<\eta$
such that, for $b_2<b<\eta$, the only preimage of $1$ under $\Kk_b^W$
in $B(0,r)$ is $0$; both $\Kk_b^W$ and $1/\Kk_b^W$ are uniformly
bounded there.
\end{corollary}

\begin{proof}
Theorem~\ref{thm:C}(1) gives locally uniform convergence of
$A_b(u):=\Rr_b^{-1}(\eee(1+u))$ to
$A_\eta(u):=\Phi_{\mathrm{att}}(u)-a$ near $u^*$, with $A_b(u^*)=0$.
The derivative $A_\eta'(u^*)$ is nonzero: the series in the proof of
Theorem~\ref{thm:C}, Lemmas~\ref{lem:F} and~\ref{lem:petal}, and Cauchy's
estimate give $\Phi_{\mathrm{att}}'(u)=2/u^2+O(1/u)$ on a narrower
attracting petal.  The Abel identity and $f_\eta'(u)\ne0$ propagate
nonvanishing back along the orbit of $u^*$.  Uniform
convergence of these local inverses therefore gives
$\Rr_b\to F_\eta^{\mathrm{att}}$ on a fixed disk about $0$.

In Theorem~\ref{thm:W2}, write
$P_b(z)=\zeta_b+P_b'(z+z_0(b))$, where
$\zeta_b=z_0(b)+t_0(b)-\ii h(b)\to0$.  The proof of that theorem gives
$p'_0(b)\to0$ and $|p'_n(b)|\le\rho\eee^{-2\pi nH(b)}$ for $n\ge1$,
while $H(b)\to\infty$ and $\operatorname{Im}z_0(b)=h(b)/2+o(1)$.
Hence $P_b\to0$ uniformly on every fixed bounded $z$-disk on which
the upper representation is defined.  That representation contains a
fixed disk about $0$ for all sufficiently large $b$, and
$\Kk_b^W(z)=\Rr_b(z+P_b(z))$ there.  This proves the asserted convergence;
Cauchy's formula gives convergence of derivatives.  Choose a small
closed disk on which $F_\eta^{\mathrm{att}}$ is nonzero and
$F_\eta^{\mathrm{att}}-1$ has only its simple zero at $0$.
Rouch\'e's theorem and $\Kk_b^W(0)=1$ give the final assertions.
\end{proof}

\noindent
The construction is directly computable (\texttt{weld\_solve.py}).  At
$b=1.10$ the iteration contracts by $\sim10^{-4}$ per step and gives
$\hat c_1=0.0864272029\,\eee^{1.5510921403\ii}$, against
$\tau_1=0.0864272022\,\eee^{1.5510921408\ii}$.  The relative difference is
$8.8\cdot10^{-9}=0.38\Lambda$, and at $b=1.30$ it is $3.9\cdot10^{-19}=0.40\Lambda$.
At these bases $h$ is too small for the hypotheses of Theorem~\ref{thm:W2}, so
the agreement is not an instance of the theorem.  It does match the first-order
expansion of the correction,
$c_1/(\Lambda\tau_1)-1\approx\Lambda\bigl[-4\pi^2|\tau_1|^2-2\pi\ii\tau_1
-4\pi\ii\tau_2\bar\tau_1/\tau_1\bigr]$, whose modulus is $0.3834\Lambda$ and
$0.3984\Lambda$ respectively .

\section{Dependence on the base}\label{sec:W1}

This section proves that $\Kk^W_b$ depends holomorphically on the base near
$(b_1,\eta)$, with quantitative bounds; the identification with the boundary
value $\Kk_b$ of~\eqref{eq:K} is completed in Section~\ref{sec:cusp}.

\begin{theorem}[W1, local]\label{thm:W1}
Let $b_0\in(b_1,\eta)$.  There is a disc $\Delta\ni b_0$ in the base plane such
that:
\begin{enumerate}
\item For $b'\in\Delta$, with $L(b'),L_2(b')$ the continued fixed points,
  $\Rr_{b'},S_{b'}$ the regular superfunctions, $\varpi(b')=2\pi\ii/\log\lambda(b')$
  the period of $\Rr_{b'}$, and $\hat T_{b'}=\Rr_{b'}^{-1}\circ S_{b'}+\varpi/2$,
  the contraction of Theorem~\ref{thm:W2} works with $\rho$ and $D$
  unchanged, $2M$ in place of $M$, and $H-\tfrac12$ in place of $H$.  The solution is
  $\Kk^W_{b'}=\Rr_{b'}\circ(\varphi_++\varpi)$ above, where
  $\varphi_+=T_{b'}\circ\varphi_-$, and $S_{b'}\circ\varphi_-$ below.  It is
  normalised by the root of $\varphi_+(z_0)=-\varpi$ in the disc of radius $1$
  about $-\varpi-t_0$.  (For real $b'$, $\varpi=-\ii h$.)
  Its fixed point, the resulting solution $\Kk^W_{b'}$, and the normalising root
  $z_0(b')$ are holomorphic in $b'$.
\item For $b'\in\Delta$ with $\operatorname{Im}b'>0$ one has $\arg\lambda>0$,
  $\arg\lambda_2<0$, and $\Kk^W_{b'}(z)\to L(b')$ as $\operatorname{Im}z\to+\infty$,
  $\Kk^W_{b'}(z)\to L_2(b')$ as $\operatorname{Im}z\to-\infty$, uniformly for
  $\operatorname{Re}z$ in compacts.  In the upper and lower half-planes (of
  the translated coordinate) it has the representations
  $\Rr_{b'}\circ(\mathrm{id}+\theta_{\mathrm{up}})$ and
  $S_{b'}\circ(\mathrm{id}+\theta_{\mathrm{dn}})$.  These are the defining
  representations of the complex-base construction in \texttt{\_cbuild.py}.
\item (Local uniqueness.)  If a solution $F$ with $F(0)=1$ admits such
  representations and, after the translation that puts its seam
  representation in the form $T\circ(\mathrm{id}+Q)=\mathrm{id}+t_0+P'$, has
  $\|Q\|_H\le\rho$ and its normalising root in that disc, then
  $F=\Kk^W_{b'}$.  (Without the root condition the statement is false: the
  translates of $\Kk^W_{b'}$ by the other roots $c_j$, $j\ge1$, have the same $Q$.)
\end{enumerate}
\end{theorem}

\begin{proof}
(1) By Lemma~\ref{lem:band}(iii), $T$ is holomorphic on the strip of half-width $d+1$,
so there is slack $1$ in $d$; we use half of it.  The Koenigs maps at $L(b')$, $L_2(b')$ and the principal logarithms
defining $\Rr^{-1}$ depend holomorphically on $b'$.  To choose one
parameter disc for the unbounded strip, first work on the compact fundamental
rectangle $0\le\operatorname{Re}z\le1$,
$|\operatorname{Im}z|\le d(b_0)+\tfrac12$.  Its image under $S_{b_0}$
is compactly contained in the common domain of the two chosen Abel charts;
their local parameter continuations therefore give a single disc $\Delta$
on this rectangle.  The identity $T_{b'}(z+1)=T_{b'}(z)+1$ extends the
result to the whole strip.  Thus $\hat T_{b'}-\mathrm{id}$ is jointly
holomorphic on $\Delta\times\{|\operatorname{Im}z|<d(b_0)+\tfrac12\}$.
Shrinking $\Delta$ again,
its modes satisfy $|t_n(b')|\le2M\eee^{-2\pi|n|(d(b_0)-1/2)}$ (Cauchy), and the
estimates of Theorem~\ref{thm:W2}(1) hold with $2M$ in place of $M$:
the original choice of $D$ gives a Lipschitz constant at most
$2\rho=\tfrac12$ and a ball-image norm at most
$\rho(1-\eee^{-2\pi D})/\pi<\rho$.  The iterates $\mathcal T^k(0)$ are holomorphic in $b'$
and converge uniformly, so $Q^*$ is holomorphic in $b'$.  So are $P'$ and,
by the implicit function theorem (Rouch\'e), $z_0$.

(2) From $\lambda\eee^{-\lambda}=\log b'$,
$\partial\lambda/\partial b=\eee^{\lambda}/(b(1-\lambda))$ is positive for
$\lambda<1$ and negative for $\lambda_2>1$.  Hence $\operatorname{Im}b'>0$
gives $\operatorname{Im}\lambda>0>\operatorname{Im}\lambda_2$ for $\Delta$
small.  Above, $\Kk^W=\Rr(z+c+o(1))$, and
$\Rr(w)=\sigma^{-1}(c'\lambda^w)\to L$ when
$\operatorname{Re}(w\log\lambda)\to-\infty$.  This is the case as
$\operatorname{Im}w\to+\infty$, since
$\operatorname{Re}(w\log\lambda)=\operatorname{Re}w\log|\lambda|-\operatorname{Im}w\arg\lambda$.
Below, the same argument with $\arg\lambda_2<0$ applies.  The representations
are those of Theorem~\ref{thm:W2}(2), transported by $z\mapsto z+z_0$.

(3) The fixed point of a contraction is unique in its ball.  $P'$ and $z_0$
are then determined, the latter by $F(0)=1$ together with the root
condition.
\end{proof}

Theorem~\ref{thm:W1} gives quantitative control of the sewn family on complex
discs about real bases near $\eta$.  The base continuation from $b>\eta$ passes the
Shell--Thron boundary, where $|\Lambda|=1$ and the contraction does not apply;
Section~\ref{sec:cusp} reaches the sewn family by a different route, a tilted
initial region that is insensitive to $|\Lambda|$.

\section{Characterisation of the sewn solution}\label{sec:char}

The sewn solution is characterised by its two regular representations and
the position of its seam.  This is the uniqueness statement in
Theorem~\ref{mthm:B}; it is also what fixes the root in the identification of
Section~\ref{sec:cusp}.

\begin{definition}\label{def:class}
For real $b\in(1,\eta)$ and $Y>0$ let $\mathcal C_b(Y)$ be the set of
solutions $F$ of~\eqref{eq:abel} with $F(0)=1$ such that
(``solution'' means holomorphic off a closed singular set and satisfying the
equation where defined)
\begin{enumerate}
\item[(a)] $F=\Rr\circ(\mathrm{id}+P)$ on $\{\operatorname{Im}z>-2Y\}$, with $P$
  holomorphic, $1$-periodic and bounded there, and $P(0)=0$ (the branch
  $\Rr^{-1}(1)=0$);
\item[(b)] $F=S\circ(\mathrm{id}+Q)$ on $\{\operatorname{Im}z<-Y\}$, with $Q$
  holomorphic, $1$-periodic and bounded there, and
  $|\operatorname{Im}Q-h/2|\le1$ on the strip $O=\{-2Y<\operatorname{Im}z<-Y\}$.
\end{enumerate}
\end{definition}

Condition (b) says that the real line of repelling time, i.e.\ the segment
$(L,L_2)$, sits half a period of attracting time below the real axis of $F$.
Heuristically, some such condition is necessary.  For the sewing solution the
roots of $\varphi^W_+(c)\in\ii h\mathbb Z$ near the axis are
$c_j=c_0+\ii jh+O(\Lambda)$.  The translates $\Kk^W(\cdot+c_j)$ keep an
$\Rr$-representation above and an $S$-representation on some lower half-plane,
and give $|c_1|\asymp\Lambda^{1+j}$.  Condition (b) excludes $j\ne0$, as the
proof shows.

\begin{theorem}[Characterisation]\label{thm:char}
There are $b_1<\eta$ and $Y_B$ such that for $b\in(b_1,\eta)$ and
$Y_0+D+3<Y<(h-Y_0-D-3)/2$ (with $Y_0,D$ as in Theorem~\ref{thm:W2}),
\[
 \mathcal C_b(Y)=\{\Kk^W_b\},
\]
with $\Kk^W_b$ normalised as in Theorem~\ref{thm:W2}(3).  In particular
$\hat c_n\to B_n\eee^{2\pi\ii na}$ for every $n$ holds for the unique element
of $\mathcal C_b(Y)$.
\end{theorem}

\begin{proof}
\emph{$\Kk^W_b\in\mathcal C_b(Y)$.}  In the normalised coordinate,
$\varphi_-(z)=z+z_0+Q^*(z+z_0)$ is defined for
$\operatorname{Im}(z+z_0)<H$, i.e.\ for $\operatorname{Im}z<-(Y_0+1+D+\rho)+O(\Lambda)$,
which contains $\{\operatorname{Im}z<-Y\}$.  There $|Q-z_0|\le\rho$, so
$|\operatorname{Im}Q-h/2|\le\rho+O(\Lambda)<1$.  The $\Rr$-representation holds
for $\operatorname{Im}(z+z_0)>-(H-1)$, i.e.\ above
$-(h-Y_0-D-2-\rho)$, which contains $\{\operatorname{Im}z>-2Y\}$.  $P$ and $Q$
are bounded there.

\emph{Uniqueness.}  Let $F\in\mathcal C_b(Y)$, $\varphi_+=\mathrm{id}+P$ and
$\varphi_-=\mathrm{id}+Q$.  Work with the strip
$\Sigma_q=\{|\operatorname{Im}w'|<d-D\}$ of $S$-time.  By Lemma~\ref{lem:modes} and
the choice of $D$, on $\overline{\Sigma_q}$
\[
 |T-\mathrm{id}-t_0|\le\frac{2M\eee^{-2\pi D}}{1-\eee^{-2\pi D}}\le\frac{\rho}{2\pi},\qquad
 |T'-1|\le\rho<1 ,
\]
so $T$ is univalent on the convex set $\Sigma_q$ (Noshiro--Warschawski),
continuous up to the boundary, and within $\rho/2\pi$ of a translation.

(i) \emph{The overlap.}  Let $O=\{-2Y<\operatorname{Im}z<-Y\}$.  For $z\in O$,
$\operatorname{Im}\varphi_-(z)\in(h/2-2Y-1,\,h/2-Y+1)\subset(-(d-D),d-D)$,
because $2Y<h-Y_0-D-2$ and $Y>Y_0+D+2$.  Since
$\sigma\circ\Rr(w)=\sigma(1)\lambda^w$ wherever $\Rr$ is defined,
$\Rr\circ\varphi_+=S\circ\varphi_-=\Rr\circ T\circ\varphi_-$ gives
$\lambda^{\varphi_+}=\lambda^{T\circ\varphi_-}$.  Hence
$\varphi_+=T\circ\varphi_--\ii h+k\ii h$ on $O$ for a constant $k\in\mathbb Z$.
No injectivity is needed for this step.

(ii) \emph{The lifted sewn surface.}  Let $V_-=\{\operatorname{Im}w'<d-D\}$
(repelling time) and let $V_+$ be the region above the curve
$T(\{\operatorname{Im}w'=-(d-D)\})-\ii h$ (attracting time).  Let $\tilde X$ be
$V_-\sqcup V_+$ modulo $w=T(w')-\ii h$ for $w'\in\Sigma_q$.  Since $T$ is a
homeomorphism of $\overline{\Sigma_q}$ onto its image, the identification
graph is closed and $\tilde X$ is a Hausdorff Riemann surface.  Translation by
$1$ acts freely on it, and $X=\tilde X/\mathbb Z\cong\mathbb C^*$, both ends
being punctured discs.  The function $\tilde K$, equal to $\Rr$ on $V_+$ and
$S$ on $V_-$, is well defined, because $\Rr(T(w')-\ii h)=S(w')$ on the band.
The sewing solution is $\tilde K\circ\tilde\Phi^W$, where $\tilde\Phi^W$ is the
lift of its uniformisation; sewing along the real line or along any other
essential curve in $\Sigma_q$ gives the same $\tilde X$.

Define $\tilde\Phi:\mathbb C\to\tilde X$ by $\varphi_+-k\ii h$ on
$\{\operatorname{Im}z>-2Y\}$ and $\varphi_-$ on $\{\operatorname{Im}z<-Y\}$.
These agree on $O$ by (i).  They land in $V_\mp$ as follows.
$\operatorname{Im}Q$ is bounded and harmonic on a half-cylinder, hence extends
harmonically over the puncture.  So on $\{\operatorname{Im}z\le-Y-\delta\}$ it is
bounded by its values on the line $\operatorname{Im}z=-Y-\delta\subset O$.
This gives $\operatorname{Im}\varphi_-\le-Y+h/2+1<d-D$.  Likewise, on a line
in $O$ near $\operatorname{Im}z=-2Y$, (i) gives
$\operatorname{Im}(\varphi_+-k\ii h)\ge\operatorname{Im}z-1-\rho/2\pi$.  The minimum
principle then gives $\operatorname{Im}(\varphi_+-k\ii h)\ge-2Y-1-\rho/2\pi$
above that line.  This exceeds the height
$-h/2-d+D+\rho/2\pi$ of the lower boundary of $V_+$, because
$2Y<h-Y_0-D-3$.

(iii) \emph{$\tilde\Phi$ is a translate of $\tilde\Phi^W$.}  $\tilde\Phi$
commutes with $z\mapsto z+1$ and descends to a holomorphic
$\Phi:\mathbb C/\mathbb Z\to X\cong\mathbb C^*$.  $\Phi$ is proper, since
$\varphi_\pm-\mathrm{id}$ are bounded, and has local degree one at each end,
since $P\to p_0$ and $Q\to q_0$.  A proper holomorphic map of $\mathbb C^*$ that
fixes the ends with local degree one is linear.  Hence
$\tilde\Phi=\tilde\Phi^W(\cdot+c)$ for some $c\in\mathbb C$, and
$F=\tilde K\circ\tilde\Phi=\Kk^W_b(\cdot+c)$ in the coordinate of
Theorem~\ref{thm:W2}(1)--(2).  Univalence of $\varphi_\pm$ is a conclusion, not
an assumption.

(iv) \emph{The root.}  Take a line $\ell\subset O$ with
$\operatorname{Im}z=-Y''$ and $Y''$ close to $2Y$.  There
$\operatorname{Im}\varphi_-\le-2Y+h/2+1+o(1)<H-\rho$, so $\varphi_-(\ell)$ lies
in $\varphi_-^W(\{\operatorname{Im}<H\})$.  The inclusion
$\varphi_-^W(\{\operatorname{Im}<H\})\supset\{\operatorname{Im}<H-\rho\}$ holds because
$|Q^*|\le\rho$: for $w$ in the smaller set, $\varphi^W_--w$ has winding
number $1$ around $0$ on the boundary of a large fundamental rectangle.  By injectivity of
$\tilde\Phi^W$, the points $z+c$ with $z\in\ell$ lie in
$\{\operatorname{Im}<H\}$.  There $Q(z)=c+Q^*(z+c)$, so
$|\operatorname{Im}c-h/2|\le1+\rho$.  On the upper chart,
$\varphi_+(z)-k\ii h=\varphi^W_+(z+c)-\ii h$, and $P(0)=0$ gives
$\varphi^W_+(c)=\ii h(1-k)$.  Since $|P'|\le\rho$ where $\varphi^W_+$ is defined,
$|c-(\ii h(1-k)-t_0)|\le\rho$.  Comparing imaginary parts with the bound on
$\operatorname{Im}c$ gives $|k|h\le1+2\rho$, so $k=0$.  Then $c$ lies in the
Rouch\'e disc of Theorem~\ref{thm:W2}(3), so $c=z_0$ and $F=\Kk^W_b$.  The last
sentence of the theorem is Corollary~\ref{cor:Kw}.
\end{proof}

% Section: general real saddle-node unfoldings.  \input from main.tex after Section~\ref{sec:char}.
% Source of the audit: docs/general-statements-zh.md; numerics: docs/generality-check-zh.md,
% docs/quad-kneser-zh.md.
\section{General real saddle-node unfoldings}\label{sec:general}

Sections~\ref{sec:transition}--\ref{sec:char} use the exponential only through a
few constants.  This section states their results for a general real
saddle-node unfolding, lists the changes in the proofs, and adds one exact
invariance statement.  Tetration is the case $a_2=\tfrac12$, $\gamma=1$,
$\rho=\tfrac13$ below.

\subsection{Setting}
Let $f_\mu(w)$ be a family with the following properties.
\begin{enumerate}
\item[(H1)] $f_\mu(w)$ is entire in $w$, holomorphic in $\mu$ near $\mu_0$, and
  real for real $(\mu,w)$.
\item[(H2)] $f_{\mu_0}(w^*)=w^*$, $f_{\mu_0}'(w^*)=1$ and
  $a_2:=f_{\mu_0}''(w^*)/2>0$ (replace $w$ by $-w$ if necessary).
\item[(H3)] $\gamma:=-\partial_\mu f_\mu(w^*)|_{\mu_0}>0$ (reverse $\mu$ if
  necessary).  For $s:=\mu_0-\mu>0$ small there are two real fixed points
  $L<L_2$ near $w^*$, $L$ attracting with multiplier $\lambda$ and $L_2$ repelling
  with multiplier $\lambda_2$.
\item[(H4)] The base point $w_0<w^*$ is real, and $f_{\mu_0}$ is increasing with
  $f_{\mu_0}(w)>w$ on $[w_0,w^*)$.  Then for $s$ small, $w_0<L$ lies in the
  immediate basin of $L$, and the Koenigs map satisfies $\sigma(w_0)<0$.
\item[(H5)] $B_1\ne0$, where $B_n$ are the Fourier coefficients of the inverse upper
  horn map of the germ at $s=0$ (used only where indicated).
\end{enumerate}
Write the germ as $g(u)=f_{\mu_0}(w^*+u)-w^*=u+a_2u^2+a_3u^3+\cdots$, with
iterative residue $\rho=1-a_3/a_2^2$.  Its formal Abel function is
$\alpha(u)=-1/(a_2u)+\rho\log(-u)+\sum_{k\ge1}d_ku^k$, and $\Phi_{\mathrm{att}}$,
$\Phi_{\mathrm{rep}}$, $B_n$ are defined as in Section~\ref{sec:horn}.  The
normalisation $F(0)=1$ becomes $F(0)=w_0$: $\Rr^{-1}(w_0)=0$,
$u^*:=w_0-w^*$ and $a:=\Phi_{\mathrm{att}}(u^*)\in\mathbb R$.  The definitions of
$S$, $T$, $t_n$, $\tau_n$, $h$, $\Lambda$, $\Kk^W$ and $\mathcal C(Y)$ are those of
Sections~\ref{sec:transition} and~\ref{sec:char}, with $\mu$ in place of $b$.

No global hypothesis enters.  All estimates take place in a fixed disc about
$w^*$; (H1) is used only to make $S$ entire.

\begin{proposition}[Scales]\label{prop:gen-scale}
As $s\to0^+$,
\[
 |\log\lambda|,\ \log\lambda_2=2\sqrt{a_2\gamma s}\,\bigl(1+O(\sqrt s)\bigr),\qquad
 \Lambda=\exp\Bigl(-\frac{2\pi^2}{\sqrt{a_2\gamma s}}\bigl(1+o(1)\bigr)\Bigr),
\]
and $p:=|\log\lambda|\log\lambda_2$ is holomorphic in $s$ with
$p=4a_2\gamma s+O(s^2)$.  Moreover, with $A_1=1/\log\lambda$, $A_2=1/\log\lambda_2$
and $u_1=L-w^*$, $u_2=L_2-w^*$,
\[
 A_1+A_2\to\rho,\qquad A_1(u_1-u_2)\to1/a_2,\qquad h_2-h\to2\pi\rho,
\]
where $h_2=2\pi/\log\lambda_2$.
\end{proposition}

\begin{proof}
Write the Weierstrass preparation~\eqref{eq:weier} as
$f(u)-u=(u^2-\sigma_1u+\sigma_2)k_s(u)$ with $k_0(0)=a_2$.  Then
$\sigma_2=f_s(0)/a_2+O(s^2)=-\gamma s/a_2+O(s^2)$, so the discriminant is
$D(s)=4\gamma s/a_2+O(s^2)$, and $\lambda_{1,2}=1\mp a_2\sqrt D+O(D)$.  The
limits of $A_1+A_2$ and $A_1(u_1-u_2)$ are those of Lemma~\ref{lem:res}, with
$1-k'(0)/k(0)^2=1-a_3/a_2^2$ and $1/k(0)=1/a_2$.  Finally
$h_2-h=2\pi(A_1+A_2)$.  For tetration this gives $C_\eta$ of
Proposition~\ref{prop:saddle} and $p=2s+\cdots$.
\end{proof}

\subsection{The theorems}

\begin{theorem}[Confluence, general]\label{thm:gen-A}
Under \textup{(H1)--(H4)}, $\tau_n(\mu)\to B_n\eee^{2\pi\ii na}$ as $\mu\to\mu_0^-$ for every
$n\ge1$, with $\tau_n$ taken on the branch $\operatorname{Im}t_0=h/2$.  Under
\textup{(H5)} also $t_n/t_1^n\to B_n/B_1^n$.
\end{theorem}

\begin{theorem}[Sewing solution, general]\label{thm:gen-B}
Under \textup{(H1)--(H4)} there is $\mu_1<\mu_0$ such that for
$\mu\in(\mu_1,\mu_0)$ the following hold.  The sewn solution $\Kk^W_\mu$,
normalised by $\Kk^W_\mu(0)=w_0$, is the unique element of $\mathcal C_\mu(Y)$.  Its
coefficients satisfy $|c_n/\Lambda^n-\tau_n(\mu)|\le C_n\Lambda$, with $C_n$ as in
Theorem~\ref{thm:W2}(4) and constants $Y_B,Y_0,D,M$ that depend on the family.
The family is holomorphic in $\mu$ near $(\mu_1,\mu_0)$, and for $\operatorname{Im}\mu>0$
small it tends to $L$ and $L_2$ as $\operatorname{Im}z\to+\infty$ and $-\infty$,
respectively.
\end{theorem}

\begin{theorem}[First and higher orders, general]\label{thm:gen-D}
Under \textup{(H1)--(H4)}, and for $n$ with $B_n\ne0$,
\[
 \log\frac{\tau_n}{B_n\eee^{2\pi\ii na}}\sim\sum_{j\ge1}\kappa^{(n)}_jp^j,\qquad
 \kappa^{(n)}_1=\frac{\delta t_n}{4a_2\gamma\,B_n}-\frac{2\pi\ii n\,\delta t_0}{4a_2\gamma},
\]
with $\delta t_m$ defined as in Theorem~\ref{thm:D} from the orbits of the germ $g$
and of its reflected local inverse.  Each $\kappa^{(n)}_j$ is given by convergent
parabolic series as in Theorem~\ref{thm:E}.
\end{theorem}

\begin{proof}[Proof of Theorems~\ref{thm:gen-A}--\ref{thm:gen-D}]
The proofs of Sections~\ref{sec:transition}--\ref{sec:char} apply verbatim,
with the following replacements.
\begin{itemize}
\item In Lemmas~\ref{lem:res}--\ref{lem:petal}, Theorem~\ref{thm:C} and
  Lemma~\ref{lem:finite}, replace
  $\alpha_2=-2/u+\tfrac13\log(-u)$ by $-1/(a_2u)+\rho\log(-u)$ and $-2/u$ by
  $-1/(a_2u)$.  The limiting petal is $\{\operatorname{Re}(-1/(a_2u))>R\}$.
\item In Theorem~\ref{thm:C}(2), the reflected inverse
  $G_\mu(v)=-f_\mu^{-1}(-v)$ is the local inverse near $w^*$.  It is jointly
  holomorphic, its germ is $v+a_2v^2+(2a_2^2-a_3)v^3+\cdots$, and its iterative
  residue is $-\rho$.
\item In Lemma~\ref{lem:realT}, $1<L$ is replaced by \textup{(H4)}: $S$ maps
  $\mathbb R$ onto $(L,L_2)$, where $\sigma>0$, while $\sigma(w_0)<0$.
\item In Lemma~\ref{lem:band}, the bound $\pi/3$ becomes $\pi|\rho|$ and
  $h_2=h+2\pi/3+o(1)$ becomes $h_2=h+2\pi\rho+o(1)$.  Injectivity of $f_\mu$ on
  $D_{r'}$ follows from $f'_{\mu_0}(w^*)=1$ after shrinking $r'$.
\item In Proposition~\ref{prop:weld}, the singular set of $\Rr$ is the closed
  set $\Sigma_\Rr$ of inverse-branch images of the singular values of $f$.
\item In Lemma~\ref{lem:model}, $d_i=\tfrac i2u^{i-1}+\cdots$ becomes
  $d_i=ia_2u^{i-1}+\cdots$, and $D(s)=8s+\cdots$ becomes $4\gamma s/a_2+\cdots$.
  In Theorem~\ref{thm:D}, $p'(0)=2$ becomes $p'(0)=4a_2\gamma$
  (Proposition~\ref{prop:gen-scale}).
\item In Theorem~\ref{thm:W1}(2), $\lambda_{1,2}=1\mp a_2\sqrt{D(s)}+O(D)$ gives
  $\partial_\mu\lambda_1>0>\partial_\mu\lambda_2$, hence
  $\operatorname{Im}\mu>0\Rightarrow\arg\lambda_1>0>\arg\lambda_2$.
\item Statements involving ratios or logarithms use \textup{(H5)} in place of
  Proposition~\ref{prop:horn-nonzero}.  Corollary~\ref{cor:sewn-cusp-local}
  needs, in addition, that $w_0$ is not precritical for $f_{\mu_0}$.
\end{itemize}
Everything else is literally unchanged.  In particular the welding identity,
the Wiener-algebra contraction and the characterisation use only the two
Koenigs charts and the relation $\sigma\circ\Rr(w)=\sigma(w_0)\lambda^w$.
\end{proof}

\begin{remark}[On \textup{(H5)}]\label{rem:gen-B1}
The proof of Proposition~\ref{prop:horn-nonzero} gives $B_1\ne0$ whenever
$f_{\mu_0}$ restricted to its immediate parabolic basin has exactly one singular
value.  This covers $w^2+c$, whose basin contains the single critical point.
For $v-\sin v+c$ the critical values $2\pi k+1$ are real and only $k=0$ lies in
the basin, but the map is not of finite type, so applying the argument requires
a check.  For $w+(w^2-\epsilon)\eee^w$ we only know $B_1\ne0$ numerically.
\end{remark}

\subsection{Independence of the base point}

\begin{proposition}\label{prop:gen-base}
Let $w_0$ and $w_0'$ both satisfy \textup{(H4)}, and mark with a prime the
quantities computed from $w_0'$.  Then for every $\mu<\mu_0$ near $\mu_0$ and every
$n\ge1$,
\[
 \tau_n'(\mu)=\tau_n(\mu)\,\eee^{2\pi\ii n\Delta(\mu)},\qquad
 B_n\eee^{2\pi\ii na'}=B_n\eee^{2\pi\ii na}\,\eee^{2\pi\ii n\Delta_0},
\]
with $\Delta(\mu)=\Rr_\mu^{-1}(w_0')$ and $\Delta_0=\Phi_{\mathrm{att}}(u^{*\prime})-\Phi_{\mathrm{att}}(u^*)$
both real.  Hence $|\tau_n(\mu)|$, $|c_n|$ and $|\hat c_n|$ do not depend on the
base point, and neither does $\operatorname{Re}\kappa^{(n)}_j$ for any $j$.
\end{proposition}

\begin{proof}
By \textup{(H4)}, $\sigma(w_0)$ and $\sigma(w_0')$ are negative, so
$\Delta=\log(\sigma(w_0')/\sigma(w_0))/\log\lambda$ is real.  Then
$\Rr'^{-1}=\Rr^{-1}-\Delta$ and $T'=T-\Delta$, so $t_n'=t_n$ for $n\ne0$ and
$t_0'=t_0-\Delta$.  The branch condition $\operatorname{Im}t_0=h/2$ is preserved,
and $\tau_n'=\tau_n\eee^{2\pi\ii n\Delta}$.  At $\mu_0$, $u^*$ and $u^{*\prime}$ lie
on the real attracting axis, where $\Phi_{\mathrm{att}}$ is real.  The sewn
solutions for the two base points differ by a real translation, which does not
change $|c_n|$ (Remark~\ref{rem:invariance}).  Taking real parts of the logarithm of
the first identity divided by the second proves the last assertion.
\end{proof}

\begin{corollary}\label{cor:gen-invariant}
$|B_n|$ and the real parts $\operatorname{Re}\kappa^{(n)}_j$ are invariant under
orientation-preserving real-analytic conjugacy of the unfolding near
$(\mu_0,w^*)$ and under reparametrisation of $\mu$.
\end{corollary}

\begin{proof}
A conjugacy $\phi_\mu$ transports Koenigs charts, $\tilde\Rr=\phi\circ\Rr$ and
$\tilde S=\phi\circ S$.  So $\tilde T=T$ up to the normalising translations.  Take a
base point inside the domain of $\phi$ and apply
Proposition~\ref{prop:gen-base}.  The parameter $p$ depends only on the
multipliers.
\end{proof}

\begin{remark}[Numerical illustration]\label{rem:gen-numerics}
The script \texttt{generality\_check.py} implements the construction of this
section for any family given by $f$ and its Taylor coefficients at the fixed
points.  The table lists the values at $\lambda=0.97$, from runs over $\lambda=0.5,\dots,0.97$; the fourth column tends to $\operatorname{Re}\kappa^{(1)}$ as $p\to0$.
\begin{center}\small
\begin{tabular}{lllll}
\toprule
family & $\rho$ & $|B_1|$ & $\operatorname{Re}\log(\tau_1/B_1\eee^{2\pi\ii a})/p$ & $|\tau_1|/|B_1|-1$\\
\midrule
$b^w$ & $1/3$ & $0.0890584364122$ & $-0.0101990$ & $-9.4\cdot10^{-6}$\\
$w^2+c$ & $1$ & $0.0597942361690$ & $-0.0150481$ & $-1.4\cdot10^{-5}$\\
$v-\sin v+c$ & $1$ & $0.0723786416245$ & $-0.0142133$ & $-1.3\cdot10^{-5}$\\
$w+(w^2-\epsilon)\eee^w$ & $0$ & $2744.09379131$ & $+1.08542$ & $+1.0\cdot10^{-3}$\\
\bottomrule
\end{tabular}
\end{center}
In each family $\tau_n$, $t_n/t_1^n$ and $\log(\tau_1/B_1\eee^{2\pi\ii a})/p$ converge at
the rate $O(p)$.  The branch $\operatorname{Im}t_0=h/2$ holds to $10^{-31}$ or
better, and the sewing iteration of Theorem~\ref{thm:W2} gives
$\hat c_n=\tau_n(1+O(\Lambda))$.  For $w^2+c$ with the base points $w_0=0$, $0.2$ and
$0.35$, the real part in the fourth column agrees to all ten digits shown at each $\lambda$, as
Proposition~\ref{prop:gen-base} requires.  The imaginary part ranges from
$-0.07$ to $116$.  The germs of $w^2+c$ and $v-\sin v+c$ have the same
$\rho$ but different $|B_1|$ and $\operatorname{Re}\kappa^{(1)}$.
\end{remark}

\begin{remark}[Numerical evidence for the analogue of Theorem~\ref{mthm:F} for $w^2+c$]\label{rem:gen-quad}
For $f_c(w)=w^2+c$ the parameter $c=1/4$ is the cusp of the main cardioid, which
plays the role of the Shell--Thron region.  For real $c>1/4$ the two fixed points
are complex conjugate and repelling, and the real axis escapes.  The
two-fixed-point construction of Kneser type converges at the sampled
parameters with the normalisation $F(0)=f_c^3(0)$, corresponding to the intended
critical-point marking $F(-3)=0$ (script \texttt{quad\_kneser.py}).
The separate proof document \texttt{docs/proofs/quadratic-continuation.tex}
establishes its quadratic F and G using marked tilted domains, an explicit
polynomial multiplier deformation with open-overlap comparison, and an
independently replayed 4,013-box Arb exterior cover. Its conclusion concerns
regular height germs with this normalisation; it does not imply convergence
of the present numerical iteration. The following earlier observations
provide numerical evidence only.
\begin{itemize}
\item Along the arc $c=\tfrac14+0.1\eee^{\ii\varphi}$ the sampled values are
  consistent with continuation across the cardioid boundary.  Interpolation
  through points on both sides of the
  neutral band predicts held-out points to $2\cdot10^{-8}$, as accurately as
  interpolation within one side.
\item On vertical ladders $c=c_r+\ii y$, extrapolation towards $y=0$ supports
  agreement with the sewn coefficient.  At $\lambda=0.1$ the extrapolated value gives
  $(\hat c_1/\tau_1-1)/\Lambda=0.2272-0.1120\,\ii$, against $0.2279-0.1115\,\ii$
  from the first-order formula at the end of Section~\ref{sec:W2}.  So the
  data support agreement with $\Kk^W$ including its $O(\Lambda)$ correction;
  they do not prove existence or identification of the boundary limit.  The ratio
  $c_2/c_1^2$ agrees with $t_2/t_1^2$ to $4\cdot10^{-8}$.
\end{itemize}
For the analogue of Theorem~\ref{mthm:G} the numerics are only partial.  The
construction converges at sampled points in a wedge to the lower right of the
cusp.  In the multiplier coordinate $c=\mu/2-\mu^2/4$, at
$\arg\mu=0.25\pi$ and $0.4\pi$, interpolation across the neutral band errs by
$\le10^{-7}$ against a variation of order $10^{-1}$.  Further left the iteration
does not converge, and above it converges only after the $\theta$ lines are
raised in some tested cases.  These observations are consistent with local
holomorphic continuation, but do not establish it. Moving the Taylor centre
or fixing the constant Fourier mode of $\theta_{\mathrm{up}}$ may enlarge
the regular numerical region; these changes remain untested. They cannot
alone give direct coverage of the whole upper half-plane while retaining
both Schr\"oder charts: at $c=1/4+i/2$, $\lambda_{\mathrm{up}}=i$, the
fifth inverse-Schr\"oder coefficient requires $0=(3-5i)/2$.
The separate quadratic proof document supplies the exact obstruction and
a conditional local convergence theorem. A full numerical method must
also replace the neutral-end representation.
A pointwise algorithm that does not record paths does not by itself establish
single-valued analytic continuation; parameter holomorphy and agreement of
local branches must also be established.
\end{remark}


\section{Continuation of Kneser's solution around the cusp}\label{sec:cusp}

This section proves Theorem~\ref{mthm:F}, for the continuation that passes
near $\eta$.  The continuation leaves the real
exterior ray, crosses the upper Shell--Thron boundary arc close to the
cusp, and returns to the real axis below $\eta$ from the upper
half-plane.  Near the cusp the two fixed points are close, and Kneser's
initial region (the chord between the fixed points together with its
image) can be tilted continuously into an initial region for the
attracting--repelling configuration below $\eta$.  The multiplier of the
upper fixed point crosses the unit circle on the way, but this plays no
role.  In the logarithmic coordinate of the pair of fixed points the
tilted region is a thin strip on which the map is a translation up to a
relative error $O(|\varepsilon|)$, so its quotient is quasiconformally a flat
cylinder, whether or not the fixed points are linearisable.  At the end
of the path the tilted region is a fundamental domain for the sewn surface
of Proposition~\ref{prop:weld}, and a proper-map argument identifies the
continued solution with $\Kk^W_b$.

\emph{Conventions for this section.}  Here $\varepsilon$ is the complex parameter
$1-\lambda_1$ defined below; for real $b<\eta$ it differs from the quantity
$|\log\lambda|$ called $\varepsilon$ in Sections~\ref{sec:W2}--\ref{sec:char} by $O(\varepsilon^2)$, and
we use results of those sections only in the form stated there.  The letters
$\alpha$ (a tilt angle), $\varsigma$ (a normal coordinate), $s\in[0,1]$ (an
interpolation parameter), $H$, $X$, $M$, $q$ and $f$ have local meanings (in the subsection
``Consequences at the cusp'' $f=F_\eta^{\mathrm{att}}$).  $L_1$ denotes
the fixed point with multiplier $1-\varepsilon$: it is the upper fixed point $L_+$ for
real $b>\eta$ and the attracting fixed point $L$ for real $b<\eta$.  Constants
$C,c$ do not depend on $\varepsilon$; $r_0$ is decreased finitely many times.

\subsection{The cusp parameter}
For small complex $\varepsilon$ put
\[
 \lambda_1=1-\varepsilon,\qquad
 \log b(\varepsilon)=\lambda_1\eee^{-\lambda_1}=(1-\varepsilon)\eee^{\varepsilon-1},
\]
and let $\lambda_2(\varepsilon)$ be the other root of
$\lambda\eee^{-\lambda}=\log b(\varepsilon)$ near $1$.  It is holomorphic,
$\lambda_2=1+\varepsilon+\tfrac23\varepsilon^2+O(\varepsilon^3)$.  By
Lemma~\ref{lem:param}, $L_j=\eee^{\lambda_j}$ are fixed points of
$E_b$ with multipliers $\lambda_j$.  In the coordinate $w=\eee(1+u)$,
\[
 f_\varepsilon(u)=\eee^{\mu(1+u)-1}-1,\qquad \mu=\eee\log b=(1-\varepsilon)\eee^{\varepsilon},
\]
with fixed points $u_j=\eee^{\lambda_j-1}-1$, so that
$u_1-u_2=-2\varepsilon(1+O(\varepsilon))$.  For $0<r\le\tfrac1{10}$ let
\[
 \mathcal S^{\rm c}_r=\Bigl\{\varepsilon:\ 0<|\varepsilon|<r,\
   -\tfrac{7\pi}{12}<\arg\varepsilon<\tfrac\pi{12}\Bigr\},\qquad
 \mathcal S^{\rm c\prime}_r=\Bigl\{\varepsilon\in\mathcal S^{\rm c}_r:\
   -\tfrac{7\pi}{12}+\tfrac1{10}<\arg\varepsilon<\tfrac\pi{12}-\tfrac1{10}\Bigr\}.
\]
The lemmas are stated on $\mathcal S^{\rm c}_r$; the family is built on
$\mathcal S^{\rm c\prime}_r$, so that the parameter discs of
Proposition~\ref{prop:cusp-holo} stay in $\mathcal S^{\rm c}_r$.  Put
$u^*=\eee^{-1}-1$, the base point $w=1$, and $u^*_n=f_\varepsilon^n(u^*)$.

\begin{lemma}[Cusp parameter]\label{lem:cusp-param}
There are $r_0,r_1>0$ with the following properties.
\begin{enumerate}
\item $\varepsilon\mapsto b(\varepsilon)$ is injective on $\mathcal S^{\rm c}_{r_0}$.  For every
  $r\le r_0$ there is $r'>0$ such that $b(\mathcal S^{\rm c\prime}_r)$ contains
  $\{b:|b-\eta|<r',\ \operatorname{Im}b>0\}$ and the segments $(\eta-r',\eta)$ and
  $(\eta,\eta+r')$.  Real $b\in(\eta-r_1,\eta)$ corresponds to real
  $\varepsilon\in(0,r_0)$, and then $L_1=L$ is the attracting and $L_2$ the repelling
  real fixed point.  Real $b\in(\eta,\eta+r_1)$ corresponds to
  $\varepsilon(\theta)=1-\theta\cot\theta-\ii\theta=\tfrac13\theta^2-\ii\theta+O(\theta^4)$,
  $0<\theta<\theta_1$; then $L_1=L_+$ is the fixed point in the upper
  half-plane used in Proposition~\ref{prop:exterior-real-parameter}, and
  $\lambda_2(\varepsilon(\theta))=\overline{\lambda_1}$, $u_2=\overline{u_1}$.
\item On $D_{\rho_0}$, $\rho_0=\tfrac12$,
  $f_\varepsilon(u)-u=(u-u_1)(u-u_2)k_\varepsilon(u)$ with $k_\varepsilon$ jointly holomorphic in
  $(\varepsilon,u)$ and $|k_\varepsilon-\tfrac12|\le\tfrac14$.  Moreover $f_\varepsilon'(u_j)=\lambda_j$.
\item $g_\varepsilon(v)=(\operatorname{Log}(1+v)+1)/\mu-1$ is jointly holomorphic on
  $\{\operatorname{Re}(1+v)>0\}$, and $g_\varepsilon(f_\varepsilon(u))=u$ whenever
  $|\operatorname{Im}(\mu(1+u))|<\pi$, in particular on $D_{\rho_0}$ and near the real
  segment $[u^*,0]$.
\end{enumerate}
\end{lemma}

\begin{proof}
(1) $(1-\varepsilon)\eee^{\varepsilon}-1=-\tfrac12\varepsilon^2(1+O(\varepsilon))$, so
$\log b-\log\eta=-(2\eee)^{-1}\psi_0(\varepsilon)^2$ with $\psi_0$ holomorphic and
$\psi_0(\varepsilon)=\varepsilon(1+O(\varepsilon))$.  For small $r_0$, $\psi_0$ is injective and
maps $\mathcal S^{\rm c}_{r_0}$ into a sector of opening $2\pi/3+o(1)<\pi$, on which
squaring is injective.  The image of $\mathcal S^{\rm c\prime}_r$ under $-\psi_0^2$
contains, near $0$, all directions in $(-\tfrac\pi6+0.2+o(1),\tfrac{7\pi}6-0.2-o(1))\supset[0,\pi]$.
Since $b\mapsto\log b$ is biholomorphic near $\eta$, (1) follows except for
the identification of the real segments.  For real $\varepsilon\in(0,1)$,
$\lambda_1\in(0,1)$ and Lemma~\ref{lem:param} applies.  For real $b>\eta$ the upper
fixed point has multiplier $L_+\log b=\theta\cot\theta+\ii\theta$ with
$\theta=\operatorname{Im}(L_+\log b)\in(0,\pi)$ (the computation in the proof of
Proposition~\ref{prop:exterior-real-parameter}), and $\lambda_1=1-\varepsilon$ gives
$\varepsilon(\theta)$, whose argument lies in $(-\pi/2,-\pi/2+\theta)$.  The conjugate
fixed point $L_-=\overline{L_+}$ is the other one near $\eee$, so
$\lambda_2=\overline{\lambda_1}$.

(2) $k_0(u)=(\eee^u-1-u)/u^2=\tfrac12+\sum_{n\ge1}u^n/(n+2)!$ differs from
$\tfrac12$ by less than $0.1$ on $D_{1/2}$.  Division of $f_\varepsilon(u)-u$ by
$(u-u_1)(u-u_2)$ (Cauchy's formula on $|u|=1$) gives joint holomorphy and
continuity at $\varepsilon=0$.  Differentiating the product at $u_j$ gives the
multipliers.

(3) $\operatorname{Log}(\eee^{\mu(1+u)-1})=\mu(1+u)-1$ when $|\operatorname{Im}(\mu(1+u))|<\pi$.
\end{proof}

\subsection{The logarithmic coordinate}
Put $q=(u-u_1)/(u-u_2)$ and $\chi=\log q$, so that
\[
 u=u(\chi)=u_2+\frac{u_1-u_2}{1-\eee^{\chi}} .
\]
The points $\chi\in2\pi\ii\mathbb Z$ (the \emph{holes}) correspond to $u=\infty$.
Let $\mathcal G^{\rm c}=\{\chi:|\chi-2\pi\ii k|\ge\tfrac12\ \text{for all }k\in\mathbb Z\}$, and put
$A_1=1/\operatorname{Log}\lambda_1=-\varepsilon^{-1}(1+O(\varepsilon))$ and $Z=A_1\chi$.

\begin{lemma}[The step in the logarithmic coordinate]\label{lem:cusp-step}
For $0<|\varepsilon|<r_0$ the following hold (the proofs of (1)--(5) do not use the
sector).
\begin{enumerate}
\item On $\mathcal G^{\rm c}$, $|1-\eee^\chi|\ge\tfrac16\max(1,|\eee^\chi|)$; hence
  $|u(\chi)|\le C|\varepsilon|$ and
  $|u-u_1||u-u_2|\le C|\varepsilon|^2\min(|\eee^\chi|,|\eee^{-\chi}|)$.
  On $\{|\chi|\le1\}$, $|1-\eee^\chi|\ge0.63|\chi|$.
\item For $u\in D_{\rho_0}\setminus\{u_1,u_2\}$,
  \[
   \frac{q(f(u))}{q(u)}=1+\beta(u),\qquad
   \beta=\frac{(u_1-u_2)k(u)}{1+(u-u_1)k(u)},
  \]
  and $\Delta(u):=\operatorname{Log}(1+\beta(u))=-\varepsilon(1+\vartheta(u))$ with
  $|\vartheta|\le C(|u|+|\varepsilon|)$ and $|\Delta'(u)|\le C|\varepsilon|$ on $D_{\rho_0/2}$.  Thus
  $\chi(f(u))=\chi(u)+\Delta(u)$ on the branch continued along the short
  segment from $\chi(u)$, and $A_1\Delta=1+O(|u|+|\varepsilon|)$.
\item On $\mathcal G^{\rm c}$, regarded as functions of $\chi$,
  $|\vartheta|\le C|\varepsilon|$ and $|\partial_\chi\Delta|\le C|\varepsilon|^2$.  On
  $\{|\chi|\le1\}\cap\{|u(\chi)|\le\rho_0/2\}$, $|\partial_\chi\Delta|\le C|\varepsilon|^2/|\chi|^2$.
\item $\Delta(u_1)=\operatorname{Log}\lambda_1$ and $\Delta(u_2)=-\operatorname{Log}\lambda_2$.
  For $\chi\in\mathcal G^{\rm c}$, $|\Delta-\operatorname{Log}\lambda_1|+|\partial_\chi\Delta|\le
  C|\varepsilon|^2|\eee^\chi|$ if $\operatorname{Re}\chi\le0$, and
  $|\Delta+\operatorname{Log}\lambda_2|+|\partial_\chi\Delta|\le C|\varepsilon|^2|\eee^{-\chi}|$
  if $\operatorname{Re}\chi\ge0$.
\item $|\partial_\varepsilon\Delta|\le1+C|\varepsilon|$ at fixed $\chi\in\mathcal G^{\rm c}$.
\end{enumerate}
All functions here are jointly holomorphic in $\varepsilon$ and $u$ (or $\chi$).
\end{lemma}

\begin{proof}
(1) On the circle $|\chi|=\tfrac12$ the minimum of $|\eee^\chi-1|$ is
$1-\eee^{-1/2}>0.39$, attained at $\chi=-\tfrac12$.  The function
$1/(\eee^\chi-1)$ is holomorphic on the part of the strip
$|\operatorname{Im}\chi|\le\pi$ outside the disc, is bounded by $1$ on its horizontal
edges (where $|\eee^\chi-1|=1+\eee^{\operatorname{Re}\chi}$), and tends to $-1$ and $0$ at
the two ends; the maximum principle and periodicity give
$|\eee^\chi-1|\ge0.39$ on $\mathcal G^{\rm c}$.  For $\operatorname{Re}\chi\ge\log2$,
$|\eee^\chi-1|\ge\tfrac12|\eee^\chi|$, and for $0\le\operatorname{Re}\chi\le\log2$,
$0.39\ge\tfrac16|\eee^\chi|$.  The bounds on $u$ follow from
$u-u_2=(u_1-u_2)/(1-\eee^\chi)$ and $u-u_1=\eee^\chi(u-u_2)$.  The last bound
is the minimum of $|(\eee^\chi-1)/\chi|$ on the closed unit disc,
$1-\eee^{-1}$.

(2) The identity follows from $(f(u)-u_j)/(u-u_j)=1+(u-u_{3-j})k(u)$.  By
Lemma~\ref{lem:cusp-param}(2) and $u_1-u_2=-2\varepsilon(1+O(\varepsilon))$,
$\beta=-\varepsilon(1+O(|u|+|\varepsilon|))$, and $\operatorname{Log}(1+\beta)=\beta(1+O(|\beta|))$.
Cauchy's estimate on $D_{\rho_0}$, where $\Delta=O(|\varepsilon|)$, bounds $\Delta'$.

(3) The first bound is (2) with (1).  For the derivatives,
$\partial_\chi\Delta=\Delta'(u)\,du/d\chi$ with
$du/d\chi=(u_1-u_2)\eee^\chi/(1-\eee^\chi)^2$, which is $O(|\varepsilon|)$ on
$\mathcal G^{\rm c}$ by (1) and $O(|\varepsilon|/|\chi|^2)$ on $|\chi|\le1$.

(4) $1+\beta(u_1)=1+(u_1-u_2)k(u_1)=f'(u_1)=\lambda_1$, and
$1+\beta(u_2)=1/(1+(u_2-u_1)k(u_2))=1/\lambda_2$; all logarithms are
principal and close to $0$.  Since $\Delta'(u)=O(|\varepsilon|)$,
$|\Delta(u)-\Delta(u_1)|\le C|\varepsilon||u-u_1|\le C|\varepsilon|^2|\eee^\chi|$ by (1)
when $\operatorname{Re}\chi\le0$, and
$|\partial_\chi\Delta|=|\Delta'(u)|\,|u_1-u_2|\,|\eee^\chi|/|1-\eee^\chi|^2\le
C|\varepsilon|^2|\eee^\chi|$.  The case of $u_2$ is symmetric, with
$\eee^{-\chi}$ in place of $\eee^{\chi}$.

(5) At fixed $\chi$, $\Delta=-\varepsilon(1+\vartheta)$ with $\vartheta$ holomorphic in $\varepsilon$ and
$O(|\varepsilon|)$; Cauchy's estimate on a disc of radius $|\varepsilon|/2$ gives
$\varepsilon\,\partial_\varepsilon\vartheta=O(|\varepsilon|)$.
\end{proof}

\subsection{Tilted initial regions}
Write $\omega=\eee^{\ii\alpha}$.  For $x_0\in[-\tfrac12,\tfrac12]$ let
\[
 \ell=\ell_{\alpha,x_0}=\{\ii\pi+x_0+t\omega:\ t\in\mathbb R\},
\]
and use the coordinates $\chi=\ii\pi+x_0+(t+\ii\varsigma)\omega$, $t,\varsigma\in\mathbb R$.
We call $\ell$ \emph{admissible for $\varepsilon$} if
\begin{equation}\label{eq:admissible}
 \cos\alpha\ge\tfrac35,\qquad -\operatorname{Im}(\varepsilon\bar\omega)\ge\tfrac12|\varepsilon| .
\end{equation}
For $\varepsilon\in\mathcal S^{\rm c}_r$ put
$\alpha(\varepsilon)=\tfrac12(\arg\varepsilon+\tfrac\pi2)\in(-\tfrac\pi{24},\tfrac{7\pi}{24})$.
Then $\cos\alpha(\varepsilon)>\cos\tfrac{7\pi}{24}>0.6$ and
$-\operatorname{Im}(\varepsilon\bar\omega)/|\varepsilon|=\sin(\tfrac\pi4-\tfrac12\arg\varepsilon)>\sin\tfrac{5\pi}{24}>0.6$.
Hence, for $\varepsilon_0\in\mathcal S^{\rm c\prime}_r$, the line $\ell_{\alpha,x_0}$ with
$|\alpha-\alpha(\varepsilon_0)|\le\tfrac1{20}$ is admissible for every $\varepsilon\in\mathcal S^{\rm c}_r$ with
$|\arg\varepsilon-\arg\varepsilon_0|\le\tfrac1{20}$: then $\alpha\in(-\tfrac\pi{24},\tfrac{7\pi}{24})$, so
$\cos\alpha>\cos\tfrac{7\pi}{24}>\tfrac35$, and
$-\operatorname{Im}(\varepsilon\bar\omega)/|\varepsilon|\ge\sin(\tfrac{5\pi}{24}-\tfrac1{10})>\tfrac12$.  At
$\varepsilon(\theta)$ (real $b>\eta$), $\alpha=0$ is admissible; at real $\varepsilon>0$,
$\alpha(\varepsilon)=\tfrac\pi4$.  For fixed $\varepsilon$ the admissible pairs $(\alpha,x_0)$ form a
rectangle.  The distance from $2\pi\ii k$ to $\ell$ is
$|(2k-1)\pi\cos\alpha+x_0\sin\alpha|\ge\tfrac35\pi-\tfrac25>1.48$, so the strip
$\{|\varsigma|\le\tfrac14\}$ lies in $\mathcal G^{\rm c}$.

\begin{lemma}[Tilted initial region]\label{lem:crescent}
Let $\varepsilon\in\mathcal S^{\rm c}_{r_0}$, let $\ell$ be admissible for $\varepsilon$, and let
$f_\chi(\chi)=\chi+\Delta(\chi)$ be $f_\varepsilon$ in the coordinate $\chi$.
\begin{enumerate}
\item On $\mathcal G^{\rm c}$ the normal component
  $\operatorname{Im}(\Delta\bar\omega)$ lies in $[\tfrac25|\varepsilon|,\tfrac65|\varepsilon|]$, and $f_\chi$
  is injective on every convex subset of $\mathcal G^{\rm c}$.
\item $f_\chi(\ell)$ is the graph $\varsigma=g(t)$ of a real-analytic function with
  $\tfrac25|\varepsilon|\le g\le\tfrac65|\varepsilon|$ and $|g'|\le C|\varepsilon|^2$.
\item $M(t,s)=\ii\pi+x_0+t\omega+s\Delta(\ii\pi+x_0+t\omega)$ is an
  orientation-preserving real-analytic diffeomorphism of $\mathbb R\times[0,1]$
  onto the closed strip $\hat H=\{0\le\varsigma\le g(t)\}$ between $\ell$ and $f_\chi(\ell)$.
\item $\chi\mapsto u(\chi)$ is injective on $\hat H$.  Its image $H=H^\ell_\varepsilon$
  together with $u_1,u_2$ is a closed Jordan region in $D_{C|\varepsilon|}$, bounded by
  the arcs $\gamma=u(\ell)$ and $f(\gamma)=u(f_\chi(\ell))$, which join $u_2$
  ($t\to+\infty$) to $u_1$ ($t\to-\infty$) and meet only there.
\item $f_\varepsilon$ maps a neighbourhood of $\gamma$ biholomorphically onto a
  neighbourhood of $f(\gamma)$; in the coordinate $\chi$ it is $f_\chi$, which is
  injective on the convex strip $\{|\varsigma|\le\tfrac14\}$.
\end{enumerate}
\end{lemma}

\begin{proof}
(1) By Lemma~\ref{lem:cusp-step}(3),
$\operatorname{Im}(\Delta\bar\omega)=-\operatorname{Im}(\varepsilon\bar\omega)-\operatorname{Im}(\varepsilon\vartheta\bar\omega)
\in[\tfrac12|\varepsilon|-C|\varepsilon|^2,|\varepsilon|+C|\varepsilon|^2]$.  If $\chi,\chi'$ lie in a convex
subset of $\mathcal G^{\rm c}$, then
$|f_\chi(\chi)-f_\chi(\chi')-(\chi-\chi')|\le C|\varepsilon|^2|\chi-\chi'|<|\chi-\chi'|$.

(2) The $t$-coordinate of $f_\chi(\ii\pi+x_0+t\omega)$ is
$t+\operatorname{Re}(\Delta\bar\omega)$, whose derivative is $1+O(|\varepsilon|^2)>0$; it is
therefore a real-analytic bijection of $\mathbb R$, and (1) and
$|\partial_\chi\Delta|\le C|\varepsilon|^2$ bound $g$ and $g'$.

(3) $\partial_tM=\omega(1+O(|\varepsilon|^2))$ and $\partial_sM=\Delta$, so the Jacobian is
$\operatorname{Im}(\Delta\bar\omega)(1+O(|\varepsilon|^2))\ge\tfrac13|\varepsilon|>0$.  $M$ maps $s=0$ onto $\ell$ and
$s=1$ onto $f_\chi(\ell)$ bijectively, each segment $t=\mathrm{const}$ to a segment
of length $\le2|\varepsilon|$ joining them, and is proper.  A proper local
diffeomorphism that is a bijection on the boundary lines is a
diffeomorphism onto the region between them.

(4) Points of the strip $\{|\varsigma|\le\tfrac14\}$ that differ by
$2\pi\ii k$, $k\ne0$, have $\varsigma$-coordinates differing by
$2\pi|k|\cos\alpha>1$; so $\exp$ is injective on $\hat H$, and the Möbius map
$q\mapsto u$ is injective.  Along $\ell$,
$\operatorname{Re}\chi=x_0+t\cos\alpha-\varsigma\sin\alpha$ tends to $\mp\infty$ as $t\to\mp\infty$, so
$q\to0$ ($u\to u_1$) and $q\to\infty$ ($u\to u_2$).  The two boundary curves
of $\hat H$ are disjoint by (2), and their images are disjoint arcs with
common endpoints $u_1,u_2$; the image of $\hat H$ is the closure of the
complementary component they bound that contains $u(\operatorname{int}\hat H)$.
The inclusion in $D_{C|\varepsilon|}$ is Lemma~\ref{lem:cusp-step}(1).

(5) By Lemma~\ref{lem:cusp-step}(2), $f_\chi$ represents $f_\varepsilon$ in $\chi$, and (1)
applies to the convex strip.
\end{proof}

Let $X=X^\ell_\varepsilon$ be the Riemann surface obtained from $H$ by identifying
$u\in\gamma$ with $f_\varepsilon(u)\in f(\gamma)$, with charts near the seam given by
Lemma~\ref{lem:crescent}(5).  Since $\chi$ is injective on $\hat H$ we use it
(or $Z=A_1\chi$) as the coordinate on $H$.  We call the branch of $\chi$ obtained
by applying $f_\chi^{\pm1}$ to points of $\hat H$ the \emph{continued branch}.

\begin{lemma}[Flat model]\label{lem:crescent-quotient}
Let $\ell$ be admissible for $\varepsilon\in\mathcal S^{\rm c}_{r_0}$.  In the coordinate $Z$ let
$\hat H_0$ be the straight strip between the line $A_1\ell$ and $A_1\ell+1$, and
let $C_0=\hat H_0/(\zeta\sim\zeta+1)=\mathbb C/\mathbb Z$.  The map
\[
 \mathcal H\bigl(A_1(\ii\pi+x_0+t\omega)+s\bigr)=A_1M(t,s),\qquad t\in\mathbb R,\ 0\le s\le1,
\]
induces a quasiconformal homeomorphism $\mathcal H:C_0\to X$ with
$|\mathcal H(\zeta)-\zeta|\le C|\varepsilon|$ on $\hat H_0$ and Beltrami coefficient
$|\upsilon_{\mathcal H}|\le C|\varepsilon|$.  Consequently:
\begin{enumerate}
\item $X$ is biholomorphic to $\mathbb C^*$, the ends at $u_1$ and $u_2$
  corresponding to the punctures.  Let $Z_X:X\to\mathbb C^*$ be a uniformisation
  with the $u_1$-end at $0$, and $z=(2\pi\ii)^{-1}\log Z_X$.  Crossing the seam
  from $\gamma$ to $f(\gamma)$ increases $z$ by $1$.  The inverse $z\mapsto u$, extended
  by $\tilde F(z+1)=f_\varepsilon(\tilde F(z))$ and by one application of $f_\chi^{-1}$
  on the strip $\{|\varsigma|\le\tfrac14\}$, is a holomorphic solution of this
  equation on the region to the right of the lift of $f_\chi^{-1}(\ell)$.
\item If $\mathsf p\in X$ and the lift $z$ is normalised at $\mathsf p$, then for
  $|\operatorname{Im}(z-z(\mathsf p))|\le1$ and $|\operatorname{Re}(z-z(\mathsf p))|\le2$ the
  $Z$-coordinate $Z(z)$ of $\tilde F(z)$ on the continued branch satisfies
  $|Z(z)-Z(\mathsf p)-(z-z(\mathsf p))|\le C|\varepsilon|$.
\end{enumerate}
\end{lemma}

\begin{proof}
$A_1\Delta=1+\vartheta_1$ with $|\vartheta_1|\le C|\varepsilon|$ on $\mathcal G^{\rm c}$
(Lemma~\ref{lem:cusp-step}(2),(3)).  In the source coordinate
$\zeta=A_1(\ii\pi+x_0+t\omega)+s$ one has $\mathcal H(\zeta)=\zeta+s\vartheta_1$, so
$|\mathcal H(\zeta)-\zeta|\le C|\varepsilon|$.  The partial derivatives are
$\partial_t\mathcal H=A_1\omega(1+O(|\varepsilon|^2))$ and $\partial_s\mathcal H=1+\vartheta_1$, against
$A_1\omega$ and $1$ for the identity.  The unit vectors $A_1\omega/|A_1\omega|$ and $1$
make an angle whose sine is
$|\operatorname{Im}(\omega\overline{\operatorname{Log}\lambda_1})|/|\operatorname{Log}\lambda_1|\ge\tfrac12-C|\varepsilon|$, by
$\operatorname{Log}\lambda_1=-\varepsilon(1+O(\varepsilon))$ and~\eqref{eq:admissible}.  So, in this basis,
$D\mathcal H=I+O(|\varepsilon|)$ and $|\upsilon_{\mathcal H}|\le C|\varepsilon|$.  $\mathcal H$ is a homeomorphism
of the strips by Lemma~\ref{lem:crescent}(3); it maps the edge $s=0$ to
$A_1\ell$ and $s=1$ to $A_1f_\chi(\ell)$ with the same $t$, and hence respects the
identifications $\zeta\sim\zeta+1$ and $u\sim f_\varepsilon(u)$.  Across $k$ seams,
$\mathcal H(\zeta+k)$ is represented on the continued branch by
$A_1f_\chi^k(\mathcal H(\zeta)/A_1)$, and $A_1f_\chi^k(\chi)-A_1\chi-k=O(|k||\varepsilon|)$ for
$|k|\le3$ by Lemma~\ref{lem:cusp-step}(2).

(1) $C_0\cong\mathbb C^*$.  Each end of $X$ is the image of an end of $C_0$, a
punctured disc; quasiconformal maps preserve the infinite modulus of an
annular end~\cite{LehtoVirtanen}, so each end of $X$ is a punctured disc and
$X\cong\mathbb C^*$.  In $C_0$ the passage from the edge $s=0$ to $s=1$ adds $1$.
The $u_1$-end is $t\to-\infty$, where
$\operatorname{Im}\zeta=t\operatorname{Im}(A_1\omega)+O(1)\to+\infty$ because
$\operatorname{Im}(A_1\omega)=-\operatorname{Im}(\operatorname{Log}\lambda_1\bar\omega)/|\operatorname{Log}\lambda_1|^2<0$.  Since $\mathcal H$
preserves orientation and ends, the uniformisation with the $u_1$-end at $0$
has the same properties.  The last assertion is the usual lift of the
uniformisation; the extra step $f_\chi^{-1}$ is allowed by Lemma~\ref{lem:crescent}(5).

(2) Put $\xi=\exp(2\pi\ii(\zeta-\zeta_{\mathsf p}))$ on $C_0$, $\zeta_{\mathsf p}=\mathcal H^{-1}(\mathsf p)$, and
let $\upsilon$ be the Beltrami coefficient of $\mathcal H$ transported to $\hat{\mathbb C}$.  For
complex $\mathsf t$ with $|\mathsf t|<R:=1/(2\|\upsilon\|_\infty)$ let $\mathcal W_{\mathsf t}$ be the solution
with coefficient $\mathsf t\upsilon$ fixing $0,1,\infty$.  By~\cite[Thm.~11]{AhlforsBers1960},
$\mathsf t\mapsto\mathcal W_{\mathsf t}(\xi)$ is holomorphic; $\mathcal W_0=\mathrm{id}$; and the $\mathcal W_{\mathsf t}$ are
$3$-quasiconformal and normalised, hence uniformly bounded, by some $M_0$, on
the annulus $A=\{\eee^{-3\pi}\le|\xi|\le\eee^{3\pi}\}$~\cite[Ch.~II]{LehtoVirtanen}.
Schwarz's lemma applied to $\mathsf t\mapsto\mathcal W_{\mathsf t}(\xi)-\xi$ on $|\mathsf t|<R$ gives
$|\mathcal W_1(\xi)-\xi|\le2(M_0+\eee^{3\pi})\|\upsilon\|_\infty\le C|\varepsilon|$ on $A$, since $R\ge1$
for small $|\varepsilon|$.  By a degree argument on the boundary circles of $A$,
$\mathcal W_1^{-1}$ also moves points of $\{\eee^{-2\pi}\le|\xi|\le\eee^{2\pi}\}$ by $O(|\varepsilon|)$.
$\mathcal W_1\circ\exp(2\pi\ii(\mathcal H^{-1}-\zeta_{\mathsf p}))$ is a uniformisation of $X$ normalised at
$\mathsf p$, so it equals $\exp(2\pi\ii(z-z(\mathsf p)))$.  Taking logarithms on the
annulus, and combining with the bounds on $\mathcal H$ across at most three seams,
gives (2).
\end{proof}

\begin{lemma}[The base orbit and the marked point]\label{lem:cusp-base}
There are $r_2\le r_0$, a small $\rho_1>0$ and $N_1$ such that the following holds
for $\varepsilon\in\mathcal S^{\rm c}_{r_2}$ and every line $\ell$ admissible for $\varepsilon$.  Let
$\chi^*_n=\chi(u^*_n)$, $n\ge N_1$, be the branch continued from the principal
value at $n=N_1$ (so $\chi^*_{n+1}=f_\chi(\chi^*_n)$), and $Z^*_n=A_1\chi^*_n$.
\begin{enumerate}
\item For $N_1\le n\le N_1+10/|\varepsilon|$,
  $Z^*_n=Z^*_{N_1}+(n-N_1)+\mathcal E_n$ with $|\mathcal E_n|\le C\log n+C|\varepsilon|n$, and
  $Z^*_{N_1}$ is within relative distance $\tfrac15$ of a real number $\ge4/\rho_1$.
\item On that range the normal coordinate $\varsigma(\chi^*_n)$ is strictly
  increasing.  Let $n_0$ be the first index with $\varsigma(\chi^*_{n_0})\ge0$.  Then
  $n_0\le N_1+8/|\varepsilon|$ and $\chi^*_{n_0}\in\hat H\setminus f_\chi(\ell)$.
\item For $N_1\le n\le n_0+1$, the $3|\varepsilon|$-neighbourhood of $\chi^*_n$ (equivalently,
  approximately, the $3$-neighbourhood of $Z^*_n$) lies in $\mathcal G^{\rm c}\cup\{|\chi|<1\}$
  and in $\{|u(\chi)|\le2\rho_1\}$.
\end{enumerate}
The \emph{marked point} of $X^\ell_\varepsilon$ is the point represented by
$\chi^*_{n_0}\in\hat H$.
\end{lemma}

\begin{proof}
At $\varepsilon=0$ the orbit $f_0^n(u^*)$ increases to $0$ along the negative axis,
with $f_0^n(u^*)=-2/n+O(\log n/n^2)$.  Choose $N_1$ with
$|f_0^{N_1}(u^*)|\le\rho_1/2$, and then $r_2$ so small that
$|u^*_{N_1}-f_0^{N_1}(u^*)|\le\tfrac1{10}|f_0^{N_1}(u^*)|$ for $|\varepsilon|<r_2$.
Since $q(u)=1+(u_2-u_1)/(u-u_2)$,
$Z^*_{N_1}=-2/u^*_{N_1}\,(1+O(|\varepsilon|/\rho_1)+O(\varepsilon))$, which gives the statement
about $Z^*_{N_1}$.  By Lemma~\ref{lem:cusp-step}(2),
$Z^*_{n+1}-Z^*_n=A_1\Delta(u^*_n)=1+e_n$ with $|e_n|\le C(|u^*_n|+|\varepsilon|)$ while
$u^*_n\in D_{\rho_0}$.  While $|\chi|\le1$, Lemma~\ref{lem:cusp-step}(1) gives
$|u|\le|u_2|+|u_1-u_2|/(0.63|\chi|)\le3.2/|Z|+C|\varepsilon|$; on $\mathcal G^{\rm c}$, $|u|\le C|\varepsilon|$.
By induction $\operatorname{Re}Z^*_n$ increases by at least $\tfrac12$ per step, and
$|Z|\ge|Z^*_{N_1}|-3\ge3/\rho_1$ on the $3$-neighbourhoods, so $|u|\le1.1\rho_1+C|\varepsilon|$
there; this gives (1) and the last part of (3).  Multiplying by
$1/A_1=-\varepsilon(1+O(\varepsilon))$, $\chi^*_n$ lies within $C|\varepsilon|\log(1/|\varepsilon|)+C|\varepsilon|/\rho_1$ of the
ray $\{-\varepsilon\tau:\tau\ge0\}$ in the range of (1).

The origin has $\varsigma(0)=-\pi\cos\alpha+x_0\sin\alpha\le-1.48$ and
$|\varsigma(0)|\le\sqrt{\pi^2+1/4}$, and along the ray $\varsigma$ increases at rate
$-\operatorname{Im}(\varepsilon\bar\omega)\ge|\varepsilon|/2$; so the ray meets $\ell$ at $\tau\le6.4/|\varepsilon|$.
Before that crossing the ray stays at distance $\ge2\pi\sin\tfrac{5\pi}{12}>6$ from
$2\pi\ii k$, $k\le-1$, and in $\{\varsigma<0\}$, at distance $\ge1.48$ from $2\pi\ii k$,
$k\ge1$.  With the smallness of the deviation this gives (3).  The normal
component of each step is at least $|\varepsilon|(\tfrac12-C\rho_1-C|\varepsilon|)>0$, so
$\varsigma(\chi^*_n)$ is strictly increasing, and $n_0$ exists with the stated bound.
Then $-\tfrac65|\varepsilon|\le\varsigma(\chi^*_{n_0-1})<0$, so $\chi^*_{n_0-1}$ lies in the convex strip
$\{|\varsigma|\le\tfrac14\}$, on which each $\mathrm{id}+s\Delta$, $0\le s\le1$, is injective.
This isotopy carries $\ell$ to $f_\chi(\ell)$ and the side $\{\varsigma<0\}$ of $\ell$ into
the corresponding side of $f_\chi(\ell)$; hence $\chi^*_{n_0}\in\hat H\setminus f_\chi(\ell)$.
\end{proof}

The physical point $u^*_{n}$ may lie in $H$ for several $n$, because $H$
spirals around $u_1$ and $u_2$ when $\alpha\ne0$.  These other visits are points of
$X$ different from the marked point, and correspond to the translated roots
$c_j$ discussed after Definition~\ref{def:class}.  Only the first crossing on
the continued branch is used.

\begin{lemma}[Shadowing the base orbit]\label{lem:cusp-shadow}
Let $W_0=\{z:\operatorname{dist}(z,[0,1])<\tfrac14\}$.  Under the hypotheses of
Lemma~\ref{lem:cusp-base}, let $z$ be the lift of Lemma~\ref{lem:crescent-quotient}
normalised by $z=n_0$ at the marked point.  Then for $0\le k\le n_0-N_1$ the
branch $g_\varepsilon^{k}$ is holomorphic on $\tilde F(n_0+W_0)$, and the $Z$-coordinates
of $g_\varepsilon^k(\tilde F(n_0+w))$ lie within $\tfrac1{10}$ of $Z^*_{n_0-k}+w$ ($w\in W_0$).
For $|\varepsilon|$ small the $N_1$ further applications of the same branch $g_\varepsilon$
(principal logarithm) are defined and jointly holomorphic on
$g_\varepsilon^{n_0-N_1}(\tilde F(n_0+W_0))$ and its images, which stay in
$\{\operatorname{Re}(1+v)>0\}$ near the real segment $[u^*,0]$.
\end{lemma}

\begin{proof}
By Lemma~\ref{lem:crescent-quotient}(2) with $\mathsf p$ the marked point, the
$Z$-coordinates of $\tilde F(n_0+w)$ lie within $C|\varepsilon|$ of $Z^*_{n_0}+w$.  In the
coordinate $Z$, $f_\varepsilon$ is $Z\mapsto Z+A_1\Delta$, whose derivative is
$1+\partial_\chi\Delta$, and $g_\varepsilon$ is its inverse $g_Z$, with
$|g_Z'-1|\le2|\partial_\chi\Delta|$ evaluated at $g_Z(Z)$.  By Lemma~\ref{lem:cusp-step}(3)
and $|\chi|=|Z|/|A_1|$, $|g_Z'-1|\le\delta(Z):=C|\varepsilon|^2+C/|Z|^2$ on the
$3$-neighbourhood of the orbit segment, where Lemma~\ref{lem:cusp-base}(3)
applies and $|Z|\ge|Z^*_{N_1}|-3$.  Since $g_Z^k(Z^*_{n_0})=Z^*_{n_0-k}$, the
mean-value inequality gives, as long as the images stay in that
neighbourhood,
\[
 |g_Z^k(Z)-Z^*_{n_0-k}-(Z-Z^*_{n_0})|\le\Bigl(\prod_{j<k}\bigl(1+2\delta(Z^*_{n_0-j})\bigr)-1\Bigr)|Z-Z^*_{n_0}|
 \le\bigl(\eee^{C|\varepsilon|+C\rho_1}-1\bigr)\cdot2,
\]
because on the $3$-neighbourhood of $Z^*_{n_0-j}$ one has $\delta\le2\delta(Z^*_{n_0-j})$,
using $\sum_j|\varepsilon|^2\le C|\varepsilon|$ over $n_0\le N_1+8/|\varepsilon|$ steps and
$\sum_j|Z^*_{n_0-j}|^{-2}\le C/|Z^*_{N_1}|\le C\rho_1$ (the real parts grow by at least
$\tfrac12$ per step).  For small $\rho_1$ and $|\varepsilon|$ the right side is below $\tfrac1{20}$,
so by induction the images stay in the neighbourhood and the claim holds.
These sets lie in $\{|u|\le2\rho_1\}\subset D_{\rho_0}$, where
$g_\varepsilon\circ f_\varepsilon=\mathrm{id}$ (Lemma~\ref{lem:cusp-param}(3)).

At $k=n_0-N_1$ the set lies within $\tfrac1{10}$, in $Z$, of $Z^*_{N_1}+W_0$.  As
$\varepsilon\to0$, $Z=A_1\chi$ converges on this region to $-2/u$ locally uniformly, so the
set lies, for small $|\varepsilon|$, in a neighbourhood of the compact set
$\overline{K_0}$, $K_0=\{u:-2/u\in Z_0+W_0^{1/5}\}$, where $Z_0=-2/f_0^{N_1}(u^*)$ and
$W_0^{1/5}$ is the $\tfrac15$-neighbourhood of $W_0$.  At $\varepsilon=0$ the attracting Fatou
coordinate is $\Phi_{\mathrm{att}}(u)=-2/u+\tfrac13\log(-u)+O(u)$; on $\overline{K_0}$ the term
$\tfrac13\log(-u)$ varies by $O(\rho_1)$, so
$\Phi_{\mathrm{att}}(\overline{K_0})\subset c_0+N_1+W_0^{1/5+C\rho_1}$ for a real constant $c_0$.
The Fatou parametrisation $P_0=\Phi_{\mathrm{att}}^{-1}$ extends to an entire function with
$P_0(z+1)=f_0(P_0(z))$, and $g_0(P_0(z))=P_0(z-1)$ near the real axis, where
$1+P_0(z-1)$ is close to $(0,1]$.  Hence $g_0^{N_1}$ is holomorphic on a
neighbourhood of $\overline{K_0}$, and all intermediate images lie in a compact subset
of $\{\operatorname{Re}(1+v)>0\}$ near $[u^*,0]$.  This is an open condition, and
$g_\varepsilon\to g_0$ uniformly on compact subsets of $\{\operatorname{Re}(1+v)>0\}$; by
compactness it persists for small $|\varepsilon|$, jointly holomorphically in $\varepsilon$.
\end{proof}

\begin{definition}\label{def:cusp-solution}
For $\varepsilon\in\mathcal S^{\rm c}_{r_2}$ and an admissible line $\ell$, the \emph{normalised
tilted solution} is
\[
 F^\ell_\varepsilon(z)=g_\varepsilon^{\,n_0}\bigl(\tilde F(z+n_0)\bigr),\qquad z\in W_0,
\]
with $\tilde F$ lifted so that the marked point has $z=n_0$ and $g_\varepsilon$ the
principal branch of Lemma~\ref{lem:cusp-param}(3), which is holomorphic along
the way by Lemma~\ref{lem:cusp-shadow}.  It is holomorphic on $W_0$,
$F^\ell_\varepsilon(0)=u^*$, and $F^\ell_\varepsilon(z+1)=f_\varepsilon(F^\ell_\varepsilon(z))$ where both sides are
defined; in $w=\eee(1+u)$ it satisfies $F(z+1)=b^{F(z)}$ and $F(0)=1$.
\end{definition}

\begin{lemma}[Change of line]\label{lem:change-line}
Let $\ell,\ell'$ be admissible for the same $\varepsilon\in\mathcal S^{\rm c}_{r_2}$.  Then
$F^\ell_\varepsilon=F^{\ell'}_\varepsilon$.
\end{lemma}

\begin{proof}
Since the admissible pairs form a rectangle, it suffices to treat
$|\alpha-\alpha'|\le\tfrac1{100}$ and $|x_0-x_0'|\le\tfrac1{100}$.  Let
$\chi=\ii\pi+x_0+(t+\ii\varsigma)\omega\in\hat H$ and let $\varsigma'$ be its normal coordinate
for $\ell'$.  Then $|\varsigma'|\le\tfrac1{100}(1+|t|)+O(|\varepsilon|)$.  Each application of
$f_\chi^{\pm1}$ changes $\varsigma'$ by at least $\tfrac25|\varepsilon|$ in the same direction
(Lemma~\ref{lem:crescent}(1)) and moves the point by at most $2|\varepsilon|$.  So
after at most $(1+|t|)/(40|\varepsilon|)$ steps, forward or backward, the orbit of $\chi$
on the continued branch reaches the strip $\hat H'$ of $\ell'$, within distance
$(1+|t|)/20$ of $\chi$.  The distance from $\chi$ to the holes is at least
$\max(1.48,\,0.6|t|-0.5)$, which exceeds $(1+|t|)/20+\tfrac12$; so this part of the
orbit stays in a disc about $\chi$ contained in $\mathcal G^{\rm c}$, on which $f_\chi$ is
injective.  As in the proof of Lemma~\ref{lem:cusp-base}(2), exactly one
iterate $f_\chi^{n(\chi)}(\chi)$ lies in $\hat H'\setminus f_\chi(\ell')$.  The map
$\chi\mapsto f_\chi^{n(\chi)}(\chi)$ is locally a fixed iterate of $f_\chi$, it jumps by one
iterate exactly across $\ell'$, where the two values are identified in
$X^{\ell'}$, and it satisfies $n(f_\chi(\chi))=n(\chi)-1$ on $\ell$.  It therefore induces
a holomorphic map $X^\ell\to X^{\ell'}$, and the same construction from $\ell'$ gives
its inverse.  It preserves the two ends.  The continued base orbit
$\chi^*_n$ is common to both lines, and both marked points are its first
crossings; the orbit segment between them lies in the region just
described, so the map sends the marked point $\chi^*_{n_0}$ of $\ell$ to
$f_\chi^{n_0'-n_0}(\chi^*_{n_0})=\chi^*_{n_0'}$, the marked point of $\ell'$.  In the lifts
normalised at the marked points ($z=n_0$ and $z=n_0'$), and since $f_\chi$
corresponds to $z\mapsto z+1$, one gets $\tilde F'(z+n_0')=f_\varepsilon^{n_0'-n_0}(\tilde F(z+n_0))$
near $z=0$.  Applying the inverse branches of Definition~\ref{def:cusp-solution}
gives $F^{\ell'}_\varepsilon=F^\ell_\varepsilon$ on $W_0$.
\end{proof}

We write $F_\varepsilon$ for the common normalised tilted solution.

\subsection{Holomorphic dependence}

\begin{proposition}[The tilted family is holomorphic]\label{prop:cusp-holo}
There is $r_3\le\min(r_2,1/240)$ such that $(\varepsilon,z)\mapsto F_\varepsilon(z)$ is jointly holomorphic on
$\mathcal S^{\rm c\prime}_{r_3}\times W_0$.
\end{proposition}

\begin{proof}
Fix $\varepsilon_0\in\mathcal S^{\rm c\prime}_{r_3}$ and the disc $U_0=\{|\varepsilon-\varepsilon_0|<|\varepsilon_0|/40\}$;
it lies in $\mathcal S^{\rm c}_{r_2}$ and $|\arg\varepsilon-\arg\varepsilon_0|<\tfrac1{20}$ on it.  Choose
$\alpha\ne0$ with $|\alpha-\alpha(\varepsilon_0)|\le\tfrac1{20}$, and then $x_0\in[-\tfrac12,\tfrac12]$, such that
$\chi^*_{n_0}(\varepsilon_0)$ is not on $\ell$ or $f_\chi(\ell)$; for $\alpha\ne0$ a change of $x_0$ moves the
line in its normal direction by $x_0\sin\alpha$, so all but finitely many small
values of $x_0$ will do.  By the admissibility remark above $\ell=\ell_{\alpha,x_0}$ is
admissible for all $\varepsilon\in U_0$.  Shrink $U_0$ so that $n_0$ is constant and
$\chi^*_{n_0}(\varepsilon)$ stays in the interior of $\hat H_\varepsilon$.  By Lemma~\ref{lem:change-line}
this choice does not affect $F_\varepsilon$.

The local-triviality argument of Proposition~\ref{prop:weld-parameter} is not
available here, because at neutral parameters the ends of $X_\varepsilon$ have no a
priori holomorphic puncture coordinate; we use the Beltrami equation instead.
Let $M_\varepsilon$ be the diffeomorphism of Lemma~\ref{lem:crescent}(3) for $\ell$.  The map
$\mathcal J_\varepsilon=M_\varepsilon\circ M_{\varepsilon_0}^{-1}:\hat H_{\varepsilon_0}\to\hat H_\varepsilon$ respects the seam
identifications and induces a homeomorphism $X_{\varepsilon_0}\to X_\varepsilon$.  At each point
its derivative sends $v_1=\partial_tM_{\varepsilon_0}=\omega(1+O(|\varepsilon_0|^2))$ to
$w_1=\partial_tM_\varepsilon=\omega(1+O(|\varepsilon_0|^2))$ and $v_2=\Delta_{\varepsilon_0}$ to $w_2=\Delta_\varepsilon$, with
$|w_2-v_2|\le1.1|\varepsilon-\varepsilon_0|\le0.03|\varepsilon_0|$ by Lemma~\ref{lem:cusp-step}(5).  Writing
it as $\zeta\mapsto\mathsf a\zeta+\mathsf b\bar\zeta$,
\[
 \mathsf b=\frac{(w_1-v_1)v_2-(w_2-v_2)v_1}{2\ii\operatorname{Im}(v_2\bar v_1)},\qquad
 |\operatorname{Im}(v_2\bar v_1)|\ge\tfrac25|\varepsilon_0|(1-C|\varepsilon_0|),
\]
so $|\mathsf b|\le0.05$, $|\mathsf a-1|\le0.05$, and the Beltrami coefficient
$\upsilon_\varepsilon=\mathsf b/\mathsf a$ of $\mathcal J_\varepsilon$ satisfies $|\upsilon_\varepsilon|\le\tfrac1{10}$.  It is holomorphic in $\varepsilon$ at
each point, since $\mathsf a,\mathsf b$ are affine in $w_1,w_2$; being uniformly bounded,
it is holomorphic as an $L^\infty$-valued function.  Transport $\upsilon_\varepsilon$ to $\hat{\mathbb C}$
by a fixed uniformisation $Z_{\varepsilon_0}$ of $X_{\varepsilon_0}$, and let $\mathcal W_\varepsilon$ be the normalised
solution fixing $0,\infty$ (the ends) and $1$.  By~\cite[Thm.~11]{AhlforsBers1960},
$\mathcal W_\varepsilon(\xi)$ is holomorphic in $\varepsilon$ for each $\xi$, and
$Z_\varepsilon=\mathcal W_\varepsilon\circ Z_{\varepsilon_0}\circ\mathcal J_\varepsilon^{-1}$ is a uniformisation of $X_\varepsilon$ with the
ends at $0,\infty$.

In the interior of $\hat H_{\varepsilon_0}$, $\mathcal J_\varepsilon$ is a real-analytic diffeomorphism,
holomorphic in $\varepsilon$ at fixed source points, and $\upsilon_\varepsilon$ is real-analytic in the
source point and holomorphic in $\varepsilon$.  By the interior Schauder estimates for
the Beltrami equation~\cite[Ch.~V]{AhlforsLectures}, applied to difference
quotients in $\varepsilon$, $(\varepsilon,\xi)\mapsto\mathcal W_\varepsilon(\xi)$ is $C^1$ there.  The maps $\mathcal W_\varepsilon$ and
$\mathcal J_\varepsilon\circ Z_{\varepsilon_0}^{-1}$ have the same Beltrami coefficient $\upsilon$ (transported),
i.e.\ $Z_\varepsilon$ is conformal.  Put $\xi_\varepsilon(\chi)=Z_{\varepsilon_0}(\mathcal J_\varepsilon^{-1}(\chi))$.  Differentiating
$\mathcal J_\varepsilon(Z_{\varepsilon_0}^{-1}(\xi_\varepsilon(\chi)))=\chi$ in $\bar\varepsilon$ at fixed $\chi$, and using that
$\mathcal J_\varepsilon$ is holomorphic in $\varepsilon$ at fixed source points, gives
$\partial_{\bar\varepsilon}\xi_\varepsilon+\upsilon\,\partial_{\bar\varepsilon}\overline{\xi_\varepsilon}=0$.  With
$\partial_{\bar\xi}\mathcal W_\varepsilon=\upsilon\,\partial_\xi\mathcal W_\varepsilon$ and $\partial_{\bar\varepsilon}\mathcal W_\varepsilon=0$ this gives
$\partial_{\bar\varepsilon}[\mathcal W_\varepsilon(\xi_\varepsilon(\chi))]=\partial_\xi\mathcal W_\varepsilon\,(\partial_{\bar\varepsilon}\xi_\varepsilon+\upsilon\,\partial_{\bar\varepsilon}\overline{\xi_\varepsilon})=0$.
So $(\varepsilon,\chi)\mapsto Z_\varepsilon(\chi)$ is jointly holomorphic in the interior of the strips.

With $\chi^*_{n_0}(\varepsilon)$ in the interior,
$\zeta_0(\varepsilon)=(2\pi\ii)^{-1}\log Z_\varepsilon(\chi^*_{n_0}(\varepsilon))$ is holomorphic, and
$\tilde F_\varepsilon(z+n_0)=u_\varepsilon\bigl(Z_\varepsilon^{-1}(\eee^{2\pi\ii(z+\zeta_0(\varepsilon))})\bigr)$ is jointly
holomorphic wherever the point lies in the interior of $\hat H_\varepsilon$, by the
holomorphic inverse function theorem.  Points of $n_0+W_0$ whose images lie
on the seam, or in the extra step of Lemma~\ref{lem:crescent-quotient}(1), are
handled in the same way with finitely many further admissible angles $\alpha$
(by compactness of $\overline{W_0}$), for which they are interior; Lemma~\ref{lem:change-line} shows that the two descriptions
define the same function.  Composing with the inverse branches of
Lemma~\ref{lem:cusp-shadow}, which are jointly holomorphic, gives joint
holomorphy on $U_0\times W_0$.  On overlapping discs the local families
coincide by Lemma~\ref{lem:change-line}.
\end{proof}

\subsection{Abel coordinates on the tilted region}

\begin{lemma}[Koenigs times along orbits]\label{lem:cusp-abel}
Let $b<\eta$ be real, so $\varepsilon>0$, and put
$\Psi=A_1\log(u-u_1)+A_2\log(u-u_2)$, $A_2=1/\log\lambda_2$, and
$F_b=\Psi\circ f_b-\Psi-1$ as in Lemma~\ref{lem:F}.  Let $P$ be a path in the
$\chi$-plane such that for each $\chi\in P$ the forward orbit $f_\chi^k(\chi)$ stays in
$\mathcal G^{\rm c}$ and $\operatorname{Re}f_\chi^k(\chi)$ decreases by at least $\varepsilon/2$ per step.  Then, along
$P$, the branch of $\Rr^{-1}$ continued along $P$ satisfies
\[
 \Rr^{-1}=\Psi+\sum_{k\ge0}F_b\circ f^k+c_+,\qquad
 \sum_{k\ge0}|F_b(f^k(u))|\le C\varepsilon,
\]
with a constant $c_+$.  The same holds for $S^{-1}$ with backward orbits
($f_\chi^{-k}$, $\operatorname{Re}$ increasing by at least $\varepsilon/2$ per step) and
$\Psi-\sum_{k\ge1}F_b\circ f^{-k}+c_-$.
\end{lemma}

\begin{proof}
By Lemma~\ref{lem:F} and Lemma~\ref{lem:cusp-step}(1),
$|F_b(f^k(u))|\le MC\varepsilon^2\min(|\eee^{\chi_k}|,|\eee^{-\chi_k}|)\le C\varepsilon^2\eee^{-|\operatorname{Re}\chi_k|}$;
since $\operatorname{Re}\chi_k$ changes by at least $\varepsilon/2$ per step, the sum is at most
$C\varepsilon^2\cdot\sum_k\eee^{-|\operatorname{Re}\chi_0-k\varepsilon/2|}\le C\varepsilon$.  The orbit converges to $u_1$
because $\operatorname{Re}\chi_k\to-\infty$.  As in the proof of Theorem~\ref{thm:C}(1),
$\lim_n[\Psi(f^nu)-n]=\Psi(u)+\sum_{k\ge0}F_b(f^ku)$ exists and equals
$A_1\log\sigma_1(u)$ plus a constant, because $\sigma_1(f^nu)=\lambda_1^n\sigma_1(u)$ and
$A_1\log\lambda_1=1$; the constant depends only on the branches, which are
continued along $P$.  The backward case is identical with the reflected
inverse family of Theorem~\ref{thm:C}(2).
\end{proof}

\subsection{Kneser's initial region}
The exterior end of the corridor uses the following parameter-dependent
form of Kneser's construction.
\Needspace{22\baselineskip}
\begin{proposition}[Local base analyticity on the real exterior ray]
\label{prop:exterior-real-parameter}
For every real $b_0>\eta$, Kneser's real-normalised solution admits a
jointly holomorphic extension in $(b,z)$ to $\Delta\times W$, where
$\Delta$ is a complex disc about $b_0$ and $W$ is a complex
neighbourhood of $[0,1]$.  These local extensions agree on overlap,
so the classical exterior branch has a jointly holomorphic germ along
the entire real ray $(\eta,\infty)$.  In particular, every compact
subinterval of that ray has a common complex base neighbourhood and a
common complex height neighbourhood of $[0,1]$ on which this family is
jointly holomorphic.
\end{proposition}

\begin{proof}
Write $a=\log b$ and choose the principal fixed-point pair $L_+(b)$ and
$L_-(b)$, conjugate when $b>\eta$ and characterised there by
$0<a\operatorname{Im}L_+<\pi$.  They are simple fixed
points and hence holomorphic in $b$ near $b_0$.  Set
$u_b=(L_++L_-)/2$, $v_b=(L_+-L_-)/(2\ii)$, and
$\gamma_b(t)=u_b+\ii v_b t$ for $-1\le t\le1$.
At a real base, $u_b,v_b\in\mathbb R$, $v_b>0$, and
\[
 E_b(\gamma_b(t))=r_b\eee^{\ii\theta_b t},\qquad
 r_b=\eee^{a u_b}=|L_+|,\quad
 \theta_b=a v_b\in(0,\pi).
\]
Since $v_b=r_b\sin\theta_b$, the multipliers satisfy
$|E_b'(L_\pm)|=a r_b=\theta_b/\sin\theta_b>1$.
The straight chord $\gamma_b$ and this circular arc meet only at
$L_\pm$ and bound the initial region
$H_b=\{\operatorname{Re}w\ge u_b,\ |w|\le r_b\}\setminus\{L_+,L_-\}$.
This is Kneser's initial region as described in
\cite[\S4, Lem.~4]{TrappmannKouznetsov2011}.
The arcs depend holomorphically on $b$ through their real-analytic
parameter $t$.  They remain disjoint away from their endpoints after
shrinking a disc $\Delta$ about $b_0$; at the endpoints their tangents
remain distinct because the fixed-point multipliers are nonreal.
The bounded Jordan interiors vary continuously, so their union in
$\Delta\times\mathbb C$ is open; the ambient coordinates $(b,w)$ make
its projection a holomorphic family.  Identify thin analytic collars
of the two arcs by
$(b,w)\mapsto(b,E_b(w))$.  The map is injective on the collar after
shrinking $\Delta$, since it is injective on the compact chord at
$b_0$ and its derivative never vanishes.  This gives a Hausdorff
holomorphic family of cylinders with one end at each $L_\pm(b)$.

The ends can be filled holomorphically.  Indeed, parameter-dependent
Koenigs coordinates $\sigma_{\pm,b}$, normalised by
$\sigma_{\pm,b}(L_\pm)=0$ and $\sigma_{\pm,b}'(L_\pm)=1$,
linearise $E_b$ near the two repelling fixed points.  If
$\ell_{\pm,b}=\log E_b'(L_\pm(b))$ is
chosen continuously, then
$\exp(2\pi\ii\log\sigma_{\pm,b}/\ell_{\pm,b})$ (or its reciprocal)
is a puncture coordinate on the corresponding glued end.  Filling
both punctures yields a proper holomorphic submersion of compact
genus-zero surfaces over $\Delta$, just as in the collar construction
of Proposition~\ref{prop:weld-parameter}.

We also have a holomorphic marked section represented by $w=1$.
For real $b>\eta$, $E_b(x)>x$ for every real $x$: the difference is
positive at $x\to\pm\infty$ and has no real zero.  The adjacent real
intervals between the chord and its image are $(u_b,r_b)$, with
$r_b=E_b(u_b)>u_b$.  The forward endpoints $E_b^n(u_b)$ increase to
$+\infty$, since a finite limit would be a real fixed point.  For any
$x>0$ below $u_b$, forward iteration eventually enters $(u_b,r_b)$
or one of its boundary seams; for $x>r_b$, repeated real logarithms
eventually do the same.
Indeed, a monotone orbit staying on either side would have a finite
real fixed-point limit.  Thus these dynamical translates and their
seam collars cover the positive real axis.  In particular, $1$ is
in the interior of one such interval or on a seam.  A finite chain of
holomorphic dynamical charts represents it near $b_0$; at a seam the
analytic collar gives the same section.  The entire compact real
segment from $1$ to $b=E_b(1)$ is covered by finitely many of these
charts.

The two filled ends and this marked point are three disjoint
holomorphic sections.  Fischer--Grauert local triviality
\cite{FischerGrauert1965} and the fibrewise three-point M\"obius
normalisation give a jointly holomorphic coordinate $q_b$ on the
compactified cylinder, taking the upper end to $0$, the lower end to
$\infty$, and $w=1$ to $1$.  The upper puncture coordinate
$\exp(2\pi\ii\log\sigma_{+,b}/\ell_{+,b})$ shows that the logarithmic
lift $z=(2\pi\ii)^{-1}\log q_b$, with $z(1)=0$, gains $+1$ under the edge
identification by $E_b$.  On the real slice, reflection exchanges the
two ends and fixes the marked point, so $q_b(\bar w)=1/\overline{q_b(w)}$
in the real dynamical charts.  Thus $z$ is real there; its local
injectivity and $z(E_b(w))=z(w)+1$ make it strictly increasing along
the positive axis.  Inverting on the finite
chart chain covering the real segment from $1$ to $b$ gives a jointly
holomorphic superfunction on $\Delta\times W$, after shrinking both.

For real $b$, the lift restricted to $H_b$ is injective and its integer
translates cover the universal-cover plane.  Choosing any $d\in H_b$
and subtracting its value gives the initial-region criterion
of~\cite[Thm.~1]{TrappmannKouznetsov2011}; the injectivity and tiling
assertions for the Kneser generalisation are proved in
\cite[Thm.~5]{TrappmannKouznetsov2011}.  We normalise at $d$ here,
since $1$ need not itself belong to $H_b$.
Their Abel coordinates differ only by a constant on $H_b$, and
continuation along the real dynamical charts to $w=1$ makes that
constant zero.  Thus the inverse just
constructed is the classical real-normalised Kneser branch.
The marking fixes the local family uniquely, and the identity theorem
patches the parameter discs on overlaps.  A finite subcover of any
compact real base interval, followed by intersection of its finitely
many height neighbourhoods, gives the final uniformity assertion.
\end{proof}

\subsection{The two ends of the corridor}

\begin{proposition}[Exterior end]\label{prop:cusp-exterior}
For $0<\theta<\theta_1$ (real $b\in(\eta,\eta+r_1)$), $F_{\varepsilon(\theta)}$ is Kneser's
real-normalised solution on $W_0$.
\end{proposition}

\begin{proof}
The line $\ell=\ii\pi+\mathbb R$ ($\alpha=0$, $x_0=0$) is admissible for $\varepsilon(\theta)$, because
$-\operatorname{Im}\varepsilon(\theta)=\theta$.  Since $\chi\in\ii\pi+\mathbb R$ means $q<0$ and $u_2=\overline{u_1}$,
$\gamma=u(\ell)$ is the open segment between the conjugate fixed points, i.e.\ in
$w=\eee(1+u)$ the chord $\gamma_b$ of Proposition~\ref{prop:exterior-real-parameter}, and
$f(\gamma)$ is its image $E_b(\gamma_b)$.  Hence $H$ is the initial region $H_b$ and $X$
is the cylinder of that proposition, with the upper end at $L_1=L_+$.  There
the uniformisation takes the upper end to $0$ and is normalised so that
$z(1)=0$, with $z$ real and increasing along the positive axis.  For
conjugate $u_1,u_2$ the real axis is $\{|q|=1\}$, i.e.\ $\operatorname{Re}\chi=0$, and the
continued branch of the real base orbit moves up the imaginary $\chi$-axis
from $0$; it crosses the chord between $u^*_{n_0-1}$ and $u^*_{n_0}$, and
$u^*_{n_0}=E_b^{n_0}(1)$ has Kneser coordinate $n_0$.  So our lift is Kneser's,
and our inverse branches are the real logarithms along the orbit, as in
Kneser's normalisation.  By Lemma~\ref{lem:change-line} the admissible line
used in the family may be replaced by this one.
\end{proof}

\begin{proposition}[Interior end]\label{prop:cusp-interior}
There is $r_4>0$ such that $F_\varepsilon=\Kk^W_{b(\varepsilon)}$ on $W_0$ for real $\varepsilon\in(0,r_4)$.
Here $r_4$ is chosen so that $b(\varepsilon)\in(b_1,\eta)$ with the closeness conditions
of Theorem~\ref{thm:W2}, and so that $0.39h+C_{\rm im}<d-D$, where $C_{\rm im}$ is the
constant in~\eqref{eq:cusp-im} below and $d,D$ are as in Lemma~\ref{lem:modes}
and Theorem~\ref{thm:W2}.
\end{proposition}

\begin{proof}
Let $\varepsilon\in(0,r_4)$ be real, so $b<\eta$ is real, $\lambda_1\in(0,1)$ and $u_1<0<u_2$.
Take $\alpha=\tfrac\pi4$, $x_0=0$, so that
$\chi=\ii\pi+(t+\ii\varsigma)\eee^{\ii\pi/4}$ has
$\operatorname{Re}\chi=(t-\varsigma)/\sqrt2$ and $\operatorname{Im}\chi=\pi+(t+\varsigma)/\sqrt2$.  Put
$h=2\pi/|\log\lambda_1|$.  The real line of the $u$-plane appears as follows:
$u<u_1$ is $\{\operatorname{Im}\chi\in2\pi\mathbb Z,\ \operatorname{Re}\chi<0\}$, the segment
$(u_1,u_2)$ is $\{\operatorname{Im}\chi\in\pi+2\pi\mathbb Z\}$, and $u>u_2$ is
$\{\operatorname{Im}\chi\in2\pi\mathbb Z,\ \operatorname{Re}\chi>0\}$; the lower half-plane is
$\{0<\operatorname{Im}\chi<\pi\}$ modulo $2\pi$, because $q\mapsto u$ has negative determinant
$u_1-u_2$.  Since $f_\varepsilon$ is real and $\Delta$ is real at $u_1,u_2$,
Lemma~\ref{lem:cusp-step} gives $\operatorname{Im}\Delta=O(\varepsilon^2)$ on $\mathcal G^{\rm c}$ and
$\operatorname{Im}\Delta=O(\varepsilon^2|\eee^{\mp\chi}|)$ near the two ends.

\emph{Marked point.}  The base orbit is real, $u^*_n<u_1$, so $\chi^*_n$ is real and
moves to the left.  On the real axis $\varsigma=-(\operatorname{Re}\chi+\pi)/\sqrt2$, so the marked
point $p_c:=\chi^*_{n_0}$ lies in $[-\pi-2\varepsilon,-\pi]$, with $t\approx-\pi\sqrt2$.

\emph{The two charts.}  Let $U\subset X$ and $V\subset X$ be the images of
$M(\{t<\pi\sqrt2-1\})$ and $M(\{t>-\pi\sqrt2+1\})$.  They are open, and saturated
for the seam identification, since $M(t,0)\sim M(t,1)$ preserves $t$.  For
$\chi\in M(\{t<\pi\sqrt2-1\})$ the forward orbit $f_\chi^k(\chi)$ moves left by at least
$\varepsilon/2$ per step, and its imaginary part drifts by $O(\varepsilon)$ in total (the sum
of $O(\varepsilon^2)$ over $O(1/\varepsilon)$ steps in $|\operatorname{Re}\chi|\le4$ plus the geometric tail
at the end).  Such points with $\operatorname{Re}\chi\ge-\tfrac12$ have
$\operatorname{Im}\chi\in(\pi-1,2\pi-0.7)$, so these orbits stay in $\mathcal G^{\rm c}$.
Symmetrically, backward orbits of points of $M(\{t>-\pi\sqrt2+1\})$ move right,
stay in $\mathcal G^{\rm c}$ and converge to $u_2$.  So Lemma~\ref{lem:cusp-abel} applies
on both sets, and the attracting and repelling Koenigs coordinates $\sigma_1$ and
$\tau$ of Section~\ref{sec:transition} are defined and nonzero along these orbits.
Let $r_+$ be the branch of $\Rr^{-1}=\log(\sigma_1/\sigma_1(u^*))/\log\lambda_1$ on
$M(\{t<\pi\sqrt2-1\})$ continued from $p_c$, where $r_+(p_c)=n_0$ is real
($\sigma_1(u^*_{n_0})/\sigma_1(u^*)=\lambda_1^{n_0}>0$).  By Lemma~\ref{lem:cusp-abel},
$\partial_\chi r_+=A_1+O(1)$ on convex subsets at distance $\ge\tfrac12+\tfrac1{10}$ from the holes.  For
$\chi\in\ell$, $\sigma_1\circ f=\lambda_1\sigma_1$ gives $r_+(f_\chi(\chi))-r_+(\chi)\in1+\ii h\mathbb Z$ for the
continuation along $s\mapsto M(t,s)$, a segment of length $O(\varepsilon)$ on which $r_+$
changes by $O(1)$, far less than $h$.  Hence $r_+\circ f_\chi=r_++1$ on $\ell$, and $r_+$
descends to a holomorphic map $U\to\mathbb C/\mathbb Z$.

Let $(t_c,\varsigma_c)$ be the coordinates of $p_c$, so $0\le\varsigma_c<g(t_c)$, and
follow the curve $\varsigma=\varsigma_c\,g(t)/g(t_c)$, which lies in $\hat H$, in the direction
of increasing $t$.  Since $|g'|\le C\varepsilon^2$ and $g\ge\tfrac25\varepsilon$,
$\operatorname{Im}\chi=\pi+(t+\varsigma)/\sqrt2$ increases strictly along it, from $0$ at $p_c$.
Let $u_m\in(u_1,u_2)$ be the point where it equals $\pi$.  Let $r_-$ be the branch
of $S^{-1}$ on $M(\{t>-\pi\sqrt2+1\})$ that is real at $u_m$; it descends to $V$
in the same way.

\emph{Imaginary parts.}  In Lemma~\ref{lem:cusp-abel},
$\Psi=A_1\chi+(A_1+A_2)\log(u-u_2)$, with $A_1=-h/2\pi$,
$A_2=\tfrac h{2\pi}(1+O(\varepsilon))$ and $A_1+A_2\to\tfrac13$ (Lemma~\ref{lem:res}), and
$\arg(u-u_2)=\arg(u_1-u_2)-\arg(1-\eee^\chi)$.  On $U$ the continued
$\arg(1-\eee^\chi)$ is bounded, because $\operatorname{Im}\chi<2\pi-0.7$ wherever
$\operatorname{Re}\chi\ge-\tfrac12$ and $|\eee^\chi|\le\eee^{-1/2}$ elsewhere; on $U\cap V$ the range of
$\operatorname{Im}\chi$ is bounded.  Hence, with the branches fixed above,
\begin{equation}\label{eq:cusp-im}
 \bigl|\operatorname{Im}r_++\tfrac h{2\pi}\operatorname{Im}\chi\bigr|\le C_{\rm im}\ \text{on }U,\qquad
 \bigl|\operatorname{Im}r_-+\tfrac h{2\pi}(\operatorname{Im}\chi-\pi)\bigr|\le C_{\rm im}\ \text{on }U\cap V .
\end{equation}
On $V\setminus U$, $\operatorname{Im}\chi>2\pi-1/\sqrt2$.  There
$\arg(1-\eee^\chi)=\operatorname{Im}\chi-\pi+O(1)$ where $\operatorname{Re}\chi\ge0$, and
$A_1-(A_1+A_2)=-A_2$, so $\operatorname{Im}r_-\le-\min(|A_1|,A_2)(\operatorname{Im}\chi-\pi)+C$.  Along the
curve from $p_c$ to $u_m$, $\operatorname{Im}\chi$ increases from $0$ to $\pi$, so the path runs
through the lower half-plane: $\arg(u-u_1)$ increases by $\pi$ and
$\arg(u-u_2)$ returns to its initial value.  Hence
$\operatorname{Im}r_+(u_m)=A_1\pi+O(\varepsilon)=-h/2+O(\varepsilon)$.  By
Lemma~\ref{lem:realT}, $\sigma_1(u_m)/\sigma_1(u^*)<0$, so
$\operatorname{Im}r_+(u_m)\in-h/2+h\mathbb Z$; therefore $\operatorname{Im}r_+(u_m)=-h/2$ exactly.  This is
the seam branch of Proposition~\ref{prop:weld} and Definition~\ref{def:class}: by
item~(c) in Section~\ref{sec:proofC}, the path from the base point to the
segment passing above $u_1$ gives $+h/2$, and the one passing below gives
$-h/2$.  (The reflected branch $\overline{\Kk^W_b(\bar z)}$, denoted $\Kk^{W,\#}_b$ in Corollary~\ref{cor:conjugate-sewn-jump} below, would require the path above $u_1$,
i.e.\ a line of negative slope at real $\varepsilon$, which is not admissible.)

\emph{Transition.}  On $U\cap V$, $\Rr(r_+)=u=S(r_-)$, so $r_+-T(r_-)\in\ii h\mathbb Z$ is
constant, with $T$ the map of Section~\ref{sec:transition} (branch
$\operatorname{Im}t_0=h/2$).  At $u_m$, $r_-$ is real and $\operatorname{Im}T(r_-)=h/2$
(Lemma~\ref{lem:realT}), while $\operatorname{Im}r_+=-h/2$; hence $r_+=T\circ r_--\ii h$ on $U\cap V$.
By~\eqref{eq:cusp-im} and $|\operatorname{Im}\chi-\pi|\le\pi-1/\sqrt2$ on $U\cap V$, $r_-(U\cap V)$
lies in $\{|\operatorname{Im}w'|<0.39h+C_{\rm im}\}$, which is inside the strip
$\Sigma_q=\{|\operatorname{Im}w'|<d-D\}$ of the proof of Theorem~\ref{thm:char} by the choice of
$r_4$.  There $T$ is univalent, $T(\mathbb R)=\{\operatorname{Im}w=h/2\}$, and $T$ maps
$\{\operatorname{Im}w'\lessgtr0\}\cap\Sigma_q$ below and above that line.

\emph{The map to the sewn surface.}  Let $X^W$ be the surface of
Proposition~\ref{prop:weld}: $C_+=\{\operatorname{Im}w\ge-h/2\}/\mathbb Z$ and
$C_-=\{\operatorname{Im}w'\le0\}/\mathbb Z$ glued by $w=T(w')-\ii h$.  Define $\Theta:X\to X^W$
by $\Theta=r_+$ where $\operatorname{Im}r_+\ge-h/2$ and $\Theta=r_-$ where $\operatorname{Im}r_-\le0$.
By~\eqref{eq:cusp-im} and the estimate on $V\setminus U$, points of $U\setminus V$ have
$\operatorname{Im}\chi<1/\sqrt2$, hence $\operatorname{Im}r_+>-0.12h-C>-h/2$, and points of $V\setminus U$ have
$\operatorname{Im}r_-<-0.38h+C<0$.  On $U\cap V$, $\operatorname{Im}r_+\ge-h/2$ if and only if
$\operatorname{Im}r_-\ge0$, by the previous paragraph.  So every point of $X$ is covered, and
on the curve $\operatorname{Im}r_-=0$ both definitions give the same point of $X^W$.  Thus $\Theta$
is continuous, holomorphic off an analytic curve, hence holomorphic.  As a
point of $X$ tends to the $u_1$-end, $t\to-\infty$ and $\operatorname{Im}r_+\to+\infty$; as it tends
to the $u_2$-end, $\operatorname{Im}r_-\to-\infty$.  So $\Theta$ extends to a holomorphic map of
the compactified spheres with $\Theta^{-1}(\text{end})=\text{end}$.

It has local degree one at the $u_1$-end; in fact $r_+$ is injective on $U$.
Let $x,x'\in U$ be represented by $\chi,\chi'\in\hat H$ with $r_+(\chi')-r_+(\chi)=m\in\mathbb Z$.  By
Lemma~\ref{lem:cusp-abel}, $r_+=A_1\chi+\phi_+(\chi)$ where $\phi_+$ is bounded with
bounded derivative on the $\tfrac1{10}$-neighbourhood of $\hat H\cap U$, which lies in
$\mathcal G^{\rm c}$.  Equal
imaginary parts give $|\operatorname{Im}\chi-\operatorname{Im}\chi'|\le C\varepsilon$; since $\hat H$ is a strip of normal
width $\le\tfrac65\varepsilon$ along a line of slope $1$, this forces $|\chi-\chi'|\le C_1\varepsilon$, and then
$|m|\le|A_1||\chi-\chi'|(1+C\varepsilon)+C\le C_2$.  On the convex disc $N$ of radius
$(C_1+3C_2)\varepsilon$ about $\chi$, $\partial_\chi r_+=A_1(1+O(\varepsilon))$, so $r_+$ is injective on $N$; the
points $\chi'$ and $f_\chi^m(\chi)$ lie in $N$, and $r_+\circ f_\chi^m=r_++m$ gives
$r_+(f_\chi^m(\chi))=r_+(\chi')$, hence $\chi'=f_\chi^m(\chi)$.  Since $\hat H$ is a fundamental strip
for $f_\chi$, $x=x'$ in $X$.  Hence $\Theta$ has degree one and is biholomorphic.

\emph{Normalisation.}  $\Theta$ maps the marked point $p_c$ to the point of $C_+$
with attracting time $r_+(p_c)=n_0$.  In Proposition~\ref{prop:weld} the lift is
fixed by $\Kk^W(0)=\Rr(0)=1$ on the branch $w(0)=0$, i.e.\ attracting time $0$
at $z=0$, hence time $n_0$ at $z=n_0$.  So the lifts agree, and
$\tilde F(z+n_0)=\Rr(r_+)=\Kk^W_b(z+n_0)$ near $z=0$.  Near $z=0$ the points
$\Kk^W_b(z+k)$, $0\le k\le n_0$, stay near $u^*_k$, where the inverse branches of
Definition~\ref{def:cusp-solution} invert $f_\varepsilon$ (Lemma~\ref{lem:cusp-param}(3));
hence $F_\varepsilon=\Kk^W_b$ near $z=0$, and by the identity theorem on $W_0$.
\end{proof}

\subsection{The continuation theorem}

\begin{theorem}[Continuation around the cusp]\label{thm:cusp}
There are $r_5>0$ and an open set $\Omega\supset D^+\cup(\eta-r_5,\eta)\cup(\eta,\eta+r_5)$,
where $D^+=\{b:|b-\eta|<r_5,\ \operatorname{Im}b>0\}$, with the following properties.
\begin{enumerate}
\item Kneser's real-normalised solution $\kappa_b$, $b\in(\eta,\eta+r_5)$, continues
  analytically in the base to a family on $\Omega$ that is jointly holomorphic on
  $\Omega\times W_0$, $W_0=\{z:\operatorname{dist}(z,[0,1])<\tfrac14\}$.  In particular it
  continues across the arc of the Shell--Thron boundary contained in $D^+$,
  whatever the arithmetic of the neutral multiplier there.
\item On $(\eta-r_5,\eta)$ the continued family equals $\Kk^W_b$ on $W_0$.
\item Let $N_W$ be the neighbourhood of $(1,\eta)$ of
  Proposition~\ref{prop:weld-parameter}.  Continuing the family from a
  neighbourhood of $(\eta-r_5,\eta)$ inside $N_W$, on a $b$-dependent
  neighbourhood of $z=0$ contained in the regular domain of $\Kk^W_b$, one
  obtains $\Kk^W$.  Hence for every $b\in(1,\eta)$, if $\kappa_{b+\ii y}$ is continued first
  through $D^+$ and then inside $N_W\cap\{\operatorname{Im}b>0\}$, the limit
  $\lim_{y\downarrow0}\kappa_{b+\ii y}$ exists near $z=0$ and equals $\Kk^W_b$; by the
  functional equation and continuation in $z$ it equals $\Kk^W_b$ on its whole
  regular domain.
\end{enumerate}
\end{theorem}

\begin{proof}
(1) By Lemma~\ref{lem:cusp-param}, $\varepsilon\mapsto b(\varepsilon)$ maps $\mathcal S^{\rm c\prime}_{r_3}$
biholomorphically onto an open set $\Omega$ that contains $D^+$ and the two real
segments, for $r_5$ small.  By Proposition~\ref{prop:cusp-holo},
$(b,z)\mapsto F_{\varepsilon(b)}(z)$ is jointly holomorphic on $\Omega\times W_0$.  By
Proposition~\ref{prop:cusp-exterior} it equals $\kappa_b$ on $(\eta,\eta+r_5)$.  Since
$\Omega$ is connected, it is the analytic continuation of $\kappa$ along every path
in $\Omega$.  The Shell--Thron boundary near $\eta$ consists of
$b=\exp(\lambda\eee^{-\lambda})$ with $|\lambda|=1$.  Its arc in the upper half-plane is
$\varepsilon=1-\eee^{\ii\phi}=-\ii\phi+O(\phi^2)$ with small $\phi>0$, since then
$b-\eta$ is a positive multiple of $-\varepsilon^2=\phi^2+\ii\phi^3+O(\phi^4)$; it lies in
$\mathcal S^{\rm c\prime}_{r_3}$.

(2) is Proposition~\ref{prop:cusp-interior}.

(3) By Proposition~\ref{prop:weld-parameter}, $\Kk^W$ is jointly holomorphic on
the set of $(b,z)$ with $b\in N_W$ and $z$ in a neighbourhood of $0$ that
depends continuously on $b$; the upper chart contains $w=0$, where $\Rr$ is
regular.  Shrinking the discs $\Delta_{b_0}$ of that proposition, we may assume
$\Delta_{b_0}\subset\{|b-b_0|<(\eta-b_0)/2\}$; then $N_W\cap\{\operatorname{Im}b>0\}$ is a union
of upper half-discs centred on $(1,\eta)$, which is connected since
consecutive ones overlap.  Each $\Delta_{b_0}\cap D^+$ is convex, and if it is
nonempty its closure meets $(\eta-r_5,\eta)$.  On a neighbourhood of
$(\eta-r_5,\eta)$ times a small disc about $z=0$, the two families agree by (2)
and the identity theorem; hence they agree on every $\Delta_{b_0}\cap D^+$, and the
continuation from $D^+$ into $N_W\cap\{\operatorname{Im}b>0\}$ is unambiguous and equals
$\Kk^W$.  At every $b\in(1,\eta)$ the boundary value is therefore $\Kk^W_b$ near
$z=0$.  Both sides satisfy
$F(z+1)=b^{F(z)}$, so they agree on the whole regular domain.
\end{proof}

\begin{corollary}[Separation limits for the continued Kneser family]\label{cor:conjecture}
For the continuation of Kneser's solution in Theorem~\ref{thm:cusp},
$\Kk_b=\lim_{y\downarrow0}\kappa_{b+\ii y}$ exists for $b\in(1,\eta)$ and equals $\Kk^W_b$.  Its
separation coefficients satisfy $c_n/\Lambda^n\to B_n\eee^{2\pi\ii na}$ for every $n\ge1$
as $b\to\eta^-$.  In particular $|c_1|/\Lambda\to|B_1|=0.0890584\ldots$, with $B_1\ne0$.
This confirms the limit suggested by the numerics of Section~\ref{sec:invariant}:
the first separation mode of the continued Kneser family, scaled by $\Lambda$,
tends to the first coefficient of the inverse horn map of $\eee^u-1$.
\end{corollary}

\begin{proof}
Theorem~\ref{thm:cusp} gives $\Kk_b=\Kk^W_b$; apply Corollary~\ref{cor:limits}
and Proposition~\ref{prop:horn-nonzero}.
\end{proof}

\begin{remark}[Lower half-plane, and what is not claimed]\label{rem:cusp}
(1) Complex conjugation of the whole construction uses the sector
$\{\bar\varepsilon:\varepsilon\in\mathcal S^{\rm c}_r\}$, the lines $\overline{\ell_{\alpha,x_0}}$ through
$-\ii\pi\equiv\ii\pi\pmod{2\pi\ii}$ with slope $-\alpha$, and the conjugate admissibility
condition $+\operatorname{Im}(\varepsilon\bar\omega)\ge\tfrac12|\varepsilon|$.  It shows that continuation of
Kneser's solution through the lower half-disc yields the reflected branch
$\Kk^{W,\#}_b$ of Corollary~\ref{cor:conjugate-sewn-jump} on $(\eta-r_5,\eta)$.  By Corollary~\ref{cor:conjugate-sewn-jump} below, the two continuations differ
by at least $2c\Lambda$ at the height $x_*$ of
Corollary~\ref{cor:cusp-real-obstruction}; this is the monodromy of Kneser's family
around $\eta$.

(2) The proof does not use any arithmetic property of the multiplier on the
Shell--Thron boundary.  The ends of the tilted region are punctures for
Siegel, Cremer and parabolic multipliers alike, because the only inputs are
the step estimate $A_1\Delta=1+O(\varepsilon)$ and the transversality
$\operatorname{Im}(\operatorname{Log}\lambda_1\bar\omega)>0$.  In particular no
Brjuno condition, Siegel-disc estimate or boundary-value (Painlev\'e-type)
argument is needed near $\eta$.

(3) The theorem asserts holomorphy in $(b,z)$ for heights near $[0,1]$, and
by the functional equation on the regular domain reached from there.  No
claim is made about the behaviour of the continued solution as
$\operatorname{Im}z\to\pm\infty$ at parameters with a neutral multiplier; the
construction only uses that each end of the tilted region is a puncture of
the quotient.

(4) The theorem concerns continuation along the stated paths.  The width of
$N_W$ is not quantified; the restriction to paths near the real axis is
removed in Theorem~\ref{thm:global}, and the relation to Paulsen's family is
discussed in Remark~\ref{rem:global-paulsen}.
\end{remark}

\subsection{The whole upper Shell--Thron region}\label{sec:global}
Theorem~\ref{thm:cusp} continues Kneser's solution into a corridor about
the cusp.  Here we show that the continued family extends, single-valued,
to the whole upper half of the Shell--Thron region.  The extension is a
quasiconformal transport of one tilted initial region by a holomorphic
motion that conjugates the dynamics; the new point is a comparison lemma
identifying the transported quotient with the tilted one near the base
parameter.  In this subsection $\lambda$ denotes the multiplier of the
attracting fixed point, so $\lambda=\lambda_1=1-\varepsilon$ in the notation above, and
\[
 \mathbb H_\lambda=\{\lambda:|\lambda|<1,\ \operatorname{Im}\lambda>0\},\qquad
 b(\lambda)=\exp(\lambda\eee^{-\lambda}),\qquad E_\lambda=E_{b(\lambda)} .
\]

\begin{lemma}[Global parameter]\label{lem:global-param}
$b$ is injective on the unit disc, $E_\lambda$ has the fixed point
$L_1(\lambda)=\eee^\lambda$ with multiplier $\lambda$, and $b(\{|\lambda|<1\})$ is the
set of bases for which $E_b$ has an attracting fixed point.  The set
$\mathcal U^+=b(\mathbb H_\lambda)$ lies in $\{\operatorname{Im}b>0\}$ and is simply
connected.
\end{lemma}

\begin{proof}
For $g(\lambda)=\lambda\eee^{-\lambda}$ one has $\operatorname{Re}(\lambda g'/g)=\operatorname{Re}(1-\lambda)>0$ on the
unit disc, so $g$ is starlike and injective.  Since $|g|<\eee<\pi$, equality
$\eee^{g(\lambda_1)}=\eee^{g(\lambda_2)}$ forces $g(\lambda_1)=g(\lambda_2)$; so $b$ is injective.  If
$L$ is a fixed point of $E_b$ with multiplier $\lambda=L\log b$ of modulus $<1$,
then $L=\eee^\lambda$ and $\log b=g(\lambda)$ (Lemma~\ref{lem:param}).  By symmetry and
injectivity $g$ maps the upper half-disc into $\{\operatorname{Im}>0\}$ (check
$\lambda=\ii y$, $y>0$ small); since $|\operatorname{Im}g|<\pi$, $\operatorname{Im}b>0$ there.  An injective
image of a simply connected domain is simply connected.
\end{proof}

Fix a base parameter $\lambda_0=1-\varepsilon_0$ with $\varepsilon_0\in\mathcal S^{\rm c\prime}_{r_3}$,
$\operatorname{Im}\varepsilon_0<0$ and $|\lambda_0|<1$, such that $E_{\lambda_0}^n(1)\ne L_1(\lambda_0)$ for all $n$.  Each of these
equations is holomorphic in $\lambda$ on the unit disc and not identically
satisfied (for real $\lambda\in(0,1)$ the orbit of $1$ increases strictly to $L_1$), so
this excludes a countable set.  Only $\sigma_0(p_0)\ne0$ for the marked point
$p_0$ below is used.  We take $r_3\le1/240$ in Proposition~\ref{prop:cusp-holo}, which
only shrinks the corridor.  Write $E_0=E_{\lambda_0}$, let $B_0$ be its attracting basin and
$\sigma_0$ its Koenigs coordinate at $L_1(\lambda_0)$, normalised by $\sigma_0'(L_1)=1$ and
extended to $B_0$ by $\sigma_0=\lambda_0^{-n}\sigma_0\circ E_0^n$.  Since $E_0$ has no critical
points, $\sigma_0'\ne0$ on $B_0$.

\begin{proposition}[Transport of the dynamics]\label{prop:global-transport}
For $\lambda\in\mathbb H_\lambda$ put $s=\operatorname{Log}\lambda$, $s_0=\operatorname{Log}\lambda_0$ (so $\operatorname{Re}s,\operatorname{Re}s_0<0$),
\[
 \mathsf b(\lambda)=\frac{s-s_0}{2\operatorname{Re}s_0},\qquad
 \mathsf c(\lambda)=\frac{\mathsf b}{1+\mathsf b}=\frac{s-s_0}{s+\overline{s_0}},\qquad
 \phi_\lambda(\zeta)=\zeta|\zeta|^{2\mathsf b(\lambda)} .
\]
Then $|\mathsf c(\lambda)|<1$ on $\mathbb H_\lambda$.  Let $\nu_\lambda=\mathsf c(\lambda)\,\frac{\sigma_0}{\overline{\sigma_0}}\,\frac{\overline{\sigma_0'}}{\sigma_0'}$ on
$B_0\setminus\sigma_0^{-1}(0)$ and $\nu_\lambda=0$ elsewhere, and let $h_\lambda$ be the
quasiconformal homeomorphism of $\hat{\mathbb C}$ with Beltrami coefficient $\nu_\lambda$
fixing $0,1,\infty$.  Then:
\begin{enumerate}
\item $\lambda\mapsto h_\lambda(w)$ is holomorphic for each $w$, $h_{\lambda_0}=\mathrm{id}$, and
  $h_\lambda\to\mathrm{id}$ locally uniformly on $\mathbb C$ as $\lambda\to\lambda_0$;
\item $h_\lambda\circ E_0=E_\lambda\circ h_\lambda$, hence $h_\lambda(E_0^n(1))=E_\lambda^n(1)$ for all $n\ge0$;
\item $h_\lambda(B_0)$ is the basin of $L_1(\lambda)=h_\lambda(L_1(\lambda_0))$, and
  $\sigma_\lambda\circ h_\lambda=\mathsf k_\lambda\,\phi_\lambda\circ\sigma_0$ on $B_0$, with $\mathsf k_\lambda\ne0$ holomorphic in $\lambda$ and $\mathsf k_{\lambda_0}=1$;
\item for $\lambda$ near $\lambda_0$, $h_\lambda$ maps the repelling fixed point $L_2(\lambda_0)$ of
  Section~\ref{sec:cusp} to $L_2(\lambda)$.
\end{enumerate}
\end{proposition}

\begin{proof}
$|s+\overline{s_0}|^2-|s-s_0|^2=4\operatorname{Re}s\operatorname{Re}s_0>0$, so $|\mathsf c|<1$, uniformly on compact
sets.  In $L=\log\zeta$, $\phi_\lambda$ is the real-linear map $L\mapsto L+2\mathsf b\operatorname{Re}L$, with
determinant $\operatorname{Re}s/\operatorname{Re}s_0>0$; it fixes $2\pi\ii$ and sends $s_0$ to $s$.  So $\phi_\lambda$ is a
quasiconformal homeomorphism with $\phi_\lambda(\lambda_0\zeta)=\lambda\phi_\lambda(\zeta)$ and Beltrami
coefficient $\mathsf c(\lambda)\zeta/\bar\zeta$; $\nu_\lambda$ is its pull-back by $\sigma_0$, hence
$E_0$-invariant (the basin is completely invariant and $\sigma_0\circ E_0=\lambda_0\sigma_0$).
Item (1) is the Ahlfors--Bers theorem~\cite[Thm.~11]{AhlforsBers1960}, since
$\nu_\lambda$ depends holomorphically on $\lambda$ through the scalar $\mathsf c(\lambda)$, and
continuity of normalised solutions in $\nu$.  The map $h_\lambda\circ E_0$ has Beltrami coefficient $E_0^*\nu_\lambda=\nu_\lambda$, so
$G_\lambda=h_\lambda\circ E_0\circ h_\lambda^{-1}$ has vanishing dilatation and is holomorphic on $\mathbb C$
by Weyl's lemma; it is a covering of $\mathbb C^*$ (it omits $h_\lambda(0)=0$), so
$G_\lambda(w)=\eee^{A w+D}$, and $G_\lambda(0)=h_\lambda(1)=1$ gives $D\in2\pi\ii\mathbb Z$, which we drop.  Near $L_1(\lambda_0)$, $h_\lambda\circ\sigma_0^{-1}\circ\phi_\lambda^{-1}$ is conformal, so $G_\lambda$ is
conformally conjugate to $\zeta\mapsto\lambda\zeta$ at $h_\lambda(L_1(\lambda_0))$: that point is a
fixed point of $\eee^{Aw}$ with multiplier $\lambda$, which forces it to be $\eee^\lambda$ and
$A=\lambda\eee^{-\lambda}$.  So $G_\lambda=E_\lambda$, proving (2).  For (3), let $V=\sigma_0^{-1}(D_\varrho)$ be a small forward-invariant disc about
$L_1(\lambda_0)$.  The map $\psi=\sigma_\lambda\circ h_\lambda\circ\sigma_0^{-1}\circ\phi_\lambda^{-1}$ is defined on
$\phi_\lambda(D_\varrho)=D_{\varrho^{\operatorname{Re}s/\operatorname{Re}s_0}}$, conformal off $0$ (equal Beltrami
coefficients) and continuous at $0$, hence conformal; it fixes $0$ and
satisfies $\psi(\lambda\zeta)=\lambda\psi(\zeta)$.  Comparing Taylor coefficients ($\lambda^{n-1}\ne1$
for $n\ge2$) gives $\psi(\zeta)=\mathsf k_\lambda\zeta$.  The identity extends to $B_0$ by
$\sigma_\lambda=\lambda^{-n}\sigma_\lambda\circ E_\lambda^n$, and
$\mathsf k_\lambda=\sigma_\lambda(h_\lambda(w))/\phi_\lambda(\sigma_0(w))$ is holomorphic in $\lambda$, because
$\phi_\lambda(\zeta)=\zeta\exp(\mathsf b(\lambda)\log|\zeta|^2)$ is.  In particular, on $B_0$ near any point
where $\sigma_\lambda$ is locally injective,
$h_\lambda=\sigma_\lambda^{-1}\circ(\mathsf k_\lambda\phi_\lambda\circ\sigma_0)$ is jointly real-analytic in the space
variable, holomorphic in $\lambda$, with nonvanishing Jacobian.  For (4),
$h_\lambda(L_2(\lambda_0))$ is a fixed point of $E_\lambda$ depending continuously on $\lambda$ and equal
to $L_2(\lambda_0)$ at $\lambda_0$; fixed points are isolated.
\end{proof}

Let $H_0\subset\mathbb C$ be the tilted region of Lemma~\ref{lem:crescent} for $E_0$
and an admissible line $\ell$, written in $w=\eee(1+u)$, and $p_0=E_0^{n_0}(1)$ its
marked point; choose $\ell$ with $p_0$ in the interior.  For $\lambda\in\mathbb H_\lambda$ let
$X^{\rm tr}_\lambda$ be $h_\lambda(H_0)$ with its edges $h_\lambda(\gamma)$ and
$h_\lambda(E_0(\gamma))=E_\lambda(h_\lambda(\gamma))$ identified by $E_\lambda$, and $p_\lambda=h_\lambda(p_0)=E_\lambda^{n_0}(1)$.

\begin{proposition}[The transported family]\label{prop:global-family}
$X^{\rm tr}_\lambda$ is biholomorphic to $\mathbb C^*$, its ends at $L_{1,2}$.  Let $z$ be
the lift of its uniformisation with the $L_1$-end at $+\ii\infty$ and $z=n_0$ at
$p_\lambda$, and let $F^{\rm tr}_\lambda$ be the germ at $z=0$ given by the local inverse of
$E_\lambda^{n_0}$ at $p_\lambda$ with value $1$, applied to the inverse of $z$ at $z+n_0$.  Then
$F^{\rm tr}_\lambda(0)=1$, the lifted inverse uniformisation $\tilde F$ satisfies
$\tilde F(z+1)=E_\lambda(\tilde F(z))$, and $(\lambda,z)\mapsto F^{\rm tr}_\lambda(z)$ is jointly holomorphic near
$\mathbb H_\lambda\times\{0\}$ and single-valued there.
\end{proposition}

\begin{proof}
$h_\lambda$ induces a quasiconformal homeomorphism $J_\lambda:X_0\to X^{\rm tr}_\lambda$,
because it conjugates the edge identifications; each end of $X^{\rm tr}_\lambda$ is
the image of a punctured-disc end of $X_0\cong\mathbb C^*$
(Lemma~\ref{lem:crescent-quotient}), hence a punctured disc.  Near the seam
the complex structure is given by local inverses of $E_\lambda$, which is
injective there because $E_0$ is.  The Beltrami coefficient of $J_\lambda$ is the
descent of $\nu_\lambda$; transported to $\hat{\mathbb C}$ by a uniformisation $Z_0$ of $X_0$
normalised at the ends and $p_0$, it equals $\mathsf c(\lambda)\nu$ for a fixed measurable
$\nu$ with $|\nu|\le1$.  Let $\mathcal W_\lambda$ be the normalised solution; by Ahlfors--Bers
it is holomorphic in $\lambda$ at each point, and $Z_\lambda=\mathcal W_\lambda\circ Z_0\circ J_\lambda^{-1}$ is the
uniformisation of $X^{\rm tr}_\lambda$ normalised at the ends and $p_\lambda$.  Near $p_0$ the
coefficient $\nu_\lambda$ is real-analytic in $w$, because $\sigma_0(p_0)\ne0$.  The
$\partial_{\bar\lambda}$-cancellation of Proposition~\ref{prop:cusp-holo}, with $h_\lambda$ in place
of $\mathcal J_\varepsilon$, shows that $(\lambda,w)\mapsto Z_\lambda(w)$ is jointly holomorphic near
$(\lambda,p_\lambda)$.  The inverse function theorem and the local inverse of $E_\lambda^{n_0}$,
which is determined by the value $1$ at $p_\lambda$ and depends holomorphically on
$\lambda$, give the germ.  The $C^1$-dependence needed in the cancellation is
explicit here: near $p_0$, $h_\lambda$ is given by the formula at the end of the
proof of Proposition~\ref{prop:global-transport}, and $\mathcal W_\lambda$ is $C^1$ in
$(\lambda,\xi)$ by interior Schauder estimates~\cite[Ch.~V]{AhlforsLectures}.  All
objects are defined by a single formula for every $\lambda\in\mathbb H_\lambda$, so the family
is single-valued.
\end{proof}

\begin{lemma}[The lifted motion near the ends]\label{lem:global-ends}
For $\lambda$ near $\lambda_0$ let $\chi_\lambda=\log\frac{u-u_1(\lambda)}{u-u_2(\lambda)}$ be the coordinate of
Section~\ref{sec:cusp} for $E_\lambda$, and let $\tilde h_\lambda$ be the lift of $h_\lambda$ to these
coordinates (defined off the holes, commuting with $\chi\mapsto\chi+2\pi\ii$, continuous
in $\lambda$ and equal to the identity at $\lambda_0$).  The lift satisfies $\tilde h_\lambda\circ f_{\chi,\lambda_0}=f_{\chi,\lambda}\circ\tilde h_\lambda$ exactly (no deck shift),
and $\tilde h_\lambda\to\mathrm{id}$ uniformly on compact subsets of $\mathcal G^{\rm c}$.  Moreover,
for every $\delta>0$ there are $T$ and a neighbourhood $U_\delta$ of $\lambda_0$ such that
for $\lambda\in U_\delta$ there are real-linear maps $A^\pm_\lambda$ and constants $a^\pm_\lambda$ with
$\|A^\pm_\lambda-I\|+|a^\pm_\lambda|\le\delta$ and
\[
 |\tilde h_\lambda(\chi)-A^+_\lambda\chi-a^+_\lambda|\le\delta\quad(\operatorname{Re}\chi\le-T),\qquad
 |\tilde h_\lambda(\chi)-A^-_\lambda\chi-a^-_\lambda|\le\delta\quad(\operatorname{Re}\chi\ge T).
\]
The same holds for $\tilde h_\lambda^{-1}$.
\end{lemma}

\begin{proof}
This is a statement for $\lambda$ near $\lambda_0$, where $u_{1,2}(\lambda)$ and $\chi_\lambda$ are defined.
$h_\lambda$ maps $u_1(\lambda_0),u_2(\lambda_0),\infty$ to $u_1(\lambda),u_2(\lambda),\infty$
(Proposition~\ref{prop:global-transport}(3),(4)), so in the cross-ratio
$q=(u-u_1)/(u-u_2)$ it is an orientation-preserving homeomorphism of
$\mathbb C^*\setminus\{1\}$ fixing $0,\infty$.  Such a map acts trivially on the winding
number about $0$, so it lifts to $\mathbb C\setminus2\pi\ii\mathbb Z$ commuting with $2\pi\ii$; joint
continuity of $(\lambda,w)\mapsto h_\lambda(w)$ (the $\lambda$-lemma) lets us choose the lift
continuously with $\tilde h_{\lambda_0}=\mathrm{id}$, and gives local uniform convergence.  The
lifted conjugacy $\tilde h_\lambda\circ f_{\chi,\lambda_0}=f_{\chi,\lambda}\circ\tilde h_\lambda+2\pi\ii k$ holds with an
integer $k$ that is continuous in $(\lambda,\chi)$ and $0$ at $\lambda_0$, hence $k=0$.

Attracting end.  In $L=\log\sigma_0$ and $L'=\log\sigma_\lambda$,
Proposition~\ref{prop:global-transport}(3) says that $h_\lambda$ is the affine map
$L'=L+2\mathsf b\operatorname{Re}L+\log \mathsf k_\lambda$.  Near $u_1$, $\sigma_\lambda$ and $q$ are both local coordinates
vanishing at $u_1$, so $\chi_\lambda=L'+\psi_\lambda(\eee^{L'})$ with $\psi_\lambda$ holomorphic near $0$ and
continuous in $\lambda$; likewise for $\lambda_0$.  Composing, with $A^+_\lambda=I+2\mathsf b\operatorname{Re}$ and
$a^+_\lambda=\log\mathsf k_\lambda+\psi_\lambda(0)-A^+_\lambda\psi_{\lambda_0}(0)$, the remainder
$\tilde h_\lambda(\chi)-A^+_\lambda\chi-a^+_\lambda$ equals
$[\psi_\lambda(\eee^{L'})-\psi_\lambda(0)]-A^+_\lambda[\psi_{\lambda_0}(\eee^{L})-\psi_{\lambda_0}(0)]$.  Since
$|\eee^{L'}|\le C|\eee^{L}|^{\operatorname{Re}s/\operatorname{Re}s_0}$ and $\operatorname{Re}s/\operatorname{Re}s_0\ge\tfrac12$ near $\lambda_0$, it is
at most $C\eee^{-T/2}$ on $\operatorname{Re}\chi\le-T$, uniformly for $\lambda$ near $\lambda_0$.  Choose $T$ with
$C\eee^{-T/2}\le\delta$, and then $U_\delta$ so that $\|A^+_\lambda-I\|+|a^+_\lambda|\le\delta$.

Repelling end.  Let $M=\log\tau_0$ and $M'=\log\tau_\lambda$ be logarithms of the
repelling Koenigs coordinates at $L_2$, in which $E_0,E_\lambda$ are the
translations by $m_0=\log\mu_0$, $m_\lambda=\log\mu_\lambda$, $\operatorname{Re}m>0$.  For $\lambda$ near $\lambda_0$, $h_\lambda$
maps a linearising disc of $E_0$ into one of $E_\lambda$, and its lift $G_\lambda$ to
$\{\operatorname{Re}M<-R\}$ satisfies $G_\lambda(M-m_0)=G_\lambda(M)-m_\lambda$ and
$G_\lambda(M+2\pi\ii)=G_\lambda(M)+2\pi\ii$ (the first because $h_\lambda$ conjugates the local
inverse branches at $L_2$, with the integer fixed by continuity from $\lambda_0$; the
second by degree).  Let
$A'_\lambda$ be the real-linear map with $A'_\lambda m_0=m_\lambda$, $A'_\lambda(2\pi\ii)=2\pi\ii$.  Then
$D_\lambda=G_\lambda-A'_\lambda$ is invariant under $M\mapsto M-m_0$ and $M\mapsto M\pm2\pi\ii$, which
preserve $\{\operatorname{Re}M<-R\}$.  Every point of that half-plane is carried by these
maps into the compact set $\{-R-\operatorname{Re}m_0\le\operatorname{Re}M\le-R,\ 0\le\operatorname{Im}M\le2\pi\}$, so
$\sup|D_\lambda|$ is its supremum there, which tends to $0$ as $\lambda\to\lambda_0$.  As at the
other end, $\chi=-M+\mathrm{const}+\psi^-(\eee^{M})$ with $\psi^-$ holomorphic and vanishing
at $0$, for both parameters, and the same two-step choice of $T$ and $U_\delta$
gives the second estimate.  The inverse is treated in the same way.
\end{proof}

\begin{lemma}[Comparison]\label{lem:global-compare}
There is a neighbourhood $U$ of $\lambda_0$ such that for $\lambda\in U$ the transported
quotient $X^{\rm tr}_\lambda$ and the tilted quotient $X^\ell_{1-\lambda}$ of
Section~\ref{sec:cusp} (same admissible line $\ell$) are biholomorphic by a
map that preserves the two ends and sends $p_\lambda$ to the marked point.
Consequently $F^{\rm tr}_\lambda=F_{1-\lambda}$ near $z=0$ for $\lambda\in U$.
\end{lemma}

\begin{proof}
Let $\ell$ be the admissible line fixed in the proof of
Proposition~\ref{prop:cusp-holo} for the disc about $\varepsilon_0$; it is the line used for
$H_0$.  Take $U$ inside that disc, so that $\ell$ is admissible, $n_0$ is constant,
and the marked point of $X^\ell_{1-\lambda}$ is interior.  Take also $U\subset U_\delta$ of
Lemma~\ref{lem:global-ends} with $\delta=1/2000$, and $T'\ge(T+1)/0.6$, so that $|t|\ge T'$
implies $|\operatorname{Re}\chi|\ge T$ on the strips; for fixed $\lambda_0$ this is a genuine
neighbourhood, although its size may be small compared with $\operatorname{Re}\varepsilon_0$.  Work in the $\chi$-plane of $E_\lambda$, where $E_\lambda$ is $f_\chi=\mathrm{id}+\Delta$ and the
tilted strip is $\hat H_\lambda$ between $\ell$ and $f_\chi(\ell)$, and put
$\hat T_\lambda=\tilde h_\lambda(\hat H_0)$, the transported strip, bounded by $\Gamma=\tilde h_\lambda(\ell)$
and $f_\chi(\Gamma)$.  Let $\varsigma,t$ be the coordinates of Section~\ref{sec:cusp}
relative to $\ell$, and $\mathsf W_\rho=\{|\varsigma|\le\rho(1+|t|)\}$.

\emph{Step 1: $\hat T_\lambda\subset \mathsf W_{1/100}\cap\mathcal G^{\rm c}$ for $\lambda\in U$, after shrinking $U$.}
$\hat H_0$ lies in $\{|\varsigma|\le\tfrac65|\varepsilon_0|\}$.  On $\{|t|\le T'\}$ this is compact and
$\tilde h_\lambda\to\mathrm{id}$ uniformly.  For $|t|\ge T'$ (so $|\operatorname{Re}\chi|\ge T$),
Lemma~\ref{lem:global-ends} gives
$|\tilde h_\lambda(\chi)-\chi|\le\|A^\pm_\lambda-I\|\,|\chi|+|a^\pm_\lambda|+\delta$, and
$|\chi|<|t|+4\le5(1+|t|)$ on $\hat H_0$; with $\delta=1/2000$ the displacement is at most
$6\delta(1+|t|)=0.003(1+|t|)$ (also on $|t|\le T'$, after shrinking $U$).  The parameter
$t$ of the image point differs from that of the source by at most this amount.
Since $|\varepsilon|\le1.025|\varepsilon_0|\le0.005$ on $U$ (using $|\varepsilon_0|\le1/240$), every point of $\hat T_\lambda$ has
$|\varsigma|\le0.005+\tfrac{0.003}{0.997}(1+|t|)<0.0082(1+|t|)\le\tfrac1{100}(1+|t|)$, and its distance to the
holes is at least
$\max(1.48,0.6|t|-0.5)-0.003(1+|t|)\ge\tfrac1{20}(1+|t|)+\tfrac12+0.2$, so that the orbit
discs of Step~2, of radius $(1+|t|)/20$, stay in $\mathcal G^{\rm c}$.  Symmetrically,
$\tilde h_\lambda^{-1}(\hat H_\lambda)\subset \mathsf W_{1/100}$ in the $\chi$-plane of $E_0$.

\emph{Step 2: orbit correspondence.}  Let $P\in\hat T_\lambda$ with parameter $t$.  As in
the proof of Lemma~\ref{lem:change-line}, each application of $f_\chi^{\pm1}$ changes
$\varsigma$ by at least $\tfrac25|\varepsilon|$ in the same direction and moves the point by at most
$2|\varepsilon|$, and $|\varsigma(P)|\le\tfrac1{100}(1+|t|)$; so the orbit of $P$ reaches $\hat H_\lambda$
after at most $(1+|t|)/(40|\varepsilon|)$ steps, forward or backward, inside a disc about
$P$ contained in $\mathcal G^{\rm c}$, on which $f_\chi$ is injective.  Exactly one iterate
$f_\chi^{n(P)}(P)$ lies in $\hat H_\lambda\setminus f_\chi(\ell)$, by the monotonicity of $\varsigma$ and the
isotopy argument of Lemma~\ref{lem:cusp-base}(2).  (The edges of $\hat T_\lambda$ are $\Gamma$
and $f_\chi(\Gamma)$ by the exact lifted conjugacy of Lemma~\ref{lem:global-ends}.)  The
map $P\mapsto f_\chi^{n(P)}(P)$ is
locally a fixed iterate, jumps by one iterate exactly across $\ell$, where the
two values are identified, and satisfies $n(f_\chi(P))=n(P)-1$ for $P\in\Gamma$,
matching the identification of the edges of $\hat T_\lambda$.  It therefore induces a
holomorphic map $\Theta:X^{\rm tr}_\lambda\to X^\ell_{1-\lambda}$.  For the inverse, apply the same
construction in the $\chi$-plane of $E_0$ to the strips $\hat H_0$ and
$\tilde h_\lambda^{-1}(\hat H_\lambda)\subset\mathsf W_{1/100}\cap\mathcal G^{\rm c}$ (Step~1), and conjugate back by $\tilde h_\lambda$:
this gives a map $\Theta'$ from $X^\ell_{1-\lambda}$ to $X^{\rm tr}_\lambda$, locally a fixed iterate of
$E_\lambda$.  Both $\Theta'\circ\Theta$ and $\Theta\circ\Theta'$ send each point to a point of its own
orbit in the same half-open strip, so they are the identity.
Since orbit segments have length at most $(1+|t|)/20$ in $\chi$, $\Theta$ preserves the
ends ($\operatorname{Re}\chi\to\mp\infty$).

\emph{Step 3: marked points.}  By Proposition~\ref{prop:global-transport}(2) and
continuity of the lift, $\tilde h_\lambda$ maps the continued base orbit $\chi^*_n(\lambda_0)$ of
Lemma~\ref{lem:cusp-base} to $\chi^*_n(\lambda)$.  So $p_\lambda$ is represented by
$\chi^*_{n_0}(\lambda)$, which lies in $\hat H_\lambda\setminus f_\chi(\ell)$ because $n_0$ is the first crossing for
all $\lambda\in U$.  Hence $n(p_\lambda)=0$ and $\Theta(p_\lambda)$ is the marked point.

A biholomorphism of $\mathbb C^*$ preserving both ends and a marked point is the
identity in the normalised coordinates, so the two lifts $z$ agree; both
germs at $z=0$ are obtained by the local inverse of $E_\lambda^{n_0}$ with value $1$
at the marked point, so they agree.
\end{proof}

\begin{theorem}[Continuation to the upper Shell--Thron region]\label{thm:global}
Let $\Omega=b(\mathcal S^{\rm c\prime}_{r_3})$ be the corridor of Theorem~\ref{thm:cusp}.  The
continuation of Kneser's solution from $(\eta,\eta+r_5)$ into $\Omega\cap\mathcal U^+$ extends
analytically along every path in $\mathcal U^+=b(\mathbb H_\lambda)$, the upper half of the
Shell--Thron region, to a single-valued family of germs at $z=0$, jointly
holomorphic in $(b,z)$.  For every $b_0\in(1,\eta)$ there is a disc $D$ about
$z=0$ such that along every path in $\mathcal U^+$ ending at $b_0$ these germs
converge uniformly on $D$ to $\Kk^W_{b_0}$; hence the boundary value is $\Kk^W_{b_0}$ on
its regular domain, by the functional equation.
\end{theorem}

\begin{proof}
By Proposition~\ref{prop:global-family}, $\lambda\mapsto F^{\rm tr}_\lambda$ is a single-valued
holomorphic family of germs on $\mathbb H_\lambda$.  By Lemma~\ref{lem:global-compare} it
coincides with the cusp family $F_{1-\lambda}$ near $\lambda_0$.  The part of the
parameter domain of Theorem~\ref{thm:cusp} that lies in $\mathbb H_\lambda$,
$\{1-\varepsilon:\varepsilon\in\mathcal S^{\rm c\prime}_{r_3},\ \operatorname{Im}\varepsilon<0,\ |1-\varepsilon|<1\}$, is connected (it is the part
of a small sector lying inside the disc $|\varepsilon-1|<1$ and below the real axis) and contains $\lambda_0$, so the
two families agree on it by the identity theorem.  Hence $F^{\rm tr}$,
transferred to $b$ by Lemma~\ref{lem:global-param}, is the analytic
continuation of Kneser's solution along every path in $\mathcal U^+$ starting in
$\Omega\cap\mathcal U^+$, and it is single-valued there.

For the boundary values, shrink the discs $\Delta_{b_0}$ of
Proposition~\ref{prop:weld-parameter} to radius at most $(\eta-b_0)/2$, as in the
proof of Theorem~\ref{thm:cusp}(3).  For $b$ in such a disc,
$|\log b|\le\log b_0-\log(1-x/2)<\log b_0+\log(1+x)=\eee^{-1}$ with $x=(\eta-b_0)/b_0$,
and $g(\{|\lambda|<1\})\supset\{|y|<\eee^{-1}\}$ because $\min_{|\lambda|=1}|\lambda\eee^{-\lambda}|=\eee^{-1}$; so
$E_b$ has an attracting fixed point and $N_W\cap\{\operatorname{Im}b>0\}\subset\mathcal U^+$
(Lemma~\ref{lem:global-param}).  This set is connected and meets $\Omega\cap\mathcal U^+$
near $\eta$, where $\Kk^W$ agrees with the cusp family (Theorem~\ref{thm:cusp}(3));
by the identity theorem $\Kk^W=F^{\rm tr}$ on all of it.  The $z$-domain of $\Kk^W$
in Proposition~\ref{prop:weld-parameter} contains a disc $D$ about $0$ for all $b$
in a neighbourhood of $b_0$.  Every path in $\mathcal U^+$ ending at $b_0$ eventually lies
in $\Delta_{b_0}\cap\{\operatorname{Im}b>0\}$, where $F^{\rm tr}=\Kk^W$ is jointly holomorphic on a
neighbourhood of the closure of that half-disc times $D$; so the germs
converge uniformly on $D$ to $\Kk^W_{b_0}$.
\end{proof}

\begin{proposition}[Crossing the attracting real slice]\label{prop:attracting-real-slice}
The family of Theorem~\ref{thm:global} extends jointly holomorphically in
$(b,z)$ across every $b\in(\eee^{-\eee},1)$.  The local extensions agree on
overlaps and hence give a single-valued family of germs on a neighbourhood
of this interval.
\end{proposition}

\begin{proof}
Fix $x\in(-1,0)$ and choose a simply connected neighbourhood $V$ of $x$
in $\{0<|\lambda|<1\}$ which meets $\mathbb H_\lambda$.  Continue the branch
$s=\operatorname{Log}\lambda$ from $V\cap\mathbb H_\lambda$ to $V$.  The proof of
Proposition~\ref{prop:global-transport} uses only
$\operatorname{Re}s=\log|\lambda|<0$ to obtain
$|\mathsf c(\lambda)|<1$.  Thus its Beltrami coefficient, normalised
solution $h_\lambda$, and all four conclusions hold on $V$ by exactly the
same formulas.  The marked quotient and its uniformisation in
Proposition~\ref{prop:global-family} consequently give a jointly
holomorphic family of germs on $V\times\{z:|z|<\delta\}$ for some
$\delta>0$.  On $V\cap\mathbb H_\lambda$ this is the original transported
family, since both use the same normalised Beltrami solution and marked
quotient.  Any two such local extensions agree on an upper-half-plane
overlap, hence agree wherever both are defined by the identity theorem.
Finally $b'(x)=b(x)(1-x)\eee^{-x}\ne0$, and $b$ maps $(-1,0)$
increasingly onto $(\eee^{-\eee},1)$; local inversion transfers these germs
to the base coordinate.
\end{proof}

\begin{remark}[The punctured interior]\label{rem:stcover}
The same argument works globally in $s$.  Every object in
Propositions~\ref{prop:global-transport} and~\ref{prop:global-family} is given by one
formula in $s=\log\lambda$, and only $\operatorname{Re}s<0$ and $|\eee^s|<1$ are used.  So $s\mapsto F^{\rm tr}_s$ is
holomorphic and single-valued on the whole half-plane $\{\operatorname{Re}s<0\}$.  This half-plane is
the universal cover of the punctured interior $\mathcal{ST}^\circ\setminus\{1\}$ via $s\mapsto\exp g(\eee^s)$.
Hence the continuation of Kneser's family continues along every path in
$\mathcal{ST}^\circ\setminus\{1\}$, and $b=1$ is its only possible singularity in the interior.  A loop
about $b=1$ acts by $s\mapsto s+2\pi\ii$.  Whether this monodromy is trivial is open.
Numerically, the normalised modes $|c_1|/\Lambda$ of the boundary values $\Kk^W_b$ stay
near $0.08$ as $b\downarrow1$ ($0.0816$ at $b=1.01$), while $\Lambda\to1$.  So near $b=1$ the
separation from regular iteration is of order one.
\end{remark}

\begin{remark}[Relation to Paulsen's family]\label{rem:global-paulsen}
Theorem~\ref{thm:global} removes the restriction to paths near the real axis
in Theorem~\ref{thm:cusp}: starting from $\kappa_b$ on $(\eta,\eta+r_5)$ and entering
$\mathcal U^+$ inside $\Omega$, every further path in $\mathcal U^+$ gives $\Kk^W$ on all of
$(1,\eta)$.  The family of~\cite{Paulsen2019} may enter $\mathcal U^+$ through a boundary
arc far from $\eta$.  It coincides with ours if its germ at one base in $\mathcal U^+$
equals $F^{\rm tr}$ there.  Theorems~\ref{thm:far} and~\ref{thm:last-arc}
show that Kneser's family crosses the entire open upper boundary arc
holomorphically and enters with $F^{\rm tr}$, and Theorem~\ref{thm:upper} shows that the
whole continuation over the upper half-plane is single-valued, so this is
Paulsen's family (Corollary~\ref{cor:paulsen}).  The nine-digit agreement at $b=\sqrt2$
(Section~\ref{sec:invariant}) is consistent with this.
\end{remark}

\subsection{Crossing the boundary away from the cusp}\label{sec:far}
Theorem~\ref{thm:cusp} crosses the Shell--Thron boundary only near $\eta$.
Here we cross it along arcs far from the cusp, up to the multiplier angle
$\alpha_1=0.9353\ldots\pi$.  The tool is Kneser's own initial region: the
straight chord between the two fixed points and its image.  We show that it
stays a crescent, with the same marked point, on a collar that crosses the
boundary arc.  This needs no linearisation of the neutral fixed point, as in
Section~\ref{sec:cusp}.

\emph{Notation.}  For complex $\theta$ with $0<\operatorname{Re}\theta<\pi$ put
\[
 \mu(\theta)=\theta\cot\theta,\quad \lambda_\pm=\mu\pm\ii\theta,\quad
 r=\eee^{\mu},\quad a(\theta)=\frac{\theta\eee^{-\mu}}{\sin\theta},\quad
 b(\theta)=\eee^{a(\theta)} .
\]
We write $E_\theta(w)=\eee^{a(\theta)w}$; this is $E_{b(\theta)}$ with the logarithm of the base
continued along paths in $\theta$, and it is the map that enters the analytic
continuation of $F(z+1)=b^{F(z)}$ in the base.  $L_\pm=r\eee^{\pm\ii\theta}$ are fixed points
of $E_\theta$ with multipliers $aL_\pm=\lambda_\pm$, and $a=\lambda_+\eee^{-\lambda_+}$.  So if
$\lambda_+\in\mathbb H_\lambda$, then $a=\operatorname{Log}b$ is principal ($|\operatorname{Im}a|<\pi$) and $b(\theta)=b(\lambda_+(\theta))\in\mathcal U^+$
in the notation of Lemma~\ref{lem:global-param}.  For real $\theta\in(0,\pi)$, $b(\theta)>\eta$ is real, $L_\pm$ is the
conjugate pair of Proposition~\ref{prop:exterior-real-parameter} with $\theta_b=\theta$, and
$|\lambda_+|=\theta/\sin\theta>1$.  The affine coordinate
\[
 \zeta(w)=\frac{w/r-\cos\theta}{\ii\sin\theta}
\]
sends $L_\pm$ to $\pm1$ and the chord $\gamma_\theta=[L_-,L_+]$ to $[-1,1]$, and a direct
computation ($a w=\mu+\ii\theta\zeta$) gives
\[
 \zeta(E_\theta(w))=s_\theta(\zeta(w)),\qquad s_\theta(\zeta)=\frac{\eee^{\ii\theta\zeta}-\cos\theta}{\ii\sin\theta} .
\]
$s_\theta$ is entire, $s_\theta(\pm1)=\pm1$ and $s_\theta'(\pm1)=\lambda_\pm$, so
$g_\theta(t)=(s_\theta(t)-t)/(1-t^2)$ is entire in $t$.  Since $s_\theta(\zeta_1)=s_\theta(\zeta_2)$ only if
$\theta(\zeta_1-\zeta_2)\in2\pi\mathbb Z$, $s_\theta$ is injective on every set of diameter
$<2\pi/|\theta|$.

\emph{The parameter region.}  Let $P_0=\tfrac1{10}$, $P_1=\tfrac{11}5$, let $V(p)=\tfrac14$ for $p\le\tfrac85$
and $V(p)=\tfrac14-\tfrac7{60}(p-\tfrac85)$ for $p\ge\tfrac85$, and
\[
 \Theta_{\rm C}=\{p(1+\ii c):0<p\le P_0,\ 0\le c\le\tfrac1{10}\},\qquad
 \Theta_{\rm K}=\{p+\ii vp^2:P_0\le p\le P_1,\ 0\le v\le V(p)\},
\]
$\Theta=\Theta_{\rm C}\cup\Theta_{\rm K}$.  Its vertical fibre over $p$ is a segment
$\{p+\ii q:0\le q\le Q(p)\}$.

\begin{lemma}[Certified inequalities]\label{lem:far-cert}
For $\theta=p+\ii q\in\Theta$ and $-1\le t\le1$:
\begin{enumerate}
\item $\operatorname{Im}g_\theta(t)<0$ and $\operatorname{Im}\bigl(\overline{s_\theta'(t)}\,g_\theta(t)\bigr)<0$;
\item $\partial_q|\lambda_+(p+\ii q)|^2<0$, $|\lambda_+|<1$ on the top edge $q=Q(p)$, and $a'(\theta)\ne0$;
\item on $\Theta_{\rm C}$, $\operatorname{Re}\cot\theta>0$;
\item for $\operatorname{Re}\theta\le P_2=\tfrac{19}{10}$, $|\operatorname{Im}a(\theta)|<\pi$.
\end{enumerate}
\end{lemma}

\begin{proof}
See Appendix~\ref{app:collar}.
\end{proof}

By (1) at $t=1$, $\operatorname{Im}\lambda_+=-2\operatorname{Im}g_\theta(1)>0$.  By (2), on each fibre of $\Theta$ the
multiplier $\lambda_+$ of the continued fixed point $L_+$ crosses the unit circle exactly
once, at $q=q_{\rm c}(p)$, and
\[
 \Theta^+=\{\theta\in\Theta:|\lambda_+(\theta)|<1\}=\{p+\ii q\in\Theta:q>q_{\rm c}(p)\}
\]
is connected, with $a$ principal and $b(\Theta^+)\subset\mathcal U^+$.  (Below $q_{\rm c}$ the
logarithm $a$ need not be principal: for $p$ near $P_1$, $\operatorname{Im}a$ exceeds $\pi$ on the
fibre.  By (4) this does not happen for $p\le P_2$.)  The boundary
point $\beta(p)=b(p+\ii q_{\rm c}(p))$ has $\lambda_+=\eee^{\ii\alpha(p)}$ with $\alpha$ continuous,
$\alpha(p)\to0$ as $p\to0$ and $\alpha(P_1)=\alpha_1=0.9353\ldots\pi$; so $\beta$ runs through the whole
upper boundary arc $\{\exp(\eee^{\ii\alpha}\eee^{-\eee^{\ii\alpha}}):0<\alpha\le\alpha_1\}$.

\begin{lemma}[The crescent]\label{lem:far-crescent}
Let $\theta$ lie in a sufficiently small open neighbourhood $\Theta'$ of $\Theta\setminus\{0\}$.
\begin{enumerate}
\item $A_\theta=s_\theta([-1,1])$ is a simple arc from $-1$ to $1$, in the open lower
  half-plane except for its ends.  Let $H_\theta$ be the closed Jordan region bounded
  by $[-1,1]\cup A_\theta$.
\item $N_\theta(t,\sigma)=t+\sigma(s_\theta(t)-t)$ is a real-analytic diffeomorphism of
  $(-1,1)\times[0,1]$ onto $H_\theta\setminus\{\pm1\}$, with $N(t,0)=t$ and $N(t,1)=s_\theta(t)$.
\item $s_\theta$ maps a neighbourhood of $[-1,1]$ biholomorphically onto a neighbourhood
  of $A_\theta$, and the upper side of $[-1,1]$ to the side of $A_\theta$ facing $H_\theta$.
\item Let $X_\theta$ be $H_\theta\setminus\{\pm1\}$ with $t\in(-1,1)$ glued to $s_\theta(t)$.  In the
  coordinates $x=\log\frac{1-t}{1+t}$ and $\chi=\log\frac{\zeta-1}{\zeta+1}$ the map
  $\Phi_\theta(x+\ii\sigma)=\chi(N_\theta(t,\sigma))$ induces a quasiconformal homeomorphism from
  $C_0=\mathbb C/\ii\mathbb Z$ onto $X_\theta$.  Its Beltrami coefficient is holomorphic in $\theta$ at
  each fixed $x+\ii\sigma$ and bounded by some $k<1$ locally uniformly in $\theta$.
  Hence $X_\theta\cong\mathbb C^*$, with the ends at $\pm1$ as punctures.
\end{enumerate}
\end{lemma}

\begin{proof}
All conditions below are strict inequalities that hold on $\Theta$ by
Lemma~\ref{lem:far-cert}, uniformly on compact parts, so they persist on a
neighbourhood.  $|\theta|<2.45$ there (also on the rectangle used in
Theorem~\ref{thm:last-arc}), so $s_\theta$ is injective on every set of diameter
$<2\pi/2.45$, in particular on the disc of radius $1.28$.

(1) $\operatorname{Im}s_\theta(t)=(1-t^2)\operatorname{Im}g_\theta(t)<0$ for $|t|<1$, and $s_\theta$ is injective on $[-1,1]$.

(2) $\partial_tN=(1-\sigma)+\sigma s'$ and $\partial_\sigma N=(1-t^2)g$, so the Jacobian is
\[
 \operatorname{Im}\bigl(\overline{\partial_tN}\,\partial_\sigma N\bigr)
 =(1-t^2)\bigl[(1-\sigma)\operatorname{Im}g+\sigma\operatorname{Im}(\bar s'g)\bigr]<0
\]
for $|t|<1$.  $N$ is continuous on the square $[-1,1]\times[0,1]$ and maps its boundary
once around the Jordan curve $\partial H_\theta$, the vertical sides to the points $\pm1$.
For $\zeta_0\notin\partial H_\theta$ all preimages are interior and count with the same sign, so
their number is the winding number of $\partial H_\theta$ about $\zeta_0$, namely $1$ in
$H_\theta$ and $0$ outside.  An interior point cannot map to $\partial H_\theta$, since its image
would have nearby points outside $H_\theta$ with preimages.

(3) $s_\theta'\ne0$, and $s_\theta$ is injective near $[-1,1]$.  At $t\in(-1,1)$ the image of
the upward direction is $\ii s'(t)$, and the direction from $A_\theta$ into $H_\theta$ is
$-\partial_\sigma N(t,1)=-(1-t^2)g(t)$.  Both lie on the same side of the tangent
$s'(t)$, because $\operatorname{Im}(\overline{s'}\,\ii s')=|s'|^2>0$ and $\operatorname{Im}(\overline{s'}(-g))>0$.

(4) Near $t=1$ write $N-1=(t-1)m$ with $m=1-\sigma(1+t)g(t)$.  For $t<1$, $m\ne0$ by (2);
at $t=1$, $m=1+\sigma(\lambda_+-1)$, whose imaginary part is $\sigma\operatorname{Im}\lambda_+>0$ for $\sigma>0$.  Hence
\[
 \Phi_\theta=x+\ii\pi+R_\theta(t,\sigma),\qquad R_\theta=\log m+\log(1+t)-\log(N+1),
\]
with $R$ real-analytic in $(t,\sigma,\theta)$ up to $t=1$; similarly at $t=-1$.  Since
$dt/dx=-(1-t^2)/2$, one has $\partial_x\Phi=1-\tfrac12(1-t^2)\partial_tR\to1$ and
$\partial_\sigma\Phi\to(\lambda_+-1)/(1+\sigma(\lambda_+-1))$ as $t\to1$, whose imaginary part is positive.  So
$D\Phi_\theta$ is continuous and nonsingular on $[-\infty,\infty]\times[0,1]$, with constant
orientation.  Its Beltrami coefficient
$(\partial_x\Phi+\ii\partial_\sigma\Phi)/(\partial_x\Phi-\ii\partial_\sigma\Phi)$ is continuous there and in $\theta$, of modulus
$<1$, and holomorphic in $\theta$ because $\Phi_\theta(x+\ii\sigma)$ is.  $\Phi(x)=\chi(t)$ and
$\Phi(x+\ii)=\chi(s_\theta(t))$, so $\Phi$ respects the two identifications.  The rest is as in
Lemma~\ref{lem:crescent-quotient}(1).
\end{proof}

\emph{The marked point.}  Let $w_{-1}=0$ and $w_n=E_\theta^n(1)$, $n\ge0$, and write
$\zeta_n=\zeta(w_n)$.  We say $w$ lies \emph{above the chord} if $\operatorname{Im}\zeta(w)>0$; such points are
not in $H_\theta$.  If $w_{-1},\dots,w_{n-1}$ lie above the chord and $\zeta_n\in H_\theta\setminus A_\theta$, we
call $n=n_\theta$ the \emph{entry index} and the class $\mathsf m_\theta\in X_\theta$ of $\zeta_n$ the
\emph{marked point}.  It is defined intrinsically.

\begin{lemma}[First entry on the cone]\label{lem:far-cone}
For $\theta\in\Theta_{\rm C}$ the entry index exists.
\end{lemma}

\begin{proof}
An analytic estimate of the orbit in the coordinate $\log\frac{\zeta-1}{\zeta+1}$; see
Appendix~\ref{app:collar}.
\end{proof}

\begin{lemma}[The marked point on $\Theta$]\label{lem:far-mark}
The entry index exists for every $\theta\in\Theta$, and $\theta\mapsto\mathsf m_\theta$ is a holomorphic section
of the family $X_\theta$ near $\Theta$: near each $\theta_0$ it is represented by a fixed iterate
$\zeta_n(\theta)$, either in the interior of $H_\theta$ or in the seam chart given by $s_\theta$.
\end{lemma}

\begin{proof}
On $\Theta_{\rm C}$ this is Lemma~\ref{lem:far-cone}.  On $\Theta_{\rm K}$ the orbit is enclosed on
parameter boxes and the entry is certified box by box (Appendix~\ref{app:collar}).

Near $\theta_0$: if $\zeta_n(\theta_0)$ is in the interior of $H$, $n_\theta=n$ nearby.  If $\zeta_n(\theta_0)\in(-1,1)$,
then $n_\theta\in\{n,n+1\}$ nearby, according to the side of the chord, and $\zeta_n$ and
$\zeta_{n+1}=s_\theta(\zeta_n)$ are the same point of $X_\theta$ in the seam chart.  The conditions
$\operatorname{Im}\zeta_m>0$, $m<n$, are open.  So $\mathsf m_\theta$ is locally the class of the single holomorphic
function $\zeta_n(\theta)$.
\end{proof}

\begin{definition}\label{def:far-solution}
For $\theta$ near $\Theta$ let $Z_\theta:X_\theta\to\mathbb C^*$ be the uniformisation taking the
$L_+$-end to $0$, the $L_-$-end to $\infty$ and $\mathsf m_\theta$ to $1$, and let
$z=(2\pi\ii)^{-1}\log Z_\theta$ be lifted so that $\mathsf m_\theta$ has $z=n_\theta$.  With $\tilde F_\theta$ the inverse, extended by
$\tilde F(z+1)=E_\theta(\tilde F(z))$ as in Lemma~\ref{lem:crescent-quotient}(1), put
\[
 F^{\rm far}_\theta(z)=G_\theta\bigl(\tilde F_\theta(z+n_\theta)\bigr),
\]
where $G_\theta$ is the branch of $E_\theta^{-n_\theta}$ along the orbit $w_{n_\theta}\mapsto\dots\mapsto w_0=1$ (for
$n_\theta=-1$, $G_\theta=E_\theta$).  Then $F^{\rm far}_\theta(0)=1$, and $F(z+1)=\eee^{a(\theta)F(z)}$ holds as an identity
of germs, with the corresponding branches along $w_{n_\theta+1}\mapsto\dots\mapsto w_1$.
\end{definition}

The sign of the lift, i.e.\ that crossing the seam from the chord to $A_\theta$ adds
$+1$, is locally constant in $\theta$.  At real $\theta$ it is Kneser's, by
Proposition~\ref{prop:far-real} below.

\begin{proposition}[Holomorphic dependence]\label{prop:far-holo}
$(\theta,z)\mapsto F^{\rm far}_\theta(z)$ is jointly holomorphic on a neighbourhood of
$(\Theta\setminus\{0\})\times\{0\}$.
\end{proposition}

\begin{proof}
This is the proof of Proposition~\ref{prop:cusp-holo}, with $M_\varepsilon$ replaced by
$N_\theta$ in the coordinates of Lemma~\ref{lem:far-crescent}(4).  Fix $\theta_0$.  Then
$\mathcal J_\theta=\Phi_\theta\circ\Phi_{\theta_0}^{-1}$ is quasiconformal, respects the seams, and is holomorphic in $\theta$
at fixed source points.  Its Beltrami coefficient $\upsilon_\theta$ is holomorphic in $\theta$ and, by
the uniform continuity in Lemma~\ref{lem:far-crescent}(4), is bounded by $\tfrac1{10}$
for $\theta$ near $\theta_0$.  Ahlfors--Bers and the $\partial_{\bar\theta}$-cancellation give a
uniformisation $Z_\theta$ that is jointly holomorphic in $(\theta,\zeta)$ in the interior of
$H_\theta$.  Near a point $\zeta_0$ of the chord, which does not move with $\theta$ in these
coordinates, $Z_\theta$ is given by $Z_\theta$ below the chord and by $Z_\theta\circ s_\theta$ above it.  Both
are continuous up to the chord, and they agree there.  So $\theta\mapsto Z_\theta(\zeta)$ is a locally
uniform limit of holomorphic functions for each $\zeta$, and $Z_\theta(\zeta)$ is holomorphic in
$\zeta$ by Morera's theorem.  Hartogs' theorem then gives joint holomorphy near the
chord.  The marked point is holomorphic
by Lemma~\ref{lem:far-mark}.  The inverse branches $G_\theta$ are finitely many
logarithms along an orbit that avoids $0$, holomorphic in $\theta$.
\end{proof}

\begin{proposition}[The real slice]\label{prop:far-real}
For real $\theta\in(0,P_1]$, $F^{\rm far}_\theta$ is Kneser's real-normalised solution $\kappa_{b(\theta)}$
near $z=0$.
\end{proposition}

\begin{proof}
For real $\theta$, $\zeta^{-1}([-1,1])$ is the chord $\gamma_b$ and $\zeta^{-1}(A_\theta)=E_b(\gamma_b)$, so $H_\theta$ is
Kneser's initial region $H_b$ and $X_\theta$ is the cylinder of
Proposition~\ref{prop:exterior-real-parameter}, with the upper end at $L_+$ sent to $0$
in both.  For real $w$, $\operatorname{Im}\zeta(w)=(\cos\theta-w/r)/\sin\theta$, which is positive exactly when
$w<u_b=r\cos\theta$, and $H_b\cap\mathbb R=[u_b,r_b)$ with $r_b=E_b(u_b)$.  The real orbit
$0=w_{-1}<1=w_0<w_1<\dots$ increases to $+\infty$.  Let $n$ be the first index with $w_n\ge u_b$.  If
$n\ge0$ then $w_{n-1}<u_b$, so $w_n<E_b(u_b)=r_b$.  If $n=-1$ then $u_b\le0<r_b$.  Either way
$w_n\in[u_b,r_b)$, so $n=n_\theta$, and $\mathsf m_\theta$ is the point that
Proposition~\ref{prop:exterior-real-parameter} represents by $w=1$ through the real
dynamical charts.  Its Kneser coordinate is $n_\theta$.  The normalisations of the lifts
agree, and the branches $G_\theta$ are the real logarithms.
\end{proof}

\begin{theorem}[Crossing the upper boundary away from the cusp]\label{thm:far}
Let $p\in(0,P_1]$, and continue Kneser's real-normalised solution $\kappa_b$ at
$b=b(p)>\eta$ analytically in the base along $q\mapsto b(p+\ii q)$, $0\le q\le Q(p)$, with
$\log b=a(p+\ii q)$ continued along the path.
\begin{enumerate}
\item The continuation exists along the whole path and is holomorphic in a
  neighbourhood of it.  The multiplier $\lambda_+$ of the continued fixed point crosses
  the unit circle once, at $q=q_{\rm c}(p)$, where $b=\beta(p)$ lies on the upper
  Shell--Thron boundary arc; near $\beta(p)$ the continuation is a holomorphic family
  on a full neighbourhood of $\beta(p)$.
\item For $q>q_{\rm c}(p)$ the base lies in $\mathcal U^+$, $a=\operatorname{Log}b$, and the continuation
  equals $F^{\rm tr}_{b(p+\ii q)}$, the family of Theorem~\ref{thm:global}.
\end{enumerate}
In particular, along these paths Kneser's family crosses every point of the upper
boundary arc $\{\exp(\eee^{\ii\alpha}\eee^{-\eee^{\ii\alpha}}):0<\alpha\le\alpha_1\}$, $\alpha_1=0.9353\ldots\pi$,
holomorphically, and it enters the upper Shell--Thron region with the germ $F^{\rm tr}$.
Its boundary values on $(1,\eta)$ are therefore $\Kk^W_b$, by Theorem~\ref{thm:global}.
For $p\le P_2$ ($\alpha\le0.7531\ldots\pi$) the whole path lies on the principal sheet,
$|\operatorname{Im}\log b|<\pi$.
\end{theorem}

\begin{proof}
By Proposition~\ref{prop:far-holo} and Lemma~\ref{lem:far-cert}(2), $b$ is locally
biholomorphic on a neighbourhood $\Theta'$ of $\Theta\setminus\{0\}$ (as is $a$), and $F^{\rm far}$ is a
holomorphic family of germs there.  By Proposition~\ref{prop:far-real} it equals $\kappa$ on
the real segment.  So along every path in $\Theta'$ starting on that segment it is the
analytic continuation of $\kappa$ in the base, with the logarithm $a$ continued along
the path.

Next, $a(\theta)=\eee^{-1}(1+\theta^2/2+O(\theta^4))$, so $b'(\theta)=b\,a'(\theta)$ has positive real part for
small $p$ on $\Theta_{\rm C}$, and there $a$ is principal.  Hence $\operatorname{Im}b(p+\ii q)=\int_0^q\operatorname{Re}b'(p+\ii q')\,dq'>0$ for $q>0$, and for $p$
small the vertical fibre is carried into $D^+\cup(\eta,\eta+r_5)\subset\Omega$.  By
Theorem~\ref{thm:cusp}(1) the continuation of $\kappa$ along it agrees with the family
on $\Omega$.  By Theorem~\ref{thm:global} that family equals $F^{\rm tr}$ on $\Omega\cap\mathcal U^+$.  So
$F^{\rm far}=F^{\rm tr}\circ b$ on a nonempty open subset of the connected set $\Theta^+$ (interior), and
hence on all of it, by the identity theorem.  The remaining assertions are the
description of $\Theta^+$ and $\beta$ after Lemma~\ref{lem:far-cert}.
\end{proof}

\begin{theorem}[Completion of the open upper boundary arc]\label{thm:last-arc}
For every $0.935\pi\le\alpha<\pi$, put
\[
 \beta(\alpha)=\exp\!\bigl(\eee^{\ii\alpha}
                         \eee^{-\eee^{\ii\alpha}}\bigr).
\]
Kneser's real-normalised family admits analytic continuation, on the principal
logarithm sheet, from a real base $b>\eta$ to a full neighbourhood of
$\beta(\alpha)$.  This continuation crosses the upper Shell--Thron boundary
holomorphically and agrees with $F^{\rm tr}$ on its inner side.  Consequently
the same assertion holds at every point of the open upper boundary arc
$0<\alpha<\pi$.  The endpoint $\beta(\pi)=\eee^{-\eee}$ is excluded.
\end{theorem}

\begin{corollary}[Entering from outside]\label{cor:band}
There is an open neighbourhood $N$ of the open upper boundary arc
$\{\beta(\alpha):0<\alpha<\pi\}$ and a single-valued holomorphic family of germs on
$N\cup\mathcal U^+$ that equals $F^{\rm tr}$ on $\mathcal U^+$.  On $N\setminus\overline{\mathcal U^+}$ it is the
continuation of Kneser's solution along paths that stay outside the closed
Shell--Thron region on the principal sheet.  Thus for every $0<\alpha<\pi$ Kneser's
family, continued from $(\eta,\infty)$ through the exterior, reaches $\beta(\alpha)$ and
enters $\mathcal U^+$ with the germ $F^{\rm tr}$.
\end{corollary}

\begin{proof}
Theorems~\ref{thm:far} and~\ref{thm:last-arc} give, about each point of the open
arc, a disc on which the continuation is holomorphic and equals $F^{\rm tr}$ on the
inner side.  Two such discs meet in a connected set containing inner points, so
the local families agree on overlaps, and $F^{\rm tr}$ is single-valued on $\mathcal U^+$.
Take $N$ to be the union of the discs, shrunk so that $N\setminus\overline{\mathcal U^+}$ is a connected
band along the arc.  For $p$ small the vertical path of Theorem~\ref{thm:far}
runs from the real exterior ray to this band outside $\overline{\mathcal U^+}$, on the principal
sheet.  Moving along the band then reaches the outer side of every $\beta(\alpha)$,
and $|\operatorname{Im}a|<\pi$ near the arc because $|a|\le\eee$ there.
\end{proof}

\begin{proof}
By Appendix~\ref{app:lastarc}, $\lambda_+$ is injective on $R=[2.19,2.31]+\ii[0.71,0.78]$, the
branch $\theta(\alpha)\in R$ of $\lambda_+(\theta)=\eee^{\ii\alpha}$, $0.935\pi\le\alpha<\pi$, satisfies the inequalities of
Lemma~\ref{lem:far-cert}(1) on $(-1,1)$ on a complex neighbourhood, and $w_{-1}=0$ lies in the
crescent, so the entry index is $-1$.  Lemma~\ref{lem:far-crescent} and
Proposition~\ref{prop:far-holo} apply on a complex parameter neighbourhood
of every $\alpha<\pi$.  Also $\lambda_+'\ne0$ and
$a'=\eee^{-\lambda_+}(1-\lambda_+)\lambda_+'\ne0$, so $b=\eee^a$ is a valid
local parameter.  These neighbourhoods carry one holomorphic family of
marked crescent quotients: the inverse branch is unique on $R$, and their
constructions agree on overlaps.

At $\alpha=0.935\pi<\alpha_1$ this is exactly the crescent and the first
mark of Theorem~\ref{thm:far}.  The families therefore agree in an open
parameter overlap.  Analytic continuation along the connected arc, followed
by the identity theorem on the inner overlaps, identifies the new family
with $F^{\rm tr}$ wherever it meets $\mathcal U^+$.

Finally choose $p\le P_2$ in Theorem~\ref{thm:far}, follow its vertical path
from the real exterior ray into $\mathcal U^+$, then follow a path in
$\mathcal U^+$ to the inner side of $\beta(\alpha)$ and use the local
crescent family to cross the boundary.  The first segment is on the principal
sheet by Lemma~\ref{lem:far-cert}(4).  On $\mathcal U^+$ the logarithm is
$a=\lambda\eee^{-\lambda}=\operatorname{Log}b$; at the target boundary
$|a|\le\eee<\pi$, so the last short segment is principal as well.
This gives the asserted principal-sheet path and completes the proof.
\end{proof}

\begin{remark}\label{rem:far}
(1) The boundary point $b_c\approx1.5217+0.0148\ii$ ($\alpha=2\pi(\sqrt2-1)/5$) tested
numerically in Section~\ref{sec:bndtest} lies on the arc covered here.  The
smooth join observed there agrees with the holomorphic crossing proved here.
(2) The endpoint $b=\eee^{-\eee}$ itself is treated in Section~\ref{sec:endpoint}.  There the
crescent keeps its angle $\arg\lambda_+\to\pi$ at $L_+$; only the interpolation used here degenerates.
(3) The identification with the family of~\cite{Paulsen2019} is completed in
Section~\ref{sec:upper} (Corollary~\ref{cor:paulsen}).
\end{remark}

\subsection{The whole upper half-plane}\label{sec:upper}
Theorems~\ref{thm:global}, \ref{thm:far} and~\ref{thm:last-arc} describe the
continuation inside the upper Shell--Thron region and across its boundary arc.
Here we show that Kneser's straight chord also carries the family over the
whole exterior part of the upper half-plane.  So the continuation is a single
holomorphic family on $\mathbb H=\{\operatorname{Im}b>0\}$, and there is no path dependence.

Write $\overline{\mathcal{ST}}$ for the closed Shell--Thron region, $g(\lambda)=\lambda\eee^{-\lambda}$, and
\[
 \mathcal E^+=\mathbb H\setminus\overline{\mathcal{ST}},\qquad
 \mathbb H=\mathcal E^+\cup\{\beta(\alpha):0<\alpha<\pi\}\cup\mathcal U^+ .
\]
On $\mathbb H$ we use the principal logarithm $a=\operatorname{Log}b$, $0<\operatorname{Im}a<\pi$.

\begin{lemma}[The exterior parameter]\label{lem:upper-param}
$\mathcal E^+$ is simply connected.  The fixed point $L_+$ of $E_b$ continued from
the real exterior ray (the one with $\operatorname{Im}L_+>0$ there) is single-valued and
holomorphic on $\mathcal E^+$.  Its multiplier $\lambda(b)$ maps $\mathcal E^+$ biholomorphically onto a
domain $\Lambda$, with inverse $\lambda\mapsto\exp g(\lambda)$, and
\[
 \Lambda\subset\{\lambda:\ |\lambda|>1,\ 0<\operatorname{Im}\lambda<\pi,\ 0<\operatorname{Im}g(\lambda)<\pi,\ g(\lambda)\notin g(\overline{\mathbb D})\}.
\]
\end{lemma}

\begin{proof}
$K=\overline{\mathcal{ST}}\cap\overline{\mathbb H}$ is a closed topological disc: it is bounded by the Jordan arc
$\beta$ and the segment $[\eee^{-\eee},\eta]$.  The complement of $\mathcal E^+$ in the Riemann sphere,
$(\hat{\mathbb C}\setminus\mathbb H)\cup K$, is connected, so $\mathcal E^+$ is simply connected.  A fixed point of $E_b$ is multiple only if its multiplier is
$1$, i.e.\ $a=1/\eee$, which does not occur for $0<\operatorname{Im}a<\pi$.  So $L_+$ continues without
monodromy.  $L_+=\eee^\lambda$ and $a=g(\lambda)$
(Lemma~\ref{lem:param}), so $b=\exp g(\lambda)$ and $\lambda$ is injective.  All fixed points of $E_b$,
$b\in\mathcal E^+$, are repelling, so $|\lambda|>1$ and $g(\lambda)=a\notin g(\overline{\mathbb D})$.  If $\operatorname{Im}\lambda=0$ then $a$
is real; if $\operatorname{Im}\lambda=\pi$ then $\operatorname{Im}a=-\pi\eee^{-\operatorname{Re}\lambda}<0$.  Neither happens on $\mathcal E^+$, and
near the real ray $\lambda=\theta\cot\theta+\ii\theta$ with small $\theta>0$; by connectedness
$0<\operatorname{Im}\lambda<\pi$.
\end{proof}

For $\lambda\in\Lambda$ let $\theta=\theta(\lambda)\ne0$ be a root of $(\lambda-2\ii\theta)\eee^{2\ii\theta}=\lambda$.  (The root $\theta=0$
is always present and is excluded.)  For $\theta\ne0$ the equation is equivalent to
$\lambda_+(\theta)=\lambda$.  Then $\lambda-2\ii\theta$ is the multiplier of a second fixed point
$L_-=\eee^{\lambda-2\ii\theta}$, so the chord construction of Section~\ref{sec:far} is defined with
this $\theta$.  $g$ has no critical point on $\Lambda$, so $\lambda$ is a local holomorphic parameter.

\begin{lemma}[Scope of Section~\ref{sec:far}]\label{lem:far-scope}
Let $\Theta''$ be an open set of $\theta$ with $0<|\theta|<\pi$ on which the inequalities of
Lemma~\ref{lem:far-cert}(1) hold for $-1<t<1$, $\operatorname{Im}\lambda_+>0>\operatorname{Im}\lambda_-$, and the entry index of the
marked point exists.  Then Lemma~\ref{lem:far-crescent}, the local statement of
Lemma~\ref{lem:far-mark}, Definition~\ref{def:far-solution} and
Proposition~\ref{prop:far-holo} hold on $\Theta''$.
\end{lemma}

\begin{proof}
Those proofs use only these properties and the injectivity of $s_\theta$ on a
neighbourhood of $[-1,1]$.  The injectivity follows from $|\theta|<\pi$:
$s_\theta(\zeta_1)=s_\theta(\zeta_2)$ forces $\theta(\zeta_1-\zeta_2)\in2\pi\mathbb Z$, and $|2\pi/\theta|>2$.  The bound $|\theta|<2.45$ in
Lemma~\ref{lem:far-crescent} was used only for this.  At the ends, $\operatorname{Im}\lambda_\pm$ play the
roles of $\operatorname{Im}g(\pm1)$ in part~(4).
\end{proof}

\begin{proposition}[Certified covering of $\Lambda$]\label{prop:upper-cert}
There is a continuous choice of $\theta(\lambda)$ on $\Lambda$, real on the image of the real
exterior ray, such that for every $\lambda\in\Lambda$ the inequalities of
Lemma~\ref{lem:far-cert}(1) hold on $(-1,1)$, $\operatorname{Im}(\lambda-2\ii\theta)<0$, $|\theta|<\pi$, and the
entry index of the marked point exists.  The certificate \texttt{certify\_upper\_halfplane.py} has four parts (tails,
core, cusp, interfaces), described in Appendix~\ref{app:upper}.
\end{proposition}

\begin{proof}
The listed checks are outward-rounded Arb computations.  Every point of $\Lambda$
lies in a certified box, because of the inclusion in Lemma~\ref{lem:upper-param}.
The roots are simple, since $\partial_\theta$ of the equation is
$2\ii\eee^{2\ii\theta}(\lambda-2\ii\theta-1)$ and $\lambda-2\ii\theta\ne1$ ($L_-$ is repelling).  So the certified
local roots, which agree on all overlaps, define one continuous branch.
\end{proof}

\begin{theorem}[Kneser's tetration on the upper half-plane]\label{thm:upper}
Kneser's real-normalised solution $\kappa_b$, $b>\eta$, continues analytically along
every path in $\mathbb H=\{\operatorname{Im}b>0\}$, with the principal logarithm.  The result is a
single-valued family of germs at $z=0$, jointly holomorphic in $(b,z)$.  It is
given on $\mathcal E^+$ by the straight-chord construction with the marked point of
Section~\ref{sec:far}, on $\mathcal U^+$ by $F^{\rm tr}$, and it is holomorphic across the open
upper boundary arc.  In particular, its boundary values on $(1,\eta)$ are $\Kk^W_b$,
whatever path in $\mathbb H$ is used.  The continuation through the lower half-plane is
the complex conjugate family.
\end{theorem}

\begin{proof}
By Proposition~\ref{prop:upper-cert} and Lemma~\ref{lem:far-scope}
(with $\lambda$ as parameter), the chord construction
with $\theta(\lambda)$ and the first-entry mark gives a holomorphic family on $\Lambda$, hence on
$\mathcal E^+$.  Near the image of the real ray it is Kneser's solution, by
Proposition~\ref{prop:far-real}, since there $\theta$ is real.  Since $\mathcal E^+$ is connected,
the family is the analytic continuation of $\kappa$ along every path in $\mathcal E^+$.  Here the
image $\Gamma$ of the real ray lies in $\partial\Lambda$.  The certified cells straddle it,
$\theta(\lambda)$ extends holomorphically across $\Gamma$, it is real there by the anchor
in~(I), and Proposition~\ref{prop:far-real} applies.  Shrink the neighbourhood $N$
of Corollary~\ref{cor:band} into $\mathbb H$, so that $N\cap\mathcal E^+$ is the connected outer band.  By
that corollary, near the open arc this continuation extends
holomorphically across the arc and equals $F^{\rm tr}$ on the inner side; $F^{\rm tr}$ is
single-valued on $\mathcal U^+$ by Theorem~\ref{thm:global}.  The three pieces agree on
their overlaps, so they define one holomorphic family on $\mathbb H$.  Since $\mathbb H$ is
simply connected, analytic continuation of $\kappa$ along any path in $\mathbb H$ gives it.
The boundary values on $(1,\eta)$ are those of Theorem~\ref{thm:global}.  For the lower
half-plane use $\overline{\kappa_{\bar b}(\bar z)}=\kappa_b(z)$.
\end{proof}

\begin{corollary}[Paulsen's family]\label{cor:paulsen}
Suppose the family of~\cite{Paulsen2019} is understood, as in its definition, as the
analytic continuation in $b$ of Kneser's solution from $(\eta,\infty)$, along paths
in the upper half-plane and with the principal logarithm.  Then it is the family of
Theorem~\ref{thm:upper}.  In particular it enters the Shell--Thron region with the
germ $F^{\rm tr}$ across every point of the open boundary arc.  Its boundary values on
$(1,\eta)$ are $\Kk^W_b$, and the separation limits of
Corollary~\ref{cor:conjecture} hold for it.
\end{corollary}

\subsection{The endpoint $b=\eee^{-\eee}$ and the real segment below $1$}\label{sec:endpoint}
At $b=\eee^{-\eee}$ the fixed point $L_+$ has multiplier $\lambda_+=-1$.  In
Theorem~\ref{thm:last-arc} the straight crescent seemed to degenerate there, because
$f_1(1)=-\tfrac12\operatorname{Im}\lambda_+\to0$.  The crescent itself does not degenerate: its angle at
$L_+$ is $\arg\lambda_+\to\pi$.  What degenerates is the interpolation $t+\sigma(s_\theta(t)-t)$,
whose factor $1+\sigma(\lambda_+-1)$ vanishes at $\sigma=\tfrac12$.  In the coordinate of the pair of
fixed points the end is harmless, and the crescent can be continued through
$\lambda_+=-1$ and across the real axis.

Put $\chi=\log\frac{\zeta-1}{\zeta+1}$, so that the chord is $\{\operatorname{Im}\chi=\pi\}$, with $x=\log\frac{1-t}{1+t}$.  The
image arc is $\Gamma(x)=x+\ii\pi+\varphi(x)$, where, with the entire function $S(w)=\sinh w/w$,
\[
 \varphi(t)=\ii\theta+\log S\bigl(\tfrac{\ii\theta(t-1)}2\bigr)-\log S\bigl(\tfrac{\ii\theta(t+1)}2\bigr),
\]
with principal logarithms where $\operatorname{Re}S>0$.  Then $\eee^{\varphi(1)}=\lambda_+$ and $\eee^{-\varphi(-1)}=\lambda_-$, and under (b)
below $\varphi(\pm1)$ are the logarithms with imaginary part in $(0,2\pi)$.  This follows from $s_\theta(t)\mp1=\eee^{\pm\ii\theta}(\eee^{\ii\theta(t\mp1)}-1)/(\ii\sin\theta)$.  We write
$\varphi_x=d\varphi/dx=-\tfrac12(1-t^2)\,d\varphi/dt$, which is regular at $t=\pm1$.  The mark has the
closed form $\chi(\zeta(0))=\log\frac{\ii\cot\theta-1}{\ii\cot\theta+1}=2\ii\theta$ (mod $2\pi\ii$).

\begin{lemma}[The crescent in the pair coordinate]\label{lem:chi-crescent}
Let $\Theta''$ be an open set of $\theta$ on which, for $-1\le t\le1$:
\begin{enumerate}
\item[(f)] $|\theta|<\pi$, and $\operatorname{Re}S>0$ at both arguments above;
\item[(a)] $\operatorname{Re}(1+\varphi_x)>0$;
\item[(b)] $0<\operatorname{Im}\varphi<2\pi$;
\item[(c)] $\operatorname{Im}\bigl(\overline{(1+\varphi_x)}\,\varphi\bigr)>0$;
\item[(d)] $\operatorname{Im}\varphi<\pi$ wherever $\operatorname{Re}\Gamma=0$;
\item[(e)] the point $2\ii\theta$ ($=\chi(\zeta(0))$, $\zeta(0)=\ii\cot\theta$), which must have imaginary part in
  $(\pi,3\pi)$, lies in the strip between $\operatorname{Im}\chi=\pi$ and $\Gamma$.
\end{enumerate}
Then the region $H_\theta$ between the chord and $s_\theta([-1,1])$ is a Jordan region, and
$s_\theta$ maps the outer side of the chord to the inner side of the arc.
$M(x,\sigma)=x+\ii\pi+\sigma\varphi(x)$ is a flat model of the quotient with the properties of
Lemma~\ref{lem:far-crescent}(4) (with $M$ in place of $N_\theta$), $w_{-1}=0$ lies in $H_\theta$ (entry
index $-1$), and Definition~\ref{def:far-solution} and Proposition~\ref{prop:far-holo}
apply on $\Theta''$.  Parts (1) and (2) of Lemma~\ref{lem:far-crescent} are not asserted.  The
arc need not lie in the lower half-plane, and it may leave $\zeta=1$ on the side $\zeta>1$.
\end{lemma}

\begin{proof}
By (a) $\Gamma$ is a graph over $\operatorname{Re}\chi$.  By (b) and (c) the map $M$ is a local
diffeomorphism.  Its Jacobian is $\operatorname{Im}\varphi+\sigma\operatorname{Im}(\overline{\varphi_x}\varphi)$, linear in $\sigma$, and $M$
is bijective on the two boundary lines.  So, as in Lemma~\ref{lem:crescent}(3), it is
a diffeomorphism onto the strip between $\operatorname{Im}\chi=\pi$ and $\Gamma$.  By (b) the strip embeds
in $\mathbb C/2\pi\ii\mathbb Z$.  By (d) it avoids $\chi\equiv0$, which corresponds to $\zeta=\infty$.  So its image
under $\zeta=(1+\eee^\chi)/(1-\eee^\chi)$ is a bounded Jordan region bounded by the chord and the arc.

As $x\to\mp\infty$, $\varphi\to\log\lambda_+$ and $\varphi\to-\log\lambda_-$ exponentially fast.  Both limits have
imaginary part in $(0,2\pi)$ by (b), so $M$ is asymptotically affine and nondegenerate at
both ends.  This holds also for $\lambda_+=-1$, where $\operatorname{Im}\log\lambda_+=\pi$.  The quotient of the
strip is therefore quasiconformally $\mathbb C/\ii\mathbb Z$, with punctures at both ends, and
$M$ is holomorphic in $\theta$ at fixed $(x,\sigma)$.

The seam map $s_\theta$ is injective near $[-1,1]$ because $|\theta|<\pi$.  In the coordinate
$\chi$ the map $s_\theta$ sends the outward normal $-\ii$ of the chord to $-\ii(1+\varphi_x)$.  The
direction from $\Gamma$ into the strip is $-\varphi=-\partial_\sigma M$.  By (c) both have
$\operatorname{Im}(\overline{(1+\varphi_x)}\,v)<0$, so they lie on the same side of $\Gamma$: the outer side of the
chord goes to the inner side of the arc.  By (e), $w_{-1}=0$ lies in the crescent.  Since $-1$
is the smallest index, the entry index is $-1$, and no condition on earlier points
is needed.  The proofs of
Lemma~\ref{lem:far-crescent}(4), Definition~\ref{def:far-solution} and
Proposition~\ref{prop:far-holo} use only these facts.  The quotient does not depend on
the interpolation, so on overlaps with Section~\ref{sec:far} the two constructions
give the same family.
\end{proof}

\begin{theorem}[Holomorphy at $b=\eee^{-\eee}$]\label{thm:endpoint}
The family of Theorem~\ref{thm:upper} extends holomorphically to a neighbourhood of
the real segment
\[
 \bigl[\exp(-12\eee^{12}),\ \exp(-\tfrac15\eee^{1/5})\bigr]\approx[\eee^{-1.95\cdot10^6},\ 0.7833],
\]
which contains $b=\eee^{-\eee}$.  At every base of this neighbourhood, including $b=\eee^{-\eee}$,
it is given by the straight-chord construction with the mark $w_{-1}=0$.  In
particular its boundary values on this segment are real-analytic.
\end{theorem}

\begin{proof}
By Appendix~\ref{app:endpoint}, (a)--(f) of Lemma~\ref{lem:chi-crescent} hold for all $\lambda$
in the square $[-1.1,-0.9]\times[-0.1,0.1]$ and the strip $[-12,-\tfrac15]\times[-10^{-3},10^{-3}]$,
with one branch $\theta(\lambda)$, which lies in $R'=[2.20,2.45]\times[0.62,0.95]$ on the square.
All checks are strict, so they persist on an open neighbourhood of the closed
square and strip.  For $\lambda=\eee^{\ii\alpha}$ near $-1$ the root in $R'$ is the one used in
Theorem~\ref{thm:last-arc}.  Indeed, that theorem's contraction disc of radius $0.00019$
about $\theta_*$ lies in $R\cap R'$, and $\lambda_+$ is injective on $R'$.  So the two roots agree near
$\alpha=\pi$, and by continuity on the whole arc inside the square.  So on the upper part of the
square the new family is that of Theorem~\ref{thm:last-arc}.  It therefore equals
$F^{\rm tr}$ inside and the family of Theorem~\ref{thm:upper} outside.  The whole strip cannot
be compared with Theorem~\ref{thm:upper} in the base variable, because
$\operatorname{Im}g(x+\ii y)\approx y(1-x)\eee^{-x}$ is large for $x\ll0$.  Let $U$ be the tube of points of the strip near $[-12,-\tfrac15]$ with
$|\operatorname{Im}g|<\pi$; its width is about $\pi\eee^{x}/(1-x)$.  On the real segment $g$ is real, strictly
increasing and has $g'=(1-\lambda)\eee^{-\lambda}>0$.  So $b=\exp g$ is injective on $U$, after shrinking it,
and maps it onto a neighbourhood of the stated $b$-segment, with $b(-1)=\eee^{-\eee}$.
$U\cap\{\operatorname{Im}\lambda>0\}$ is connected, maps into $\mathbb H$ with the principal logarithm, and contains
the upper part of the square near $\lambda=-1$.  By the identity theorem the new family
agrees there with the family of Theorem~\ref{thm:upper}.  Composed with the inverse of
$\exp\circ g$ on $U$, it gives the holomorphic extension.
\end{proof}

\begin{remark}\label{rem:endpoint}
(1) The extension across the segment is the continuation of the upper family.
It need not agree with the continuation through the lower half-plane, which is
the complex-conjugate family.  At real $b$ in the segment the partner fixed point
$L_-$ is not real, so the construction is not symmetric under conjugation.
(2) For $b\in(1,\eta)$ the partner is real: $\theta$ is purely imaginary, the image arc
falls onto the chord, and the straight crescent collapses.  There the tilted
regions of Section~\ref{sec:cusp} are needed.  Near $b=1$ ($\lambda_+\to0^-$) the root
$\theta$ goes to infinity.  The interval towards $b=0$ is covered instead
by transport through the attracting two-cycle (Theorem~\ref{thm:positive-real-axis}).
\end{remark}

\subsection{The real slice below the period-doubling point}\label{sec:two-cycle}
The interval below $\eee^{-\eee}$ admits a different transport coordinate:
the multiplier of the attracting two-cycle.  This avoids a uniform
straight-crescent estimate as $b\downarrow0$.

\begin{theorem}[Continuation across the real slice below $\eta$]\label{thm:positive-real-axis}
The family of Theorem~\ref{thm:upper} extends jointly holomorphically in
$(b,z)$ across every $b\in(0,1)$.  Together with
Theorem~\ref{thm:global}, it extends across every point of
$(0,\eta)\setminus\{1\}$.  The extensions are locally single-valued and
agree wherever their domains overlap in the upper half-plane.
\end{theorem}

\begin{proof}
Proposition~\ref{prop:attracting-real-slice} covers $(\eee^{-\eee},1)$,
Theorem~\ref{thm:endpoint} covers $\eee^{-\eee}$, and
Proposition~\ref{prop:weld-parameter} covers $(1,\eta)$ by the
identification in the proof of Theorem~\ref{thm:global}.  We construct
the extension on $(0,\eee^{-\eee})$.

Write $b=\eee^{-a}$, $a>\eee$, and let $p<q$ be the real two-cycle of
$E_b(w)=\eee^{-aw}$.  Its multiplier $m=a^2pq$ has the parametrisation
\[
 ap=u(\coth u-1),\qquad aq=u(\coth u+1),\qquad
 a=\frac{u}{\sinh u}\eee^{u\coth u},\qquad
 m=\left(\frac{u}{\sinh u}\right)^2,\quad u>0.
\]
Indeed $aq/ap=\eee^{2u}$ and $aq-ap=2u$, which give
$p=\eee^{-aq}$ and $q=\eee^{-ap}$.  Moreover
\[
 \frac{d\log a}{du}=\frac1u-\frac{u}{\sinh^2u}>0,
 \qquad \frac{d\log m}{du}=2\left(\frac1u-\coth u\right)<0.
\]
Thus $a$ increases from $\eee$ to $+\infty$, while $m$ decreases from
$1$ to $0$.  In particular $m$ is a real-analytic coordinate on the
whole interval $(0,\eee^{-\eee})$, with nonzero derivative.

Choose $b_*=\exp(-2\eee^2)$, which is in the certified segment of
Theorem~\ref{thm:endpoint}, and let $m_*$ be its two-cycle multiplier.
Denote by $B_p,B_q$ the two completely invariant basins for $E_{b_*}^2$,
distinguished by convergence to $p,q$, and by $\sigma_p,\sigma_q$ their
Koenigs coordinates, normalised to have derivative $1$ at the cycle.
They satisfy $\sigma_j\circ E_{b_*}^2=m_*\sigma_j$ and have nonzero
derivative on their respective basins.  On a simply connected domain
containing $(0,1)$ and contained in $\{0<|m|<1\}$, take the logarithm
real on $(0,1)$ and set
\[
 c(m)=\frac{\operatorname{Log}m-\log m_*}
              {\operatorname{Log}m+\log m_*},\qquad
 \nu_m=c(m)\frac{\sigma_j}{\overline{\sigma_j}}
                 \frac{\overline{\sigma_j'}}{\sigma_j'}
       \quad\text{on }B_j\setminus\sigma_j^{-1}(0),
\]
and $\nu_m=0$ elsewhere.  Since $\operatorname{Re}\operatorname{Log}m<0$,
$|c(m)|<1$.  The relation $\sigma_q\circ E_{b_*}=E_{b_*}'(p)\sigma_p$
on $B_p$, and its counterpart on $B_q$, show that $\nu_m$ is
$E_{b_*}$-invariant.  The normalised solution $h_m$ fixing $0,1,\infty$
therefore conjugates $E_{b_*}$ to an entire universal cover of
$\mathbb C^*$ taking $0$ to $1$, hence to $E_{b(m)}$ for a holomorphic
$b(m)=h_m(b_*)$.  The radial-stretch calculation of
Proposition~\ref{prop:global-transport}, applied to $E_{b_*}^2$,
shows that the transported two-cycle has multiplier $m$.
For real $m\in(0,1)$ the coefficient is conjugation-symmetric, so
$h_m$ preserves the real axis in order.  In particular
$h_m(p),h_m(q)\in(0,1)$ are a distinct real two-cycle and $b(m)>0$.
The parametrisation above
then gives $b(m)=\exp(-a(u(m)))\in(0,\eee^{-\eee})$ and
$b'(m)\ne0$.

Let $H_*$ be the marked straight crescent at $b_*$ from
Theorem~\ref{thm:endpoint}; its mark is $w_{-1}=0$.
The quotients of $h_m(H_*)$, with the edges identified by $E_{b(m)}$,
are quasiconformally equivalent to the twice-punctured quotient of $H_*$.
They have the transported fixed-point ends and the same mark $0$.
The Beltrami coefficient descends to the quotient as $c(m)$ times a
fixed coefficient.  Ahlfors--Bers uniformisation and the
$\partial_{\bar m}$-cancellation of
Proposition~\ref{prop:global-family} give a jointly holomorphic family
of normalised germs, locally in $(m,z)$: give the mark height $-1$ and
compose the inverse uniformisation with $E_{b(m)}$, as in
Definition~\ref{def:far-solution}, so that $F_m(0)=E_{b(m)}(0)=1$.
At the mark the coefficient is
real-analytic in space because $0$ belongs to a two-cycle basin and
$\sigma_j(0)\ne0$.

These germs agree with the straight-crescent family near $b_*$.
Here is the local comparison.  At $m=m_*$ the two marked crescents are
identical.  On a compact part away from the fixed-point ends, their
edges are close and the orbit-selection argument of
Lemma~\ref{lem:change-line} gives a conformal map of quotients by
locally fixed iterates of $E_{b(m)}$.  At either end, use a repelling
linearising coordinate and its logarithm.  A lift of $h_m$ commutes
with the deck translation $2\pi\ii$ and conjugates translation by
$\log\lambda_j(b_*)$ to translation by $\log\lambda_j(b(m))$.
Consequently it is a real-linear map, tending to the identity as
$m\to m_*$, plus a bounded doubly periodic remainder tending uniformly
to zero on a fundamental parallelogram.  In the logarithmic coordinate
the transported and straight edges differ by at most
$O(|m-m_*|(1+|\operatorname{Re}t|))$, whereas a dynamical iterate changes their
normal coordinate by a quantity bounded away from zero.  At most
$O(|m-m_*|(1+|\operatorname{Re}t|))$ iterates select the unique point in the
other half-open crescent; the total tangential displacement is of the
same order and remains within the linearising collar after shrinking
$|m-m_*|$.  This is the end-collar argument of
Lemma~\ref{lem:global-compare}, now with repelling rather than neutral
ends.  These orbit maps agree on overlaps, preserve the ends
and the mark, and identify the normalised inverse uniformisations.
Thus the two families agree on an open neighbourhood of $b_*$.

Finally, $b'(m)>0$ on the real interval.  Take overlapping complex
neighbourhoods of its points on which $b(m)$ is biholomorphic and whose
upper halves map into the upper base half-plane.  The transported
family and that of Theorem~\ref{thm:upper} agree on the first such
neighbourhood by the preceding comparison, hence on every successive
overlap by the identity theorem.  This gives the asserted extension
along the whole interval.
\end{proof}

The proof of Theorem~\ref{thm:global} gives more than boundary values on
$(1,\eta)$: on $N_W\cap\{\operatorname{Im}b>0\}$ it identifies the upper family
with the jointly holomorphic sewn family of
Proposition~\ref{prop:weld-parameter}.  Thus the upper continuation already
extends holomorphically across every base in $(1,\eta)$.

\begin{proposition}[Obstructions at $b=1$ and $b=0$]\label{prop:zero-one-obstruction}
Let $W\subset\mathbb C$ be a connected height domain with $0\in W$ and
$W+1\subset W$.  No normalised family $F(b,z)$ satisfying
$F(b,0)=1$ and $F(b,z+1)=\exp((\operatorname{Log}b)F(b,z))$
can be jointly holomorphic on a product of $W$ with a disc about $b=1$.
Moreover, for any such family on upper-half-plane bases approaching $b=0$,
the height value $F(b,2)$ has no holomorphic extension through $b=0$.
\end{proposition}

\begin{proof}
Suppose first that the family extended jointly holomorphically across
$b=1$.  At $b=1$ the functional equation gives $F(1,z+1)=1$ for
$z\in W$.  Since $W+1$ is a nonempty open subset of the connected domain
$W$, the identity theorem yields $F(1,z)\equiv1$ on $W$.
Put $H(z)=\partial_bF(1,z)$.  Differentiating the functional equation at
$b=1$ now gives $H(z+1)=F(1,z)=1$.  The identity theorem again gives
$H\equiv1$ on $W$, whereas $F(b,0)=1$ gives $H(0)=0$, a contradiction.

For the second assertion, the functional equation and normalisation give
$F(b,1)=b$ and $F(b,2)=\exp(b\operatorname{Log}b)$.
Along $b=\ii t$, $t\downarrow0$, this value tends to $1$, but its
derivative $\exp(b\operatorname{Log}b)(1+\operatorname{Log}b)$ is unbounded.
Hence the upper-half-plane germ has no holomorphic extension through $0$.
\end{proof}

\subsection{Consequences at the cusp}
The following statements describe the sewn family for real $b<\eta$, mostly
close to $\eta$, and hence, by Theorem~\ref{thm:cusp}(3), the boundary values
$\Kk_b$ of the continued Kneser family.
\begin{corollary}[Non-reality at fractional real heights]
\label{cor:sewn-nonreal}
For every real $b<\eta$ sufficiently close to $\eta$, and every
$\delta>0$ sufficiently small, some real $x$ with $0<|x|<\delta$
satisfies $\operatorname{Im}\Kk_b^W(x)\ne0$.  In particular the sewn
solution is not real-valued on any real interval about $z=0$.
\end{corollary}

\begin{proof}
Suppose that $\Kk_b^W(x)$ were real for all real $x$ in a neighbourhood
of $0$.  Shrink that neighbourhood so its values lie in a real interval
about $1$ on which the attracting Koenigs coordinate $\sigma_b$ is
holomorphic and real, with $\sigma_b(\Kk_b^W(x))/\sigma_b(1)>0$.
The upper representation of Theorem~\ref{thm:W2} gives
\[
 \frac{\sigma_b(\Kk_b^W(x))}{\sigma_b(1)}
   =\lambda_b^{\,x+P_b(x)},\qquad P_b(0)=0.
\]
Since $\log\lambda_b$ is real and nonzero, continuity at $x=0$ makes
$\operatorname{Im}P_b(x)=0$ on this interval.  The function $P_b$ is
holomorphic and periodic on a strip containing the whole real line,
so reflection and the identity theorem make it real on that line.
But its Fourier expansion has only non-negative modes.  A real-valued
periodic function has conjugate positive and negative coefficients;
therefore all its nonconstant coefficients vanish.  This contradicts
$c_1(b)\ne0$ for $b$ sufficiently close to $\eta$, established in
Corollary~\ref{cor:Kw}.
\end{proof}

By Theorem~\ref{thm:cusp}, this non-reality statement applies to the
continued Kneser family $\Kk_b=\Kk^W_b$ near $\eta$.
The tiny non-real part observed numerically by Paulsen at
$b=\sqrt2$~\cite[\S6]{Paulsen2019} is reproduced by $\Kk^W$ in all nine quoted
digits (Section~\ref{sec:invariant}); this is numerical, not a proof that his
continuation path gives $\Kk^W$.

\begin{corollary}[Complex cusp profile and its imaginary trace]
\label{cor:imaginary-profile}
Let $A=B_1\eee^{2\pi\ii a}\ne0$ and put $f=F_\eta^{\mathrm{att}}$.
On some fixed complex disk about $z=0$, as $b\to\eta^-$ through real bases,
\[
 \frac{\Kk_b^W(z)-\Rr_b(z)}{\Lambda(b)}
 \longrightarrow f'(z)A(\eee^{2\pi\ii z}-1)
\]
locally uniformly, hence with all $z$-derivatives.  In particular, on
a fixed real interval about $0$,
\[
 \frac{\operatorname{Im}\Kk_b^W(x)}{\Lambda(b)}
 \longrightarrow
 f'(x)
   \operatorname{Im}\!\left[A(\eee^{2\pi\ii x}-1)\right]
\]
uniformly on compact subintervals.  The right-hand side is not
identically zero on any interval about $0$.
\end{corollary}

\begin{proof}
By Theorem~\ref{thm:W2}(4), with
$B=\eee^{2\pi(Y_0+D+2)}$, there is a constant $C$ independent of $b$ and
$n$ such that $|c_n(b)|\le C(\Lambda B)^n$ for all $n\ge1$ and $b$ sufficiently
close to $\eta$.  Here we used $|\tau_n|\le M\eee^{2\pi n(Y_0+1)}$ and
absorbed the factor $\Lambda$ in the error bound into $C$.
Since $P_b(0)=0$, its constant mode is $-\sum_{n\ge1}c_n$.
Choose a disk in the common upper representation supplied by
Corollary~\ref{cor:sewn-cusp-local}.  The Fourier tail is geometric
on every smaller disk once $b$ is close enough to $\eta$, and there
\[
 P_b(z)=c_1(b)(\eee^{2\pi\ii z}-1)+O(\Lambda^2)
\]
locally uniformly.  Theorem~\ref{thm:W2}(4) and
Corollary~\ref{cor:tauconv} give $c_1/\Lambda\to A$.
On a fixed disk about $0$, $\Rr_b\to F_\eta^{\mathrm{att}}$ with its
derivatives by Corollary~\ref{cor:sewn-cusp-local} and its proof.
Taylor's formula gives
$\Kk_b^W(z)=\Rr_b(z)+\Rr_b'(z)P_b(z)+O(\Lambda^2)$
on a smaller disk.  This proves the complex limit, with derivative
convergence by Cauchy's formula.  For real $x$, $\Rr_b(x)$ and $f'(x)$
are real, giving the imaginary trace.
If its last factor vanished on an interval, both the first and second
derivatives at $x=0$ would vanish, forcing $\operatorname{Re}A=
\operatorname{Im}A=0$, contrary to $A\ne0$.
\end{proof}

\begin{corollary}[Obstruction at the real cusp]
\label{cor:cusp-real-obstruction}
For every sufficiently small complex neighbourhood $V$ of $0$, one can
choose a real $x_*\in V\setminus\{0\}$ and constants $c,\delta>0$ such that
\[
 |\operatorname{Im}\Kk_b^W(x_*)|\ge c\Lambda(b)
 \qquad(\eta-\delta<b<\eta,\ b\in\mathbb R).
\]
There is therefore no jointly holomorphic family $F(b,z)$ on a product
of a disc about $b=\eta$ and $V$ which is real at real heights for real
$b>\eta$ and agrees with $\Kk_b^W$ for all real $b<\eta$ sufficiently
close to $\eta$.  In particular, direct holomorphic continuation of a
real-normalised Kneser family through the real cusp cannot identify it
with the inner sewn branch.
\end{corollary}

\begin{proof}
The limiting profile in Corollary~\ref{cor:imaginary-profile} is not
identically zero on any real interval about $0$.  Choose $x_*$ where it
is nonzero.  Since $\Lambda(b)>0$ for real $b<\eta$, the uniform limit
gives the stated lower bound after reducing $\delta$.
For fixed real $x_*$, Schwarz reflection in the parameter says that a
holomorphic function $b\mapsto F(b,x_*)$ which is real on a real interval
above $\eta$ is also real on the real interval below $\eta$ in the same
disc.  This contradicts the lower bound if $F=\Kk^W$ there.
\end{proof}

\begin{proposition}[Logarithmic monodromy at the first backward singularity]
\label{prop:backward-log-monodromy}
Let $U$ be a simply connected parameter disc on which $a(b)=\log b\ne0$.
Suppose $F_b(z)$ is jointly holomorphic near $U\times\{-1\}$,
$F_b(-1)=0$, and $F_b'(-1)\ne0$.  Locally in $(b,t)$ write
$F_b(-1+t)=t\,u_b(t)$ with $u_b(0)\ne0$.  On any simply connected
sector $0<|t|<r$, all holomorphic backward branches satisfying
$\exp(a(b)G_b(-2+t))=F_b(-1+t)$ have the form
\[
 G_b(-2+t)=\frac{\operatorname{Log}t+\log u_b(t)+2\pi\ii k}{a(b)},
 \qquad k\in\mathbb Z.
\]
Continuation once around $t=0$ adds $2\pi\ii/a(b)$.  Hence no
single-valued holomorphic or meromorphic backward branch satisfying
the equation exists on a full punctured disc about $z=-2$.
For every real $b\in(1,\eta)$, the sewn solution $\Kk_b^W$ has a
simple zero at $z=-1$, so this description applies to it and to its
local holomorphic parameter family.
\end{proposition}

\begin{proof}
The quotient $u_b(t)$ extends jointly holomorphically through $t=0$.
After shrinking $U$ and $r$, it has a joint holomorphic logarithm.
Two logarithms of $t u_b(t)$ on a connected sector differ by a constant
multiple of $2\pi\ii$, giving the formula.  A turn about $t=0$ changes
$\operatorname{Log}t$ by $2\pi\ii$.  If a single-valued backward branch existed on
the punctured disc, integrating the logarithmic derivative of
$\exp(aG)=t u$ around a small circle would give both $0$ and $2\pi\ii$.
A meromorphic branch has no pole away from the centre, since its
exponential would have an essential singularity there; it is therefore
subject to the same contradiction.

For the sewn family, the upper chart of Proposition~\ref{prop:weld}
contains $w=-1$ and its inverse uniformising coordinate $w(z)$ is
conformal, with $w(-1)=-1$.  For real $b$,
$\Rr_b'(0)=\sigma(1)\log\lambda/\sigma'(1)\ne0$, since $\sigma(1)\ne0$ and
$\sigma'>0$ on $(-\infty,L)$ (forward iteration), and differentiating the functional
equation at $z=-1$ gives
$\Rr_b'(-1)=\Rr_b'(0)/a(b)\ne0$.  Thus
$(\Kk_b^W)'(-1)=\Rr_b'(-1)w'(-1)\ne0$.
Proposition~\ref{prop:weld-parameter} preserves this simple zero on
a sufficiently small complex parameter disc about each real base.
\end{proof}

\begin{corollary}[Conjugate inner branches and their jump]
\label{cor:conjugate-sewn-jump}
For real $b\in(1,\eta)$ put
$\Kk_b^{W,\#}(z)=\overline{\Kk_b^W(\bar z)}$ on the reflected regular
domain.  This reflected family is locally jointly holomorphic in
$(b,z)$ near every such real base.  Both branches obey the tetration
equation, are normalised at $z=0$, agree at every nonnegative integer
height where defined, and have a simple zero at $z=-1$ with the same
logarithmic monodromy $2\pi\ii/\log b$ at $z=-2$.
For $x_*$ as in Corollary~\ref{cor:cusp-real-obstruction},
\[
 |\Kk_b^W(x_*)-\Kk_b^{W,\#}(x_*)|\ge2c\Lambda(b)
 \qquad(\eta-\delta<b<\eta,\ b\in\mathbb R).
\]
More precisely, on a fixed real interval about $0$ the quotient of
their difference by $\Lambda(b)$ tends locally uniformly to
$2\ii(F_\eta^{\mathrm{att}})'(x)
\operatorname{Im}[A(\eee^{2\pi\ii x}-1)]$.
\end{corollary}

\begin{proof}
For complex $b$ near a real base, define the reflected family by
$\Kk_b^{W,\#}(z)=\overline{\Kk_{\bar b}^W(\bar z)}$.
Proposition~\ref{prop:weld-parameter} makes it jointly holomorphic.
Reflection preserves the equation for real $b$ and therefore, by the
identity theorem, on its local complex parameter disc.  The integer
values follow from the equation and normalisation.  The zero at $-1$
and its simplicity follow from Proposition~\ref{prop:backward-log-monodromy},
as does the monodromy.  On real heights the difference is exactly
$2\ii\operatorname{Im}\Kk_b^W(x)$, so the lower bound and the profile
follow from Corollaries~\ref{cor:cusp-real-obstruction} and
\ref{cor:imaginary-profile}.
\end{proof}

\begin{proposition}[The cusp splitting is beyond all algebraic orders]
\label{prop:cusp-no-puiseux}
On a sufficiently small real interval $J$ about $0$, put
\[
 M(x)=(F_\eta^{\mathrm{att}})'(x)
       \operatorname{Im}[A(\eee^{2\pi\ii x}-1)].
\]
The interval $J$ may be chosen so that $M(x)\ne0$ for every
$x\in J\setminus\{0\}$.  For every such $x$
and every integer $q\ge1$, the germ along $t>0$ of
\[
 t\longmapsto\Kk^W_{\eta-t^q}(x)
\]
has no meromorphic extension through $t=0$.  In particular, no finite
ramification of the base parameter makes this value meromorphic at
the real cusp; it has no convergent Puiseux--Laurent germ there.
If it admits a single-valued holomorphic continuation on a full
punctured $t$-disc, the singularity at $t=0$ is essential.
\end{proposition}

\begin{proof}
The profile $M$ is real analytic, vanishes at $0$, and, by
Corollary~\ref{cor:imaginary-profile}, is not identically zero on any
interval about $0$.  Thus $0$ is an isolated zero, and we shrink $J$
accordingly.
Suppose such a meromorphic germ $G(t)$ existed.  Its reflection
$G^\#(t)=\overline{G(\bar t)}$ would also be meromorphic, so
$D(t)=G(t)-G^\#(t)$ would be meromorphic near $0$.
For positive real $t$, the imaginary profile and the saddle-node
scale of $\Lambda$ give
\[
 D(t)=2\ii\operatorname{Im}\Kk^W_{\eta-t^q}(x)
 =O\!\left(\exp(-c_0t^{-q/2})\right)=o(t^N)
 \qquad\text{for every }N\ge0
\]
with some $c_0>0$.  A meromorphic Laurent germ with this decay on the
positive axis is identically zero.  This contradicts the nonzero
limit $D(t)/\Lambda(\eta-t^q)\to2\ii M(x)$.
\end{proof}

\begin{corollary}[A two-jet witness at the normalisation height]
\label{cor:cusp-two-jet}
Put $f=F_\eta^{\mathrm{att}}$ and $A=B_1\eee^{2\pi\ii a}\ne0$.
As $b\to\eta^-$ through real bases,
\[
 \frac{1}{\Lambda(b)}
 \begin{pmatrix}
   \operatorname{Im}[(\Kk_b^W)'(0)]\\
   \operatorname{Im}[(\Kk_b^W)''(0)]
 \end{pmatrix}
 \longrightarrow
 \begin{pmatrix}
   2\pi f'(0)\operatorname{Re}A\\
   4\pi f''(0)\operatorname{Re}A
     -4\pi^2 f'(0)\operatorname{Im}A
 \end{pmatrix}\ne\begin{pmatrix}0\\0\end{pmatrix}.
\]
Consequently at least one of the first two $z$-derivatives at $0$,
with the same choice of derivative for every $q\ge1$, has no
meromorphic base germ through $t=0$ after $b=\eta-t^q$.
The corresponding derivative of the reflected inner branch differs
from that of $\Kk_b^W$ by at least $c\Lambda(b)$ in modulus for some
$c>0$ and all sufficiently close real $b<\eta$.
\end{corollary}

\begin{proof}
The complex limit of Corollary~\ref{cor:imaginary-profile} and
Cauchy's formula give convergence of the first two derivatives.
The derivatives of $\Rr_b$ at $0$ are real for real $b$, giving the
displayed vector.  Since $f'(0)\ne0$ by
Corollary~\ref{cor:sewn-cusp-local}, vanishing of its first entry
forces $\operatorname{Re}A=0$, and then vanishing of its second forces
$\operatorname{Im}A=0$, a contradiction.
For a nonzero entry, the reflection argument in the proof of
Proposition~\ref{prop:cusp-no-puiseux} applies to that derivative.
Reflection doubles its imaginary part, which also gives the lower bound.
\end{proof}

\begin{remark}[Conjugate bypasses of the cusp]
Thus the monodromy of Remark~\ref{rem:cusp}(1) is quantified by
Corollary~\ref{cor:conjugate-sewn-jump}.
\end{remark}

\section{Numerical results}\label{sec:numerics}

\subsection{Ladder computations of a two-sided numerical solution}\label{sec:invariant}
In this numerical section, $\Kk_b$ denotes the extrapolated values from a
pointwise two-fixed-point construction; no base continuation is assumed.
The solution is computed at complex bases $b+\ii y$ by a
two-fixed-point $\theta$-mapping construction.  The program represents it as
$\Rr_{b+\ii y}\circ(\mathrm{id}+\theta_{\mathrm{up}})$ in the upper half-plane and
$S_{b+\ii y}\circ(\mathrm{id}+\theta_{\mathrm{dn}})$ in the lower one.  The
computation runs along ladders $y\downarrow0$ at working precisions up to
$52$ digits and is extrapolated to $y=0$.  Then $c_1$ is read off from
$D=\Kk-\Rr$.  First-mode dominance is confirmed by the constancy of
$q(z)=D(z)/(\Rr'(z)(\eee^{2\pi\ii z}-1))$ over four points, including
$z=\tfrac12-\tfrac\ii2$, where the first mode is amplified by $\eee^{\pi}$
(Table~\ref{tab:qz}).  Omitting the constant mode $c_0=-\sum_{n\ge1}c_n$ produces
a spread of $1.917$ instead.

\begin{table}[htbp]
\centering
\caption{The ratio $|q|$ at four points.  The last column is $\max|q|/\min|q|$.}
\label{tab:qz}
\small\setlength{\tabcolsep}{4pt}
\begin{tabular}{lccccc}
\toprule
base & $z=\tfrac12$ & $z=-\tfrac12$ & $z=\tfrac14$ & $z=\tfrac12-\tfrac{\ii}{2}$ & spread \\
\midrule
$1.1+0.0005\ii$ & \num{1.987636e-9} & \num{1.987636e-9} & \num{1.987636e-9} & \num{1.987637e-9} & $1.0000$\\
$1.05+0.003\ii$ & \num{1.422053e-7} & \num{1.422053e-7} & \num{1.422053e-7} & \num{1.422061e-7} & $1.0000$\\
\bottomrule
\end{tabular}
\end{table}

Table~\ref{tab:chat} lists $|\hat c_1|=|c_1|/\Lambda$.  Over the range shown,
$|D|$ varies by forty-five orders of magnitude while $|\hat c_1|$ varies by less
than $5\%$ and increases towards $|B_1|$.  The row $b=\sqrt2$ uses Paulsen's
value $\operatorname{Im}\Kk_{\sqrt2}(\tfrac12)\approx-1.18899697\cdot10^{-48}$,
from a $120$-digit contour computation~\cite{Paulsen2019}.  It shares neither
code nor method with ours.  Solving the sewing equation directly at $b=\sqrt2$
(80 digits, $32$ Fourier modes, residual $<10^{-75}$;
\nolinkurl{paulsen_sqrt2_check.py}) gives
$c_1=-2.0924123\cdot10^{-50}+1.4780488648\cdot10^{-48}\,\ii$, $\Rr'_{\sqrt2}(\tfrac12)=0.4022184246$, and
$\operatorname{Im}\Kk^W_{\sqrt2}(\tfrac12)=\operatorname{Im}[\Rr'(\tfrac12)P(\tfrac12)]+O(\Lambda^2)=-1.18899697185\cdot10^{-48}$.
This matches Paulsen's value to relative accuracy $1.6\cdot10^{-9}$, the
rounding of his quoted digits.

\begin{table}[htbp]
\centering
\caption{Ladder values of $|\hat c_1|$ and their ratio to $|B_1|$.}
\label{tab:chat}
\small\setlength{\tabcolsep}{4pt}
\begin{tabular}{llllll}
\toprule
$b$ & $\lambda$ & $\Lambda$ & $|\hat c_1|$ & $|\hat c_1|/|B_1|$ & source\\
\midrule
$1.05$ & $0.051362$ & \num{1.68e-6}  & $0.085017$ & $0.95462$ & ladder, $20$ digits\\
$1.10$ & $0.105964$ & \num{2.30e-8}  & $0.086427$ & $0.97045$ & ladder, $20$ digits\\
$1.15$ & $0.164802$ & \num{3.10e-10} & $0.087205$ & $0.97919$ & ladder, $20$ digits\\
$1.20$ & $0.229312$ & \num{2.28e-12} & $0.087732$ & $0.98511$ & ladder, $20$ digits\\
$1.25$ & $0.301735$ & \num{4.91e-15} & $0.088125$ & $0.98952$ & ladder, $24$ digits\\
$1.30$ & $0.385935$ & \num{9.82e-19} & $0.088434$ & $0.99299$ & ladder, $32$ digits\\
$\sqrt2$ & $0.693147$ & \num{1.66e-47} & $0.088941$ & $0.99868$ & \cite{Paulsen2019}\\
\bottomrule
\end{tabular}
\end{table}

The ladder values agree with the Kneser-free prediction $\tau_1$ and
$t_2/t_1^2$, moduli and arguments, to the precision of the extrapolation
(Table~\ref{tab:tau}).  Solving the sewing equation of Theorem~\ref{thm:W2}
numerically gives $\hat c_1=\tau_1(1+\vartheta)$ with $|\vartheta|=0.38\Lambda$
at $b=1.10$ and $0.40\Lambda$ at $b=1.30$.  This matches the first-order formula
$\vartheta\approx\Lambda[-4\pi^2|\tau_1|^2-2\pi\ii\tau_1-4\pi\ii\tau_2\bar\tau_1/\tau_1]$.

\subsection{Approach to the parabolic limit}\label{sec:mode2}
The transition map needs no ladder.  Its modes are measured on the line
$\operatorname{Im}z=-h_2/2+Y$ of repelling time, near the gate carrying the
positive modes, so the working precision needed is independent of $h$.
Results for $Y=1,1.5,2,2.5,3$ agree to all printed digits.
Table~\ref{tab:eta} follows the invariants to $\eta-b=2.7\cdot10^{-5}$.

\begin{table}[htbp]
\centering
\caption{Kneser-free invariants near $\eta$.  Here $\delta_1=1-|\tau_1|/|B_1|$
and $p=|\log\lambda|\log\lambda_2$.}
\label{tab:eta}
\small\setlength{\tabcolsep}{4pt}
\begin{tabular}{lllllll}
\toprule
$\lambda$ & $\eta-b$ & $|\tau_1|/|B_1|$ & $\delta_1/p$ & $\arg\tau_1$ & $|t_2/t_1^2|$ & $|t_3/t_1^3-B_3/B_1^3|$\\
\midrule
$0.5$ & $\num{9.04e-2}$ & $0.99601457$ & $0.0102076$ & $1.581415$ & $1.555683$ & $0.148$\\
$0.8$ & $\num{1.21e-2}$ & $0.99952725$ & $0.0102002$ & $1.585920$ & $1.572187$ & $0.0177$\\
$0.9$ & $\num{2.84e-3}$ & $0.99989062$ & $0.0101993$ & $1.586391$ & $1.573905$ & $0.00410$\\
$0.95$ & $\num{6.87e-4}$ & $0.99997362$ & $0.0101991$ & $1.586499$ & $1.574298$ & $0.000989$\\
$0.98$ & $\num{1.08e-4}$ & $0.99999587$ & $0.0101990$ & $1.586528$ & $1.574403$ & $0.000155$\\
$0.99$ & $\num{2.67e-5}$ & $0.99999897$ & $0.0101990$ & $1.586532$ & $1.574418$ & $0.0000385$\\
\midrule
limit & $0$ & $1$ & $0.0101990063$ & $1.5865331$ & $1.5744227$ & $0$\\
\bottomrule
\end{tabular}
\end{table}

\subsection{The deficit law}\label{sec:deficit}
The natural parameter is the product of the logarithmic multipliers,
$p=|\log\lambda|\log\lambda_2=\varepsilon^2(1-\varepsilon/3+O(\varepsilon^2))$.  Up to an analytic
factor, this is the canonical parameter of the
Marde\v{s}i\'c--Roussarie--Rousseau normal form~\cite{MRR2004}.  Fits of degree
$2$, $3$ and $4$ in $p$ over twelve bases agree to eleven digits:
\begin{equation}\label{eq:kappa1}
 \log\frac{\tau_n}{B_n\eee^{2\pi\ii na}}=\kappa^{(n)}p+O(p^2),\qquad
 \begin{aligned}
 \kappa^{(1)}&=-0.0101990063459-0.0132509561825\,\ii,\\
 \kappa^{(2)}&=-0.0510565562-0.0623157805\,\ii,\\
 \kappa^{(3)}&=-0.1400443481-0.1692244639\,\ii,
 \end{aligned}
\end{equation}
with $p^2$-coefficients of order $10^{-4}$ and constant terms below
$10^{-16}$.  Hence $\delta_1=Cp+O(p^2)$ with $C=0.0101990063$.  Theorem~\ref{thm:D}
proves the first-order statement and computes $\kappa^{(n)}$ independently of any
fit (Section~\ref{sec:first}).  By
Theorem~\ref{thm:W2}, the Kneser-side $\hat c_n$ inherit~\eqref{eq:kappa1} up to
$O(\Lambda)$, which is beyond all orders in $p$.  The $\kappa^{(n)}$ are not linear
in $n$, so the first-order effect of the unfolding is not a translation of Fatou
coordinates.

\subsection{Identification with an independent Kneser-type implementation}\label{sec:fatou}
The sewing solution can also be compared with a completely independent
two-fixed-point computation at individual bases.  Sheldon Levenstein's PARI/GP
program \texttt{fatou.gp} builds the ``merged'' two-fixed-point solution for
complex bases.  At each base it constructs, independently, the solution with a
$\theta$-mapping at each of the two fixed points: the two regular representations
of~\cite[Lemma~1]{Paulsen2019}.  Its author regards this as the continuation of
Kneser's solution in the base.  We use the unmodified program
(sha256 \texttt{70559daf\ldots}): published reference values with $100$
verified digits at $b=0.8+0.4\ii$ and $b=1+\ii$, and our own runs at
\texttt{\textbackslash p 40} at $b=1.2+0.1\ii$ and $b=1.3+0.2\ii$, near the
interval $(1,\eta)$.  In those runs $\mathrm{sexp}(0)=1$ to $10^{-53}$, and the
functional equation holds to $10^{-54}$.

We solved the sewing equation of Theorem~\ref{thm:W2} directly at these complex
bases (script \nolinkurl{weld_complex.py}).  $\Rr$ is the regular solution at the
attracting fixed point and $S$ the one at the repelling fixed point; the seam
lies in the band of holomorphy of $T=\Rr^{-1}\circ S$, and the root is selected
as in Definition~\ref{def:class}.  Table~\ref{tab:fatou} shows the result.  At
the three bases where the band is not too narrow, the sewing solution agrees
with the independent value \emph{to the full working precision of $40$ digits}.
That is $27$ to $35$ orders of magnitude below the size of the separation
$|\kappa-\Rr|$ itself.  At $1+\ii$ the band is narrowest and the agreement
($2\cdot10^{-25}$) is limited by the number of samples.  The neighbouring roots
$c_{\pm1}$ reproduce $\Rr$ instead, and miss by the full separation.  The
$\theta$-mapping builder used for the ladders does not converge at $0.8+0.4\ii$,
whereas the sewing iteration converges in a few dozen steps.

\begin{table}[htbp]
\centering
\caption{Sewing solution versus the independent \texttt{fatou.gp} values of
$\kappa_b(\tfrac12)$ (working precision $40$ digits).}
\label{tab:fatou}
\small\setlength{\tabcolsep}{5pt}
\begin{tabular}{lllll}
\toprule
$b$ & $|\Lambda|$ & $|\kappa_b(\tfrac12)-\Rr_b(\tfrac12)|$ & $|\Kk^W_b(\tfrac12)-\kappa_b(\tfrac12)|$ & $|\hat c_1|$\\
\midrule
$1.2+0.1\ii$ & \num{1.91e-11} & \num{6.91e-13} & \num{4.6e-40} & $0.0869956$\\
$1.3+0.2\ii$ & \num{2.92e-11} & \num{1.81e-12} & \num{8.3e-41} & $0.0876095$\\
$0.8+0.4\ii$ & \num{1.05e-3} & \num{9.86e-5} & \num{5.5e-39} & $0.0856072$\\
$1+\ii$ & \num{2.13e-2} & \num{4.41e-3} & \num{2.0e-25} & $0.0873850$\\
\bottomrule
\end{tabular}
\end{table}

The agreement identifies the sewing solution with an independent implementation
of the two-sided regular representation, at four bases, two of them near
$(1,\eta)$.  It does \emph{not} test the analytic continuation across the
Shell--Thron boundary: both programs work at one base at a time.  The
continuation itself is proved near the cusp in Section~\ref{sec:cusp}.

\subsection{A pointwise test near the Shell--Thron boundary}\label{sec:bndtest}
Before Theorem~\ref{thm:upper} was available, we tested the pointwise output of
the unmodified \texttt{fatou.gp} (\texttt{\textbackslash p 24}, and $50$ digits in
\nolinkurl{boundary_multiheight.py}) at six bases along the normal
$\lambda=r\eee^{\ii\phi}$ through the boundary point $b_c=\exp(\lambda_c\eee^{-\lambda_c})\approx1.5217+0.0148\ii$,
$\lambda_c=\eee^{\ii\phi}$, $\phi=2\pi(\sqrt2-1)/5$: three inside ($r=0.98,0.99,0.995$) and three
outside ($r=1.005,1.01,1.02$).  At heights $z=\tfrac k6$, $k=1,\dots,5$, a degree-$4$
polynomial in $b$ through five bases predicts the sixth to $\le1.7\cdot10^{-12}$, and
quadratics fitted on the two sides differ at $b_c$ by $\le7.1\cdot10^{-9}$, which is the size
of the truncation term (raw values in \nolinkurl{data/boundary-multiheight.json}).  The
two sides use structurally different representations, so this was an
independent check of the smooth join proved in Theorem~\ref{thm:last-arc}.

\section{Open problems}\label{sec:open}

The boundary arc up to $\eee^{-\eee}$ and the identification with Paulsen's
base-continued family, under the analytic-continuation interpretation of
Corollary~\ref{cor:paulsen}, are settled by
Theorems~\ref{thm:last-arc} and~\ref{thm:upper}.  The following questions remain.

\begin{enumerate}
\item \emph{Endpoint asymptotics and uniqueness.}
  Theorems~\ref{thm:positive-real-axis} and~\ref{thm:global} extend the
  upper family across every $b\in(0,\eta)\setminus\{1\}$.
  Proposition~\ref{prop:zero-one-obstruction} shows that $b=1$ and $b=0$
  cannot be removable singularities of a normalised jointly holomorphic
  family on a common forward-invariant height domain.
  Paulsen~\cite{Paulsen2026} studies selected noninteger heights as the base
  approaches $1$.  Determine uniform noninteger-height asymptotics for the
  present upper continuation near $b=1$ and $b=0$.  A related question is a
  global uniqueness theorem for Kneser-type solutions inside the Shell--Thron
  region.  The condition by limits at $\pm\ii\infty$ alone admits solutions of the form
  $\Rr\circ(z+2\pi\ii mz/\log\lambda+\cdots)$.
\item \emph{The first-order coefficients.}  Theorem~\ref{thm:D} expresses
  $\kappa^{(n)}$ as a convergent series over parabolic orbits.  Is there a closed
  form, for instance in terms of the horn-map data $B_n$, $a$ and the iterative
  residue?  Determine the radius of convergence and the optimal Gevrey
  order of the expansion in $p$.  The $\tfrac12$-summability results
  of~\cite{ChristopherRousseau,Ribon2010} are in a different normalisation.
\item \emph{Remaining certification.}  The coordinate $a$, the moduli
  of the first three horn coefficients, and the two horn ratios used in the
  numerical comparisons have interval enclosures above.  Give interval enclosures
  of the ladder and extrapolated two-fixed-point values.  This would turn
  the remaining numerical comparisons into computer-assisted ones.
\end{enumerate}

\section*{Code and data availability}
{\small
The exploratory computations use the Python library \nolinkurl{kneser}
(mpmath); the interval certificates use python-flint/Arb.  The
transition-map invariants are computed by \nolinkurl{transition_map.py},
\nolinkurl{eta_approach.py}, \nolinkurl{eta_sweep.py}, \nolinkurl{deficit_param.py} and
\nolinkurl{deficit_modes.py}.  The parabolic data come from
\nolinkurl{parabolic_horn.py} and \nolinkurl{parabolic_horn_inverse.py}.  The
first-order series and the second-order check are \nolinkurl{mrr_derivative.py}
and \nolinkurl{mrr_second.py}.  The boundary comparison is reproduced by
\nolinkurl{boundary_multiheight.py} from the unmodified \nolinkurl{fatou.gp}.
The confluence check is \nolinkurl{confluence_check.py}, and the sewing solution is
computed by \nolinkurl{weld_solve.py}.  The complex-base sewing solution and its comparison with \nolinkurl{fatou.gp} are
computed by \nolinkurl{weld_complex.py}.  The ladder computations of $\Kk_b$ use the
complex-base builder of the library together with \nolinkurl{c1_invariant.py} and
\nolinkurl{first_mode_check.py}.  The signs and branches used in
Section~\ref{sec:cusp} (normal component of the step, the first crossing of
the base orbit, and the seam height $-h/2$ through the lower half-plane) are
checked numerically by \nolinkurl{cusp_tilted_check.py}.  The interval
certificates are \nolinkurl{certify_far_collar.py} (boundary collar),
\nolinkurl{certify_last_boundary_arc.py} (last boundary arc),
\nolinkurl{certify_upper_halfplane.py} (upper half-plane),
\nolinkurl{certify_endpoint.py} (endpoint segment),
\nolinkurl{certify_parabolic_a.py} (the coordinate $a$) and
\nolinkurl{certify_horn_coefficients.py} (horn coefficients).\par
Code and data are available from the author on request.\par
}

\section*{Acknowledgements}
The proofs in Sections~\ref{sec:proofC}--\ref{sec:char} and~\ref{sec:cusp} were
checked by AI-assisted adversarial review, and the resulting corrections are
incorporated.
The author thanks the tetration forum for the \texttt{fatou.gp} reference values.

\appendix
\section{Interval certificates}\label{app:cert}
All certificates are Python programs using python-flint/Arb ball arithmetic with
outward rounding; every strict inequality is decided on enclosures, and floating
point is used only to choose centres of Krawczyk boxes.  Krawczyk's operator
$K=m-Yf(m)+(1-Yf'(B))(B-m)$, $K\subset\operatorname{int}B$, proves existence and uniqueness of a zero in $B$
for every parameter in the given balls.  Running times are on a laptop.
\subsection{The boundary collar}\label{app:collar}
\emph{Lemma~\ref{lem:far-cert}.}  Interval arithmetic with outward rounding (\texttt{certify\_far\_collar.py},
python-flint/Arb), on boxes in $(p,c)$ and $(p,v)$ and subintervals of $t$.
On the cone the quantities in (1) are divided by $p$, and
$|\lambda_+|^2-1$ on its top edge is divided by $p$; this makes them analytic at $p=0$.
On $\Theta_{\rm K}$ each strip of width $(P_1-P_0)/440$ is certified up to the larger
value of $V$ at its ends.  The top-edge test is made at the smaller value;
this suffices because $|\lambda_+|$ decreases in $q$.  They are evaluated through the entire series
$g_\theta=-\ii\theta\,(\theta/\sin\theta)\,h$,
$h=-\sum_{n\ge2}\ii^n\theta^{n-2}P_n(t)/n!$ with $P_n(t)=(t^n-t^{n\bmod2})/(t^2-1)$,
$|P_n|\le n/2$ on $[-1,1]$, and the Bernoulli series of $\mu$, each with an explicit tail
bound.  For $p\to0$ the limits are $\operatorname{Im}g/p\to-\tfrac12$ and
$\operatorname{Im}(\bar s'g)/p\to-\tfrac12$.

\emph{Proof of Lemma~\ref{lem:far-cone}.}  $w_{-1}=0$ lies above the chord: $\zeta(0)=\ii\cot\theta$, and Lemma~\ref{lem:far-cert}(3) applies.  Put
$q_n=(\zeta_n-1)/(\zeta_n+1)$, so that $\zeta_n=(1+q_n)/(1-q_n)$ and $\operatorname{Im}\zeta_n=2\operatorname{Im}q_n/|1-q_n|^2$.
From $s_\theta$,
\[
 q_{n+1}=\frac{\eee^{\ii\theta\zeta_n}-\eee^{\ii\theta}}{\eee^{\ii\theta\zeta_n}-\eee^{-\ii\theta}}
 =\eee^{\ii\theta}\frac{\sin(\theta(\zeta_n-1)/2)}{\sin(\theta(\zeta_n+1)/2)}
 =q_n\eee^{\ii\theta}\frac{T(v_n-\ii\theta)}{T(v_n)},\qquad v_n=\frac{\ii\theta}{1-q_n},
\]
with $T(w)=\sinh w/w$, since $\theta(\zeta_n+1)/2=-\ii v_n$ and $\theta(\zeta_n-1)/2=-\ii(v_n-\ii\theta)$.
From the Bernoulli expansion of $\coth w-1/w$ and $2\zeta(2k)\le\pi^2/3$: on $|w|\le\tfrac{9}{20}$,
$|(\log T)'(w)|\le K|w|$ and $|(\log T)''(w)|\le K_2$, with $K=1/(3(1-x))<0.341$,
$K_2=(1+x)/(3(1-x)^2)$ and $x=(9/20)^2/\pi^2$.

A series enclosure gives $\chi_0=\log q_0$ with $3p<y_0=\operatorname{Im}\chi_0<3.3p$ on the whole cone.
Put $\chi_{n+1}=\chi_n+\ii\theta+\log T(v_n-\ii\theta)-\log T(v_n)$ and $y_n=\operatorname{Im}\chi_n$.  The segment
$v_n-\tau\ii\theta=\ii\theta(1/(1-q_n)-\tau)$, $0\le\tau\le1$, has modulus at most $|\theta|(1+1/|1-q_n|)$.
Now $|1-q|\ge\sin y$ if $0<y\le\pi/2$ and $|1-q|\ge1$ if $\pi/2\le y\le3\pi/2$, whatever $\operatorname{Re}\chi$ is.  So
while $3p\le y_n\le\pi+1.2p$ the modulus is at most $1.005p(1+1/\sin3p)<\tfrac9{20}$, and the
correction term has modulus at most $K|\theta|^2(1+1/\sin 3p)<0.16p$.  Hence
$0.84p<y_{n+1}-y_n<1.16p$.  Let $n_*$ be the first index with $y_{n_*}\ge\pi$.  Then
$y_{n_*}<\pi+1.16p$, and $0<y_n<\pi$, i.e.\ $\operatorname{Im}q_n>0$, for $n<n_*$.  So $w_0,\dots,w_{n_*-1}$ lie
above the chord.

If $y_{n_*}=\pi$, then $\zeta_{n_*}\in(-1,1)$.  Otherwise let
$S=\{x_{n_*-1}+\ii y:y_{n_*-1}\le y\le\pi\}$ and $f(\chi)=\chi+\ii\theta+\log T(v-\ii\theta)-\log T(v)$, the
map $E_\theta$ in the coordinate $\chi$.  On $S$ one has $|1-q|\ge1$ and $|\zeta|\le1$, and
$f'-1=[(\log T)'(v-\ii\theta)-(\log T)'(v)]\,vq/(1-q)$ has modulus at most $2K_2|\theta|^2<\tfrac1{10}$.  So
$\operatorname{Im}f$ increases along $S$, from $y_{n_*}>\pi$ to less than $\pi+1.16p$.  Thus $f(S)$ is a
curve in the open lower $\zeta$-half-plane.  It ends at the point of $A_\theta$ that is the image of
$x_{n_*-1}+\ii\pi$, and it meets $A_\theta$ nowhere else, since $s_\theta$ is injective on the unit
disc.  $S$ approaches the chord from above, so by Lemma~\ref{lem:far-crescent}(3) $f(S)$
leaves $A_\theta$ into $H_\theta$.  Hence $f(S)$ minus its end lies in the interior of $H_\theta$, and so does
$\zeta_{n_*}$.  The constants are checked at $p=P_0$ in \texttt{certify\_far\_collar.py}.  The
bounds $|\theta|(1+1/\sin3p)$ and $K|\theta|^2(1+1/\sin3p)/p$ increase with $p$, because $p/\sin3p$ does, so
the checks at $P_0$ cover the whole cone.

\emph{Lemma~\ref{lem:far-mark} on $\Theta_{\rm K}$.}  On $\Theta_{\rm C}$ this is Lemma~\ref{lem:far-cone}.  On $\Theta_{\rm K}$ the script
encloses the orbit on parameter boxes.  It iterates the exact step of the
previous proof in mean-value form about the midpoint parameter.  For each box it
certifies $\operatorname{Im}\zeta_m>0$ for $-1\le m<n$, and then one of the following.  Either a
Krawczyk test proves $\zeta_n=N_\theta(t,\sigma)$ with $|t|<1$, $0<\sigma<1$; so $\zeta_n$ lies in the interior of
$H_\theta$, by Lemma~\ref{lem:far-crescent}(2).  Or $\zeta_{n'}$ ($n'=n$ or $n-1$) lies in a rectangle
$D=[\mathsf a,\mathsf b]\times[-\rho,\rho]$ with $-1<\mathsf a<\mathsf b<1$, $\rho\le\tfrac14$, for which Krawczyk tests prove the
following: every point of $D\cap\{\operatorname{Im}\le0\}$ is $N_\theta(t,\sigma)$ with $|t|<1$, $\sigma<1$; and
every point of $s_\theta(D\cap\{\operatorname{Im}\ge0\})$ is $N_\theta(t,\sigma)$ with $|t|<1$, $\sigma>0$, the root being
unique in a Krawczyk box that contains the known roots $(t,1)$, $t\in[\mathsf a,\mathsf b]$, of
the real edge.  In the
first case $\sigma\ge0$ is automatic, since $\operatorname{Im}N(t,\sigma)=\sigma(1-t^2)\operatorname{Im}g>0$ for $\sigma<0$.  In the
second case the root is $(t,1)$ on the real edge, by uniqueness; off the real
axis $\sigma\ne1$, since $\sigma=1$ would put a point of $s_\theta(D)$ on $A_\theta$ and contradict
injectivity.  Since $D\cap\{\operatorname{Im}>0\}$ is connected and, by
Lemma~\ref{lem:far-crescent}(3), $\sigma<1$ just above the real axis, $\sigma<1$ there.  Hence $\zeta_{n'}\in H_\theta$ if it is not above the chord, and
$\zeta_{n'+1}\in H_\theta\setminus A_\theta$ otherwise.  All $33\,027$ boxes pass (about six minutes on a laptop).

\subsection{The last boundary arc}\label{app:lastarc}
\emph{Theorem~\ref{thm:last-arc}.}  We use the chord and the coordinate $s_\theta$ of Section~\ref{sec:far}.
The interval certificate \texttt{certify\_last\_boundary\_arc.py} validates
a simple root
\[
 \theta_* = 2.29857900665128663858\ldots
              +0.76604606099318995273\ldots\ii
\]
of $\lambda_+(\theta)=-1$.  On the rectangle
$R=[2.19,2.31]+\ii[0.71,0.78]$ it verifies
$\operatorname{Re}(\lambda_+'(\theta)/\lambda_+'(\theta_*))>4/5$.
The Noshiro--Warschawski criterion makes $\lambda_+$ injective on $R$.
Validated contraction boxes give its unique inverse branch
$\theta(\alpha)\in R$ for $0.935\pi\le\alpha\le0.9999\pi$;
a contraction disk of radius $0.00019$ about $\theta_*$ covers the remaining
interval $0.9999\pi\le\alpha\le\pi$.

For $-1\le t\le1$ set
$f_1=\operatorname{Im}g_\theta(t)$ and
$f_2=\operatorname{Im}(\overline{s_\theta'(t)}g_\theta(t))$.
The certificate proves $f_1,f_2<0$ on $[-1,19/20]$ and
$\partial_t f_1,\partial_t f_2>0$ on $[19/20,1]$ along this branch.
The first part uses $79$ parameter boxes up to $0.9999\pi$ and a fixed
$0.0002$ square about $\theta_*$ for the rest.  The evaluations use the
entire series for $g_\theta$ and its $t$ derivative, with explicit tails and
outward-rounded Arb intervals.  At the endpoint of the chord there is the
exact identity
\[
 f_1(1)=f_2(1)=-\tfrac12\operatorname{Im}\lambda_+
                =-\tfrac12\sin\alpha<0 \qquad(\alpha<\pi).
\]
The positive derivatives therefore close the inequalities all the way to
$t=1$, even as their endpoint margin tends to zero when $\alpha\to\pi$.
The same certificate gives, by Krawczyk inclusion,
$\zeta_\theta(0)=\ii\cot\theta\in\operatorname{int}H_\theta$
throughout the arc.  Thus the first-entry index is always $-1$; no orbit
tracking is needed.

\subsection{The upper half-plane}\label{app:upper}
The certificate consists of four
parts (\texttt{certify\_upper\_halfplane.py}):
\begin{enumerate}
\item[(T)] \emph{Tails}, $\operatorname{Re}\lambda\le-12$.  With $\kappa=1/\lambda$ and $\theta=\pi-\delta$, $D=\delta/\kappa$ is the
  fixed point of the contraction $D=-(\pi-\kappa D)\,\ell(2\ii\kappa(\pi-\kappa D))$,
  $\ell(z)=-\log(1-z)/z$.  After multiplication by $1/|\lambda|$ all quantities are
  analytic at $\kappa=0$.  $F$ is a contraction on the fixed convex set
  $|\operatorname{Re}D+\pi|,|\operatorname{Im}D|\le\tfrac32$, so $D$ is unique there.  The quantities are checked on boxes in $(1/|\operatorname{Re}\lambda|,\operatorname{Im}\lambda)$ that
  include the limit.  The mark is $w_{-1}=0$: $\zeta(0)=\ii\cot\theta=N_\theta(t,\sigma)$ with $(t,\sigma)$
  near $(0,\tfrac12)$, by a scaled Krawczyk test.  There are $6230$ boxes.
\item[(K)] \emph{Core}, $-12\le\operatorname{Re}\lambda\le5$, $|\lambda-1|\ge0.08$.  Boxes are discarded when they
  certainly violate $|\lambda|\ge1$ or $0\le\operatorname{Im}g\le\pi$, or satisfy $g(\lambda)\in g(\mathbb D)$.  The last
  test uses that $g(\mathbb D)$ is starlike with boundary $\rho=\eee^{-\cos\phi}$ at argument
  $\psi=\phi-\sin\phi$.  For $\operatorname{Re}\lambda\ge5$, $|g|\le(5+\pi)\eee^{-5}<\eee^{-1}$, so $g(\lambda)\in g(\mathbb D)$.
  On each remaining box a Krawczyk test gives the unique root $\theta$ in a
  small box.  Adjacent boxes are checked to carry the same root, by uniqueness in
  the hull of their $\theta$-boxes.  The chain starts at a real $\theta$.  Then (1) of
  Lemma~\ref{lem:far-cert} is checked, together with the endpoint identities
  $f_1(1)=f_2(1)=-\tfrac12\operatorname{Im}\lambda$ and $f_1(-1)=f_2(-1)=\tfrac12\operatorname{Im}(\lambda-2\ii\theta)$ and a sign
  of $\partial_tf$ near $t=\pm1$.  This handles the degenerate edge $\operatorname{Im}\lambda\to0$
  ($b\to(0,\eee^{-\eee})$).  The mark is checked as in Lemma~\ref{lem:far-mark}.  A box
  whose $\theta$-box lies in $\Theta$ or in its complex conjugate needs no new check.
  Indeed $s_{\bar\theta}(t)=-\overline{s_\theta(-t)}$ and $\zeta_{\bar\theta}(\bar w)=-\overline{\zeta_\theta(w)}$, so the
  crescent, the inequalities (with $t\mapsto-t$) and the entry index of $\bar\theta$ are those of
  $\theta$ reflected in the imaginary axis.
\item[(P)] \emph{Cusp}, $|\lambda-1|<0.08$.  Write $\lambda_+(\theta)=1+\ii\theta+\theta^2D(\theta^2)$ with $|D|\le0.34$ for
  $|\theta|\le0.09$ (certified).  By Rouch\'e, $\lambda_+(\theta)=\lambda$ has exactly one root with $|\theta|<0.09$.  In polar form $\theta=\rho\eee^{\ii\phi}$, $\rho\le0.09$, the conditions
  $|\lambda|\ge1$, $\operatorname{Im}\lambda>0$, $\operatorname{Im}g>0$, scaled by $\rho$, $\rho$, $\rho^2$, force $0\le\phi\le\arctan\tfrac1{10}$.  The
  scaled $\operatorname{Im}g$ vanishes identically along the real directions $\phi=0$ and $\phi=-\pi/2$,
  and so does $\operatorname{Im}\lambda$ along $\phi=-\pi/2$; there the sign of $\partial_\phi$ is used.  So $\theta$ lies in the cone
  $\Theta_{\rm C}$, where Lemmas~\ref{lem:far-cert} and~\ref{lem:far-cone} apply.
\item[(I)] \emph{Interfaces and anchor.}  The core roots at $\operatorname{Re}\lambda=-12$ coincide with the tail
  roots.  In the tails the root is the unique fixed point of the contraction on the
  fixed set $|\operatorname{Re}D+\pi|,|\operatorname{Im}D|\le\tfrac32$, so it is one continuous branch there.  The core
  $\theta$-boxes on cells meeting $|\lambda-1|<0.08$ lie in $0<|\theta|<0.09$, so they carry the Rouch\'e
  root.  At $\lambda=\cot1+\ii$ the certified box contains the real root $\theta=1$ and not $0$.
\end{enumerate}

\subsection{The endpoint segment}\label{app:endpoint}
The certificate \texttt{certify\_endpoint.py} (python-flint/Arb, under a minute)
checks (a)--(f) of Lemma~\ref{lem:chi-crescent} for all $\lambda$ in the square
$[-1.1,-0.9]\times[-0.1,0.1]$ and in the strip $[-12,-\tfrac15]\times[-10^{-3},10^{-3}]$.
\begin{itemize}
\item \emph{The square.}  $\operatorname{Re}\bigl(\lambda_+'(\theta)/\lambda_+'(\theta_*)\bigr)>0$ on the convex rectangle
  $R'=[2.20,2.45]\times[0.62,0.95]$ ($\theta_*$ is the root of $\lambda_+=-1$), so $\lambda_+$ is injective there
  (Noshiro--Warschawski).  Krawczyk boxes inside $R'$ give the root $\theta(\lambda)$.  Conditions
  (a)--(f) are checked on $t$-intervals, and (e) by comparing $\zeta(0)$ with the arc.
\item \emph{The strip.}  The first strip box lies in the square, and its root enclosure lies in
  $R'$; this anchors the strip to the square's branch.  Adjacent boxes carry the same root:
  a Krawczyk test on the hull of their $\theta$-boxes, for $\lambda$ on the common edge, proves that
  the root is unique there.  On $t$-intervals touching $\pm1$, where $x=\mp\infty$, conditions (d) and
  (e) use the one-sided bounds $x\le x_{\max}$ and $x\ge x_{\min}$.
\end{itemize}

\subsection{The parabolic constants}\label{app:horn}
\emph{Proposition~\ref{prop:horn-intervals}.}  The certificate \texttt{certify\_horn\_coefficients.py} uses the first
$12$ formal Abel coefficients as exact rationals.  For their truncation
$\alpha_{12}$, the Abel defect vanishes to order $14$ and has modulus
$<4$ on $|u|=\tfrac12$.  Once the forward and backward orbits reach
$\operatorname{Re}(-2/u)>80$ and $\operatorname{Re}(2/u)>80$, respectively,
the defect estimate and the increase of these real parts by at least
$3/4$ per step bound each Fatou-coordinate tail after $100$ iterates
by $2.25\cdot10^{-17}$, and after $1200$ iterates by
$4.08\cdot10^{-31}$.

Outward-rounded Arb arithmetic covers
$0\le\operatorname{Re}z\le1$, $1.9\le\operatorname{Im}z\le2.1$ with
$256\times40$ boxes.  On each box a Krawczyk test gives a unique local
inverse of $\Phi_{\mathrm{rep}}$; adjacent roots agree on their common
edges, and the roots at $\operatorname{Re}z=0,1$ are related by
$f(u)=\eee^u-1$.  The same boxes give
$|D(z)|<1$ for
$D=\Phi_{\mathrm{att}}\circ\Phi_{\mathrm{rep}}^{-1}-\mathrm{id}$.
A chain of $400$ further boxes joins this branch to $u=11\ii/20$ in
the upper gate, fixing the horn-map normalisation.

At $z_j=j/128+2\ii$ the inverse and $D(z_j)$ are enclosed using
$1200$ iterates.  For $n=1,2,3$, multiply the discrete Fourier average
by $\eee^{4\pi n}$.  The analytic periodic trapezoid bound on the
strip of half-width $\delta=1/10$ makes its alias error at most
\[
 2\eee^{2\pi n(2+\delta)}
 \frac{\eee^{-2\pi(128)\delta}}
      {1-\eee^{-2\pi(128)\delta}}.
\]
Adding this bound to the Arb sample enclosures gives the stated
intervals.

\emph{The coordinate $a$.}  \texttt{certify\_parabolic\_a.py} computes the first ten formal Abel
coefficients as exact rationals.  For their truncation $\alpha_{10}$, the
residual $\alpha_{10}(\eee^u-1)-\alpha_{10}(u)-1$ vanishes to order $12$;
Arb bounds its modulus by $0.438$ on $|u|=\tfrac12$.  Put
$u_0=1/\eee-1$, $u_{k+1}=\eee^{u_k}-1$, and $x=-u_{1000}<0.002$.
Since $-u_{k+1}\le(-u_k)-(-u_k)^2/3$ for $-u_k\le1$, Cauchy's estimate
bounds the tail after $\alpha_{10}(u_{1000})-1000$ by
\[
 0.438\,2^{12}\left(x^{12}+\frac{3x^{11}}{11}\right)<10^{-27}.
\]
The orbit and truncated value are enclosed by Arb; together they give the
stated interval.

\begin{thebibliography}{99}

\bibitem{AhlforsLectures}
L.~V.~Ahlfors,
\emph{Lectures on Quasiconformal Mappings}, 2nd ed.,
University Lecture Series~38, Amer.\ Math.\ Soc., Providence, RI, 2006.

\bibitem{AhlforsBers1960}
L.~Ahlfors and L.~Bers,
\emph{Riemann's mapping theorem for variable metrics},
Ann.\ of Math.\ (2) \textbf{72} (1960), 385--404.

\bibitem{BuffEcalleEpstein2013}
X.~Buff, J.~\'Ecalle and A.~Epstein,
\emph{Limits of degenerate parabolic quadratic rational maps},
Geom.\ Funct.\ Anal.\ \textbf{23} (2013), 42--95.

\bibitem{BuffEpstein2002}
X.~Buff and A.~L.~Epstein,
\emph{A parabolic Pommerenke--Levin--Yoccoz inequality},
Fund.\ Math.\ \textbf{172} (2002), 249--289.
\href{https://www.impan.pl/shop/en/publication/transaction/download/product/89261}{Publisher PDF}.

\bibitem{Cheritat2022}
A.~Ch\'eritat,
\emph{Near parabolic renormalization for unicritical holomorphic maps},
Arnold Math.\ J.\ \textbf{8} (2022), 169--270.
\href{https://doi.org/10.1007/s40598-020-00172-6}{doi:10.1007/s40598-020-00172-6}.

\bibitem{ChristopherRousseau}
C.~Christopher and C.~Rousseau,
\emph{The moduli space of germs of generic families of analytic
diffeomorphisms unfolding a parabolic fixed point},
Int.\ Math.\ Res.\ Not.\ IMRN (2014); arXiv:0809.2167.

\bibitem{DudkoSauzin2014}
A.~Dudko and D.~Sauzin,
\emph{On the resurgent approach to \'Ecalle--Voronin's invariants},
C.\ R.\ Math.\ Acad.\ Sci.\ Paris \textbf{353} (2015), 265--271;
\href{https://arxiv.org/abs/1307.8095}{arXiv:1307.8095}.

\bibitem{Ecalle1981}
J.~\'Ecalle,
\emph{Les fonctions r\'esurgentes}, Publ.\ Math.\ d'Orsay, 1981.

\bibitem{Epstein1993}
A.~L.~Epstein,
\emph{Towers of finite type complex analytic maps},
Ph.D. thesis, City University of New York, 1993.

\bibitem{FischerGrauert1965}
W.~Fischer and H.~Grauert,
\emph{Lokal-triviale Familien kompakter komplexer Mannigfaltigkeiten},
Nachr.\ Akad.\ Wiss.\ G\"ottingen Math.-Phys.\ Kl.\ II (1965), 89--94.

\bibitem{GelfreichBrannstrom}
V.~Gelfreich and N.~Br\"annstr\"om,
\emph{Asymptotic series for the splitting of separatrices near a Hamiltonian
bifurcation}, preprint, 2008,
\href{https://arxiv.org/abs/0806.2403}{arXiv:0806.2403}.

\bibitem{GelfreichSauzin2001}
V.~Gelfreich and D.~Sauzin,
\emph{Borel summation and splitting of separatrices for the H\'enon map},
Ann.\ Inst.\ Fourier (Grenoble) \textbf{51} (2001), 513--567.

\bibitem{GelfreichSimo2008}
V.~Gelfreich and C.~Sim\'o,
\emph{High-precision computations of divergent asymptotic series and
homoclinic phenomena},
Discrete Contin.\ Dyn.\ Syst.\ Ser.\ B \textbf{10} (2008), 511--536.

\bibitem{Glutsyuk2001}
A.~A.~Glutsyuk,
\emph{Confluence of singular points and the nonlinear Stokes phenomenon},
Trans.\ Moscow Math.\ Soc.\ \textbf{62} (2001), 49--95.

\bibitem{InouShishikura}
H.~Inou and M.~Shishikura,
\emph{The renormalization for parabolic fixed points and their perturbation},
preprint, 2008,
\url{https://www.math.kyoto-u.ac.jp/~mitsu/pararenorm/}.

\bibitem{Kneser1950}
H.~Kneser,
\emph{Reelle analytische L\"osungen der Gleichung
$\vartheta(\vartheta(x))=\eee^{x}$ und verwandter Funktionalgleichungen},
J.\ Reine Angew.\ Math.\ \textbf{187} (1950), 56--67.

\bibitem{Kouznetsov2009}
D.~Kouznetsov,
\emph{Solution of $F(z+1)=\exp(F(z))$ in complex $z$-plane},
Math.\ Comp.\ \textbf{78} (2009), 1647--1670.

\bibitem{KouznetsovTrappmann2010}
D.~Kouznetsov and H.~Trappmann,
\emph{Portrait of the four regular super-exponentials to base $\sqrt2$},
Math.\ Comp.\ \textbf{79} (2010), 1727--1756.

\bibitem{Lavaurs1989}
P.~Lavaurs,
\emph{Syst\`emes dynamiques holomorphes: explosion de points p\'eriodiques
paraboliques}, Th\`ese, Universit\'e Paris-Sud, Orsay, 1989.

\bibitem{LehtoVirtanen}
O.~Lehto and K.~I.~Virtanen,
\emph{Quasiconformal Mappings in the Plane}, 2nd ed.,
Grundlehren der mathematischen Wissenschaften~126, Springer, 1973.

\bibitem{MRR2004}
P.~Marde\v{s}i\'c, R.~Roussarie and C.~Rousseau,
\emph{Modulus of analytic classification for unfoldings of generic parabolic
diffeomorphisms},
Mosc.\ Math.\ J.\ \textbf{4} (2004), no.~2, 455--502.

\bibitem{Nixon2022}
J.~D.~Nixon,
\emph{Asymptotic solutions of the tetration equation},
preprint, 2022, \href{https://arxiv.org/abs/2208.05328}{arXiv:2208.05328}.

\bibitem{Oudkerk1999}
R.~Oudkerk,
\emph{The parabolic implosion for
$f_0(z)=z+z^{\nu+1}+O(z^{\nu+2})$},
Ph.D.\ thesis, University of Warwick, 1999;
see also \emph{The parabolic implosion: Lavaurs maps and strong convergence for
rational maps}, in: Value Distribution Theory and Complex Dynamics,
Contemp.\ Math.\ \textbf{303}, Amer.\ Math.\ Soc., 2002, 79--105.

\bibitem{PaulsenCowgill2017}
W.~Paulsen and S.~Cowgill,
\emph{Solving $F(z+1)=b^{F(z)}$ in the complex plane},
Adv.\ Comput.\ Math.\ \textbf{43} (2017), 1261--1282.
\href{https://doi.org/10.1007/s10444-017-9524-1}{doi:10.1007/s10444-017-9524-1}.

\bibitem{Paulsen2019}
W.~Paulsen,
\emph{Tetration for complex bases},
Adv.\ Comput.\ Math.\ \textbf{45} (2019), 243--267.
\href{https://doi.org/10.1007/s10444-018-9615-7}{doi:10.1007/s10444-018-9615-7}.

\bibitem{Paulsen2026}
W.~Paulsen,
\emph{Analyzing the function $z\uparrow\uparrow a$ for fractional $a$},
Adv.\ Comput.\ Math.\ \textbf{52} (2026), article~25.
\href{https://doi.org/10.1007/s10444-026-10300-z}{doi:10.1007/s10444-026-10300-z}.

\bibitem{Ribon2010}
J.~Rib\'on,
\emph{Multisummability of unfoldings of tangent to the identity
diffeomorphisms}, preprint, 2010,
\href{https://arxiv.org/abs/1009.3518}{arXiv:1009.3518}.

\bibitem{RousseauTeyssier2008}
C.~Rousseau and L.~Teyssier,
\emph{Analytical moduli for unfoldings of saddle-node vector fields},
Mosc.\ Math.\ J.\ \textbf{8} (2008), no.~3, 547--614.

\bibitem{Shishikura2000}
M.~Shishikura,
\emph{Bifurcation of parabolic fixed points},
in: \emph{The Mandelbrot Set, Theme and Variations},
London Math.\ Soc.\ Lecture Note Ser.\ \textbf{274},
Cambridge Univ.\ Press, 2000, 325--363.

\bibitem{KouznetsovTrappmann2012}
H.~Trappmann and D.~Kouznetsov,
\emph{Computation of the two regular super-exponentials to base
$\exp(1/\eee)$},
Math.\ Comp.\ \textbf{81} (2012), 2207--2227.

\bibitem{TrappmannKouznetsov2011}
H.~Trappmann and D.~Kouznetsov,
\emph{Uniqueness of holomorphic Abel functions at a complex fixed point pair},
Aequationes Math.\ \textbf{81} (2011), 65--76.

\bibitem{Voronin1981}
S.~M.~Voronin,
\emph{Analytic classification of germs of conformal mappings
$(\mathbb C,0)\to(\mathbb C,0)$},
Funct.\ Anal.\ Appl.\ \textbf{15} (1981), 1--13.

\end{thebibliography}

\end{document}
